% concatenated from Z_SBP-correction.tex
\documentclass[12pt, oneside, psamsfonts]{amsart}

%-------------pdflatex or latex--------------------------------------
\newif\ifPDF
\ifx\pdfoutput\undefined\PDFfalse
\else \ifnum \pdfoutput > 0 \PDFtrue
        \else \PDFfalse
        \fi
\fi

%-------------packages---------------------------------------------
%\usepackage{showkeys}
\usepackage[centertags]{amsmath}
\usepackage{amsfonts}
\usepackage{mathrsfs}
\usepackage{textcomp}
\usepackage{amssymb}
\usepackage{amsthm}
\usepackage{newlfont}
\usepackage[all]{xy}


%------------packages for pdflatex------------------------------

\ifPDF
  \usepackage[pdftex]{color, graphicx}
  \usepackage[pdftex, bookmarks, colorlinks]{hyperref}
  \hypersetup{colorlinks=false}

%------------\packages for latex---------------------------------

\else
  \usepackage{color}
  \usepackage[dvips]{graphicx}
  \usepackage[dvips]{hyperref}
\fi


%-------------------------------------PAGE LAYOUT---------------------------------------------------

\usepackage[scale=0.8]{geometry}

%------------------------------------ the space between lines -----------------------------------

\renewcommand{\baselinestretch}{1.1}

% SHOW LINE NUMBERS---------------------------------------

\usepackage[pagewise, mathlines, displaymath]{lineno}



%------------------------------------ THEOREMS -------------------------------------------------------
\newtheorem{thm}{Theorem}[section]
\newtheorem{cor}[thm]{Corollary}
\newtheorem{lem}[thm]{Lemma}
\newtheorem{prop}[thm]{Proposition}
\theoremstyle{definition}
\newtheorem{defn}[thm]{Definition}
\theoremstyle{remark}
\newtheorem{rem}[thm]{Remark}
\newtheorem{example}[thm]{Example}
\numberwithin{equation}{section}

\newtheorem{NN}[thm]{}

% -------------------------------------MATH -----------------------------------------------------------
\newcommand{\norm}[1]{\| #1\|}
\newcommand{\abs}[1]{| #1 |}
\newcommand{\set}[1]{\left\{#1\right\}}
\newcommand{\Real}{\mathbb R}
\newcommand{\Int}{\mathbb Z}
\newcommand{\Comp}{\mathbb C}
\newcommand{\Ratn}{\mathbb Q}
\newcommand{\eps}{\varepsilon}
\newcommand{\To}{\longrightarrow}
\newcommand{\F}{\mathcal{F}}
\newcommand{\Kzero}{\textrm{K}_0}
\newcommand{\Kone}{\textrm{K}_1}
\newcommand{\tr}{\mathrm{T}}
\newcommand{\MC}[2]{\mathrm{M}_{#1}(\mathrm{C}(#2))}
\newcommand{\MCO}[2]{\mathrm{M}_{#1}(\mathrm{C}_0(#2))}
\newcommand{\Mat}[1]{\mathrm{M}_{#1}(\mathbb C)}
\newcommand{\aff}{\mathrm{Aff}}
\newcommand{\St}{\mathrm{S}}
\newcommand{\KC}[1]{{\mathrm{C}(#1)}\otimes{\mathcal K}}

%------------------------------------------------------title------------------------------------------------

\begin{document}

%OPTIONS FOR LINE NUMBERS----------------------------------
%\linenumbers


\title{On the small boundary property and $\mathcal Z$-absorption}

\author{George A. Elliott}
\address{Department of Mathematics, University of Toronto, Toronto, Ontario, Canada~\ M5S 2E4}
\email{elliott@math.toronto.edu}


\author{Zhuang Niu}
\address{Department of Mathematics and Statistics, University of Wyoming, Laramie, Wyoming 82071, USA}
\email{zniu@uwyo.edu}


%\thanks{}
\keywords{$\mathcal Z$-Absorbing C*-Algebras, Small Boundary Property, Mean Topological Dimension}
\date{\today}
%\dedicatory{}
%\commby{}
\subjclass{46L35, 37B02}

%------------------------------------------abstract--------------------------------------

\begin{abstract}
We introduce Property (C) for a unital commutative sub-C*-algebra $D$ of a unital C*-algebra $A$, which is a version of the relative comparison property using almost normalizers. In the case that $D = \mathrm{C}(X)$ and $A = \mathrm{C}(X) \rtimes\Gamma$, where $(X, \Gamma)$ is a free and minimal dynamical system with uniform Rokhlin property, it turns out that this Property (C) is equivalent to the small boundary property of $(X, \Gamma)$.


%Under the assumption of this property, the $\mathcal Z$-absorption of $A$ is shown to imply the small boundary property of $(D, \mathrm{T}(A)|_D)$, where $A =\mathrm{C}(X) \rtimes \Int^d $ and $D = \mathrm{C}(X)$. %Property $(E)$ holds for all simple AH-algebras with diagonal maps and all $A=\mathrm{C}(X) \rtimes \Int^d$, where $(X, \Int^d)$ is a minimal and free topological dynamical systems. %But we do not know if Property (C) always holds.

\end{abstract}

\maketitle

\setcounter{tocdepth}{1}

%\tableofcontents

\section{Introduction}

This is a continuation of our study \cite{EN-SBP-I} of the relation between the small boundary property of dynamical systems and the $\mathcal Z$-absorption of C*-algebras.

The Jiang-Su algebra $\mathcal Z$ is an infinite-dimensional unital simple separable amenable C*-algebra which has the same value of the Elliott invariant as $\Comp$. A C*-algebra $A$ is said to be $\mathcal Z$-absorbing if $A \cong A \otimes \mathcal Z$, and the class of $\mathcal Z$-absorbing C*-algebras (which includes $\mathcal Z$ itself) is considered to be well behaved. In fact, the class of $\mathcal Z$-absorbing unital simple separable amenable C*-algebras which satisfy the Universal Coefficient Theorem of KK-theory (possibly redundant) is  classified by the conventional Elliott invariant (see \cite{GLN-TAS-1}, \cite{GLN-TAS-2}, \cite{EN-K0-Z}, \cite{EGLN-DR}, \cite{EGLN-ASH}, \cite{TWW-QD}, \cite{CETWW-dim-n}).

On the other hand, the small boundary property of a topological dynamical system was introduced in \cite{Lindenstrauss-Weiss-MD} as a dynamical system analogue of the usual definition of zero-dimensional space. It is closely related to the mean (topological) dimension, which was introduced by Gromov (\cite{Gromov-MD}), and then was developed and studied systematically by Lindenstrauss and Weiss (\cite{Lindenstrauss-Weiss-MD}), as a numerical invariant which measures the complexity of a dynamical system in terms of the dimension growth with respect to partial orbits. The small boundary property always implies the zero mean dimension (\cite{Lindenstrauss-Weiss-MD}), and the converse holds for $\mathbb Z^d$-actions with the marker property (\cite{Lind-MD}, \cite{GLT-Zk}).

It was shown in \cite{EN-MD0} that the small boundary property of $(X, \Int)$ implies $\mathcal Z$-absorption of the crossed product C*-algebra $ A = \mathrm{C}(X) \rtimes \Int$. In this paper, we investigate the converse: how to recover the small boundary property of $(X, \Int)$ in terms of C*-algebra information.

We consider the following property for a pair of C*-algebras $D \subseteq A$:  
\newtheorem*{defn-n}{Definition}
\begin{defn-n}[Definition \ref{defn-C}]
Let $A$ be a unital C*-algebra and let $D$ be a unital commutative sub-C*-algebra. Then the pair $(D, A)$ is said to have Property (C) if for any positive contractions $f, g, h \in D$ satisfying $f, g \in \overline{hDh}$, and $$\mathrm{d}_\tau(f) < \mathrm{d}_\tau(g),\quad \tau \in \tr(A),$$ and for any $\eps>0$, there is a contraction $u \in \overline{hAh} + \Comp 1_A$ such that
$$ u f u^* \in_\eps^{\norm{\cdot}_2} \overline{gAg},$$
$$\mathrm{dist}_{2, \tr(A)}(udu^*, (D)_1) < \eps,\quad \mathrm{dist}_{2, \tr(A)}(u^*du, (D)_1) < \eps, \quad d \in (D)_1, $$ 
and
$$\norm{uu^* - 1}_{2, \tr(A)},\  \norm{u^*u - 1}_{2, \tr(A)} < \eps.$$
\end{defn-n}

Property (C) can be regarded as a relative comparison property for $D$ inside $A$, with respect to the uniform trace norm, but requires the comparison being implemented by almost normalizers. This property for C*-pair turns out to be equivalent to the small boundary property of the dynamical system $(X, \Gamma)$: 

 %While property (E) is an existence property for affine functions of the trace simplex.

%Under the assumption of Property (C), $\mathcal Z$-absorption of $A$ implies that $(X, \Int^d)$ has the small boundary property:
%\theoremstyle{theorem}
%\newtheorem*{thmA}{Theorem A}
\begin{thm}[Corollary \ref{Z-SBP}]
Let $(X, \Int^d)$ be a free and minimal topological dynamical system, and let $A = \mathrm{C}(X)\rtimes \Int^d$. Then $(\mathrm{C}(X), A)$ has Property (C) if, and only if $(X, \Int^d)$ has the (SBP). % If $(\mathrm{C}(X), A)$ has Property (C), then $$A \cong A \otimes \mathcal Z \quad \Longrightarrow \quad \textrm{$(X, \Int^d)$ has the (SBP)}. $$
\end{thm}


This theorem also holds for a free and minimal action of an arbitrary  amenable group with the uniform Rokhlin Property (Definition \ref{defn-URP-COS}). % (properties which may always hold).


In order to show the (SBP) for the pair $(D, \tr(A))$, by \cite{EN-SBP-I}, it is enough to show that every self-adjoint element of $D$ can be approximated by self-adjoint elements of $D$ with a neighbourhood of $0$ uniformly small under all tracial spectral measures (see Theorem \ref{SBP-2-norm}). Such approximating elements exist (Property (S)), but only in the ambient C*-algebra $A$. However, this can be fixed by using Properties (C): Upon using  Property (E) (Definition \ref{defn-D}), an existence property which always holds for the C*-algebras in question, one obtains a self-adjoint element in the subalgebra $D$ which almost has the same trace spectral measure distributions as the self-adjoint element in $A$ provided by Property (S), and then upon using Property (C), this element can be twisted inside $D$ to approximate the given self-adjoint element and still have a neighbourhood of $0$ uniformly small under all tracial spectral measures. This shows the (SBP) for $(D, \tr(A))$.

Nevertheless, Property (C) still involves the pair of C*-algebras $D \subseteq A$. The question is still open whether the small boundary property of $(X, \Gamma)$ can be recovered just from the ambient C*-algebra $A$, or in particular, whether $A \cong A\otimes \mathcal Z$ implies that the small boundary property of $(X, \Gamma)$. This question seems to be close related to  possible uniqueness of the Cartan subalfgebras of $A$ with respect to the uniform trace norm. 


\subsection*{Acknowledgements} The research of the first named author was supported by a Natural Sciences and Engineering Research Council of Canada (NSERC) Discovery Grant, and the research of the second named author was supported by a Simons Foundation grant  (MP-TSM-00002606). Part of the work was carried out during the visit of the second named author to Kyushu University in December 2024; he  thanks Yasuhiko Sato for the discussions and the hospitality during the visit. AI is not used during the research of this work.

\section{Preliminaries and notation}

In this section, let us collect some notation and definitions concerning C*-algebras and dynamical systems.


\subsection{Comparison of positive elements and $\mathcal Z$-absorbing C*-algebras}

Let $A$ be a C*-algebra, and let $a, b \in A$ be positive elements. One says that $a$ is Cuntz subequivalent to $b$, denoted by $a \precsim b$, if there is a sequence $(x_n)$ in $A$ such that
$$ \lim_{n \to \infty} x_n^*b x_n = a.$$
The following lemma will be frequently used.
\begin{lem}[\cite{RorUHF2}]
If $a, b\in A$ are positive elements such that $\norm{a - b} < \eps$, then $(a-\eps)_+ \precsim b$, where $(a-\eps)_+ = f(a)$ with $f(t) = \max\{t-\eps, 0\}$, $t \in \Real$.
\end{lem}


Let $\tau \in \tr(A)$. For each positive element $a \in A$, define
$$\mathrm{d}_\tau(a) = \lim_{n\to\infty} \tau(a^{\frac{1}{n}}).$$
Then, if $a \precsim b$, one has $$\mathrm{d}_\tau(a) \leq \mathrm{d}_\tau(b),\quad \tau \in \tr(A).$$ The converse in general does not hold. 

\begin{defn}[\cite{JS-Z}]
The Jiang-Su algebra $\mathcal Z$ is the (unique) simple unital inductive limit of dimension drop C*-algebras such that $$(\Kzero(\mathcal Z), \Kzero^+(\mathcal Z), [1_\mathcal Z]_0) \cong (\Int, \Int^+, 1),\quad  \Kone(\mathcal Z) = \{0\},\quad \mathrm{and} \quad \tr(\mathcal Z) = \{\mathrm{pt}\}.$$ 

A C*-algebra $A$ is said to be $\mathcal Z$-absorbing if $A \cong A \otimes \mathcal Z$.
\end{defn}

If $A$ is simple and $\mathcal Z$-absorbing, then the strict order induced by the Cuntz subequivalence relation is determined by the rank functions; that is, for any positive elements $a, b \in A$,
$$ \mathrm{d}_\tau(a) < \mathrm{d}_\tau(b),\quad \tau \in \tr(A) \quad \Longrightarrow \quad  a \precsim b.  $$
The Toms-Winter conjecture asserts that strict comparison implies $\mathcal Z$-absorption for simple separable amenable C*-algebras (this was verified for C*-algebras with finitely many extreme traces in \cite{Matui-Sato-CP}, and then it was generalized to C*-algebras with extreme traces being compact and finite dimensional; see \cite{Sato-CP}, \cite{KR-CenSeq}, \cite{TWW-Z}).

The class of simple separable amenable $\mathcal Z$-absorbing C*-algebras which satisfy the Universal Coefficient Theorem can be classified by the conventional Elliott invariant, which in the unital case consists of the K-groups and the pairing with the trace simplex (the order on the K-group, not redundant in more general cases, is determined by the pairing) (see \cite{GLN-TAS-1}, \cite{GLN-TAS-2}, \cite{EN-K0-Z}, \cite{EGLN-DR}, \cite{EGLN-ASH}, \cite{TWW-QD}, \cite{CETWW-dim-n}):
\begin{thm}\label{classification-Z}
Let $A, B$ be unital simple separable amenable $\mathcal Z$-absorbing C*-algebras which satisfy the UCT. Then $$ A \cong B \quad \Longleftrightarrow \quad \mathrm{Ell}(A) \cong \mathrm{Ell}(B), $$ where $\mathrm{Ell}(\cdot)$ denotes the Elliott invariant. Moreover, any isomorphism between the Elliott invariant can be lifted to an isomorphism between the C*-algebras.
\end{thm}

\subsection{Uniform trace norm}

\begin{defn}
Let $A$ be a unital C*-algebra, and let $\tau \in \tr(A)$. Define
$$  \norm{a}_{2, \tau} = (\tau(a^*a))^{\frac{1}{2}}, \quad a \in A.$$
For any set $\Delta \subseteq \tr(A)$,  define the uniform trace norm
$$ \norm{a}_{2, \Delta} = \sup\{ \norm{a}_{2, \tau}: \tau \in \Delta\},\quad a \in A.$$
\end{defn}
The uniform trace norm satisfies
$$ \norm{ab}_{2, \Delta} \leq \min\{ \norm{a} \norm{b}_{2, \Delta}, \norm{a}_{2, \Delta} \norm{b}\} \quad \mathrm{and}\quad \abs{\tau(a)} \leq \norm{a}_{2, \Delta},\quad  a, b \in A,\ \tau \in \Delta.$$ %(see \cite{CETW-Gamma}) $$\tau(a) \leq \norm{a}_{2, \Delta},\quad a \in A,\ \tau \in \Delta.$$

We shall use $l^\infty(A)$ to denote the C*-algebra of bounded sequences of $A$, i.e., $$l^\infty(A) = \{(a_n): a_n \in A,\ \sup\{\norm{a_n}: n=1, 2, ...\} < +\infty \}.$$
Let $\omega$ be a free ultrafilter; then the trace-kernel is the ideal
$$J_{2, \omega, \tr(A)}:=\{(a_n) \in l^\infty(A): \lim_{n \to\omega} \norm{a_n}_{2, \tr(A)} = 0\}.$$

\begin{defn}[Definition 2.1 of \cite{CETW-Gamma}]\label{UGamma-defn}
A unital C*-algebra $A$ is said to have uniform property $\Gamma$ if for each $n\in \mathbb N$, there is a partition of unity $$p_1, p_2, ..., p_n \in (l^\infty(A)/J_{2, \omega, \Delta}) \cap A'$$ such that
$$\tau(p_iap_i) = \frac{1}{n}\tau(a),\quad a\in A,\ \tau\in \tr(A)_\omega,$$
where $\tr(A)_\omega$ denotes the set of limit traces of $l^\infty(A)$, i.e., the traces of the form $$\tau((a_i)) = \lim_{i \to \omega} \tau_i(a_i),\quad \tau_i \in \tr(A),$$
and $a$ is regarded as the constant sequence $(a) \in l^\infty(A)$. 

 
All $\mathcal Z$-absorbing C*-algebras have uniform property $\Gamma$ (see Theorem 5.6 of \cite{CETW-Gamma}). (Indeed, all unital simple  amenable C*-algebras with unique trace have (uniform) property $\Gamma$, and if a C*-algebra $U$ has uniform property $\Gamma$, then the tensor product C*-algebra $A \otimes U$ has uniform property $\Gamma$ (an extreme trace on a tensor product is a product trace).)
\end{defn}



\subsection{AH algebras with diagonal maps}

\begin{defn}
An AH algebra with diagonal maps is the limit of an inductive sequence
$$\xymatrix{
A_1 \ar[r] & A_2 \ar[r] & \cdots \ar[r] & A = \varinjlim A_n,
}$$
where $A_i = \bigoplus_{j} \mathrm{M}_{n_{i, j}}(\mathrm{C}(X_{i, j}))$, and each connecting map preserves the diagonal subalgebras, i.e., it has the form $$f \mapsto \mathrm{diag}\{f\circ\lambda_1, ..., f\circ\lambda_m\},$$ where the $\lambda$s are continuous maps between the $X$s.
\end{defn}

All simple unital AH algebras with diagonal maps have stable rank one (\cite{EHT-sr1}), but not all AH algebras with diagonal maps are $\mathcal Z$-absorbing. In the pioneering work \cite{Vill-perf},  Villadsen constructed simple AH algebras with diagonal maps which have perforation in the ordered $\Kzero$-group. The construction was then used in \cite{Toms-Ann} to obtain a simple AH algebra with diagonal maps which has  the same value of the conventional Elliott invariant as an AI algebra, but is not isomorphic to this AI algebra. Although Villadsen algebras are not $\mathcal Z$-absorbing, a preliminary classification is obtained in \cite{ELN-Vill}.


\subsection{The small boundary property and the mean dimension}

\begin{defn}
A topological dynamical system $(X, \Gamma)$ is free if $x \gamma = x$ implies $\gamma = e$ where $x\in X$ and $\gamma \in \Gamma$. It is said to be minimal if the only closed invariant subspaces of $X$ are $\O$ and $X$. 

A topological dynamical system induces an action of $\Gamma$ on $\mathrm{C}(X)$ by $$\gamma(f)(x) = f(x\gamma),\quad x \in X.$$ We shall assume $\Gamma$ is discrete. The (universal) crossed product C*-algebra $\mathrm{C}(X) \rtimes \Gamma$ is the universal C*-algebra generated by $\mathrm{C}(X)$ and unitaries $u_\gamma$, $\gamma \in \Gamma$, with respect to the relations
$$ u^*_{\gamma} f u_{\gamma} = \gamma(f) \quad \mathrm{and} \quad u_{\gamma_1} u^*_{\gamma_2} = u_{\gamma_1 \gamma^{-1}_2},\quad f\in \mathrm{C}(X),\ \gamma, \gamma_1, \gamma_2 \in \Gamma.$$
\end{defn}

If $\Gamma$ is amenable and the dynamical system $(X, \Gamma)$ is free and minimal, the C*-algebra $\mathrm{C}(X) \rtimes \Gamma$ is simple, unital, amenable, stably finite, and satisfies the UCT. However, the C*-algebra $\mathrm{C}(X) \rtimes \Gamma$ may fail to be $\mathcal Z$-absorbing, even for $\Gamma =\Int$ (\cite{GK-Dyn}).

Let us consider the following property of dynamical systems.
\begin{defn}
A topological dynamical system $(X, \Gamma)$ is said to have the small boundary property (SBP) if for any $x \in X$ and any open neighbourhood $U$ of $x$, there is a neighbourhood $V$ of $x$ such that $V \subseteq U$ and $\mu(\partial V) = 0$ for all invariant measures $\mu$ (\cite{Lindenstrauss-Weiss-MD}).

More generally, consider a metrizable compact space $X$ and a collection $\Delta$ of Borel probability measures on $X$. The pair $(X, \Delta)$ is said to have the (SBP) if for any $x \in X$ and any open neighbourhood $U$ of $x$, there is a neighbourhood $V$ of $x$ such that $V \subseteq U$ and $\mu(\partial V) = 0$ for all $\mu \in \Delta$.  (\cite{EN-SBP-I})
\end{defn}

We have the following criterion for the (SBP):
\begin{thm}[Theorem 2.9 of \cite{EN-SBP-I}]\label{SBP-2-norm}
Let $X$ be a metrizable compact space, and let $\Delta$ be a compact set of Borel probability measures on $X$. Then the pair $(X, \Delta)$ has the (SBP) if, and only if, for any continuous real-valued function $f: X \to \Real$ and any $\eps>0$, there is a continuous real-valued function $g: X\to \Real$ such that
\begin{enumerate}
\item $\norm{f - g}_{2, \Delta} < \eps$, and 
\item there is $\delta>0$ such that $\tau_\mu(\chi_{\delta}(g)) < \eps$, $\mu \in\Delta$,
where $\tau_\mu$ is the tracial state of $\mathrm{C}(X)$ induced by $\mu$, and
\begin{equation*}
\chi_\delta(t) = \left\{ \begin{array}{ll}
1, & \abs{t} < \delta, \\
2-\abs{t}/\delta, & \delta\leq \abs{t} < 2\delta, \\
0, & \mathrm{otherwise}.
\end{array}\right.
\end{equation*}
\end{enumerate} 
\end{thm}


Mean topological dimension was introduced by Gromov (\cite{Gromov-MD}), and then was developed and studied systematically by Lindenstrauss and Weiss (\cite{Lindenstrauss-Weiss-MD}):
\begin{defn}
Consider a topological dynamical system $(X, \Gamma)$, where $\Gamma$ is discrete and amenable. Its mean dimension is defined as
$$\mathrm{mdim}(X, \Gamma):= \sup_{\mathcal U} \lim_{n \to\infty} \frac{1}{\abs{\Gamma_n}}\mathcal D( \bigwedge_{\gamma \in \Gamma_n} \mathcal U \gamma^{-1} ),$$
where $\Gamma_1, \Gamma_2, ...$ is a F{\o}lner sequence of $\Gamma$, the supremum is taken over all finite open covers $\mathcal U$ of $X$, and $\mathcal D(\mathcal U) = \min\{\mathrm{ord}(\mathcal V): \mathcal V \prec \mathcal U\}$ ($\mathrm{ord}$ is the maximal number of the mutually overlapping sets minus $1$).
\end{defn}

By \cite{Lindenstrauss-Weiss-MD}, the small boundary property of $(X, \Gamma)$ implies zero mean dimension. The converse was shown in \cite{MR3614036} and \cite{GLT-Zk} for $\Gamma = \Int^d$, and in \cite{Niu-MD-Zd} for actions with the (URP). 

Zero mean dimension (or small boundary property) implies the $\mathcal Z$-absorption of the C*-algebra:
\begin{thm}[\cite{EN-MD0}]\label{old-thm}
Let $(X, \Int)$ be a free and minimal dynamical system. If $\mathrm{mdim}(X, \Int) = 0$, then the C*-algebra $\mathrm{C}(X) \rtimes \Int$ is $\mathcal Z$-absorbing.
\end{thm}

The main motivation of this work is the converse question of this theorem.
 

\subsection{Uniform Rokhlin property and Cuntz comparison of open sets}

The following two properties were introduced in \cite{Niu-MD-Z}.
\begin{defn}[Definition 3.1 and Definition 4.1 of \cite{Niu-MD-Z}]\label{defn-URP-COS}
A topological dynamical system $(X, \Gamma)$, where $\Gamma$ is a discrete amenable group, is said to have the uniform Rokhlin property (URP) if for any $\eps>0$ and any finite set $K\subseteq \Gamma$, there exist closed sets $B_1, B_2, ..., B_S \subseteq X$ and $(K, \eps)$-invariant sets $\Gamma_1, \Gamma_2, ..., \Gamma_S \subseteq \Gamma$ such that the transformed sets 
$$B_s\gamma,\quad \gamma\in \Gamma_s,\quad s=1, ..., S, $$
are mutually disjoint and  
$$\mathrm{ocap}(X\setminus\bigsqcup_{s=1}^S\bigsqcup_{\gamma\in \Gamma_s}B_s\gamma) < \eps,$$
where the abbreviation $\mathrm{ocap}$ stands for orbit capacity (see, for instance, Definition 5.1 of \cite{Lindenstrauss-Weiss-MD}).

The dynamical system $(X, \Gamma)$ is said to have $(\lambda, m)$-Cuntz-comparison of open sets, where $\lambda\in (0, 1]$ and $m\in \mathbb N$, if for any open sets $E, F\subseteq X$ with $$ \mu(E) < \lambda \mu(F),\quad \mu \in\mathcal M_1(X, \Gamma),$$ 
where $\mathcal{M}_1(X, \Gamma)$ is the simplex of all invariant probability measures on $X$, it follows that 
$$\varphi_E \precsim \underbrace{\varphi_F\oplus\cdots \oplus \varphi_F}_m\quad\mathrm{in}\ \mathrm{C}(X)\rtimes\Gamma,$$ 
where $\varphi_E$ and $\varphi_F$ are continuous functions with open supports $E$ and $F$ respectively.

The dynamical system $(X, \Gamma)$ is said to have Cuntz comparison of open sets (COS) if it has $(\lambda, m)$-Cuntz-comparison on open sets for some $\lambda$ and $m$.
\end{defn}

\begin{thm}[\cite{Niu-MD-Z} and \cite{Niu-MD-Zd}]
Let $(X, \Gamma)$ be a minimal and free dynamical system.
\begin{itemize}
\item If $\Gamma = \Int^d$, then $(X, \Gamma)$ has the (URP) and (COS).
\item If $\Gamma$ is finitely generated and has sub-exponential growth, and if $(X, \Gamma)$ has a Cantor factor, then $(X, \Gamma)$ has the (URP) and (COS).
\end{itemize}
\end{thm}

The (UPR) implies that the C*-algebra $\mathrm{C}(X)\rtimes\Gamma$ can be weakly tracially approximated by the homogeneous C*-algebras generated by the Rokhlin towers. Together with the (COS), it has the following implications for the C*-algebra $\mathrm{C}(X) \rtimes \Gamma$:
\begin{thm}[\cite{Niu-MD-Z}, \cite{Niu-MD-Z-absorbing}, \cite{LN-sr1}]
Let $(X, \Gamma)$ be a minimal and free dynamical system with the (URP) and (COS), and let $A = \mathrm{C}(X) \rtimes\Gamma$. Then:
\begin{itemize}
\item $\mathrm{rc}(A) \leq \frac{1}{2} \mathrm{mdim}(X, \Gamma)$, where $\mathrm{rc}(A)$ is the radius of comparison of $A$; 
\item $A$ has stable rank one, i.e., invertible elements are dense;
\item $A \cong A\otimes\mathcal Z$ if, and only if, $A$ has strict comparison for positive elements; in particular,
\item if $(X, \Gamma)$ has the (SBP), then $A \cong A\otimes\mathcal Z$.
\end{itemize}

\end{thm}

%\subsection{Strong uniform property $\Gamma$ and approximate divisibility}
%
%In contrast to uniform property $\Gamma$, strong uniform property $\Gamma$ and approximate divisibility were introduced in \cite{EN-SBP-I}:
%
%\begin{defn}\label{AD-defn}
%Let $A$ be a C*-algebra and let $D \subseteq A$ be a sub-C*-algebra. The pair $(D, A)$ is said to have  strong uniform property $\Gamma$ if for each $n\in \mathbb N$, there is a partition of unity $$p_1, p_2, ..., p_n \in (l^\infty(D)/J_{2, \omega, \Delta}) \cap A'$$ such that
%$$\tau(p_iap_i) = \frac{1}{n}\tau(a),\quad a\in A,\ \tau\in \tr(A)_\omega,$$
%where $\tr(A)_\omega$ denotes the set of limit traces of $l^\infty(A)$, i.e., the traces of the form $$\tau((a_i)) = \lim_{i \to \omega} \tau_i(a_i),\quad \tau_i \in \tr(A),$$
%and $a$ is regarded as the constant sequence $(a) \in l^\infty(A)$.
%
%
%Consider a commutative C*-algebra $D \cong \mathrm{C}(X)$, and consider a collection $\Delta$ of Borel probability measures on $X$. The pair $(X, \Delta)$ is said to be (tracially) approximately divisible if there is $K>0$ such that for each $n\in \mathbb N$, there is a partition of unity $$p_1, p_2, ..., p_n \in l^\infty(D)/J_{2, \omega, \Delta}$$ such that
%$$\tau(p_iap_i) \leq \frac{1}{n}K\tau(a),\quad a\in D^+,\ \tau\in \Delta_\omega,\ i=1, ..., n.$$
%
%\end{defn}
%
%It is clear that strong uniform property $\Gamma$ for $(D, A)$ implies approximate divisibility of $(D, \tr(A)|_A)$. 
%
%\begin{thm}[\cite{EN-SBP-I}]
%If $(D, \Delta)$ is approximately divisible, then $(D, \Delta)$ has the (SBP).
%\end{thm}
%









%\vskip 1in 
%\begin{thm}\label{the-right-sandwich}
%Let $(D, A)$ be a pair of C*-algebras with Property (C). Let $\phi_1, \phi_2, \psi_1, \psi_2 \in D_1^+$ satisfying the following property:
%\begin{enumerate}
%\item\label{cond-s-1} $\phi_2 \in \overline{\phi_1 D\phi_1}$,
%\item\label{cond-s-2} $\psi_1\psi_2 = \psi_2$,
%\item\label{cond-s-3} $\mathrm{d}_\tau(\psi_1) < \mathrm{d}_\tau(\phi_1)$, for all $\tau\in\tr(A)$,
%\item\label{cond-s-4} $\inf\{\tau(\psi_2) - \mathrm{d}_\tau(\phi_2): \tau \in \tr(A)\} >0 $.
%\end{enumerate}
%Then, for any $\eps>0$, there is a contraction $u\in A$ such that
%$$ u^*\psi_1u \in  \overline{\phi_1A\phi_1},\quad \norm{(u^*\psi_1u) \phi_2 - \phi_2}_{2} < \eps,$$
%$$\mathrm{dist}_{2}(udu^*, D_1) < \eps,\quad \mathrm{dist}_{2}(u^*du, D_1) < \eps,\quad d \in D_1,$$
%and
%$$\norm{uu^* - 1}_{2}, \norm{u^*u - 1}_{2} < \eps.$$
%\end{thm}
%
%\begin{proof}
%Pick $\delta>0$ such that
%$$\delta< \tau(\psi_2) - \mathrm{d}_\tau(\phi_2),\quad \tau \in \tr(A).$$
%
%With the given $\eps$, choose $\eps' \in (0, \eps/2)$ such that if $\norm{ab - b}_2< \eps'$, where $a, b$ are positive contractions, then $$\norm{\eta_{\eps}(a)b - b}_2 < \delta/2$$
%where 
%$$
%\eta_\eps = \left\{
%\begin{array}{ll}
%0, & t \leq 1-\eps/2 \\
%\mathrm{linear}, & t \in [1-\eps/2, 1] \\
%1, & t \geq 1.
%\end{array}
%\right.
%$$
%
%By Property (C) and assumption \ref{cond-s-3}, there is a contraction $u_1\in A$ such that
%$$u^*_1 \psi_1 u_1 \in \overline{\phi_1 A \phi_1}$$
%$$\mathrm{dist}_{2}(u_1du^*_1, D_1) < \eps'/8, \quad \mathrm{dist}_{2}(u_1^*du_1, D_1) < \eps'/8,\quad d \in D_1, $$
%and
%$$\norm{u_1u_1^* - 1}_{2}, \norm{u_1^*u_1 - 1}_{2} < \eps/8.$$
%
%Choose $n$ large enough so that $$ \norm{\phi_1^{\frac{1}{n}}(u_1^*\psi_1u_1) - u_1^*\psi_1u_1}< \eps'/16 \quad \mathrm{and}\quad \norm{\phi_1^{\frac{1}{n}}(u_1^*\phi_1u_1) - u_1^*\phi_1u_1}< \eps'/16.    $$
%Choose positive contractions $d_1, d_2 \in D$ such that
%$$ \norm{d_1- u_1^*\psi_1u_1}_2< \eps'/16 \quad \mathrm{and} \quad \norm{d_1- u_1^*\psi_1u_1}_2 < \eps'/16.$$
%Then
%$$\norm{\phi_1^{\frac{1}{n}}d_1 - u_1^*\psi_1u_1}_2 < \eps'/8 \quad \mathrm{and}\quad \norm{\phi_1^{\frac{1}{n}}d_2 - u_1^*\psi_2u_1}_2 < \eps'/8.$$
%Consider $\phi_1^{\frac{1}{n}}d_1$ and $\phi_1^{\frac{1}{n}}d_1$, and still denoted by $d_1$ and $d_2$ respectively. One has $$d_1, d_2 \in \overline{\phi_1 D \phi_1}$$
%and
%$$\norm{d_1 - u_1^*\psi_1u_1}_2 < \eps'/8 \quad \mathrm{and}\quad \norm{d_2 - u_1^*\psi_2u_1}_2 < \eps'/8.$$
%
%Note that it follows that
%$$\norm{d_1d_2 - d_2}_2 < \eps',$$
%and by the choice of $\eps'$, 
%$$\norm{\eta_\eps(d_1)d_2 - d_2}_2 < \delta/2.$$
%Then $$ \tau(\eta_\eps(d_1)d_2) \approx_{\delta/2} \tau(d_2),\quad \tau\in\tr(A),  $$
%and
%$$  \mathrm{d}_\tau(\eta_\eps(d_1)d_2 ) \geq \tau(\eta_\eps(d_1)d_2) > \mathrm{d}_\tau(\phi_2),\quad \tau\in \tr(A).$$
%
%By Property (C) again, there is a contraction $u_2 \in \overline{\phi_1A\phi_1}+\Comp 1$ such that 
%$$u_2 \phi_2 u^*_2 \in \overline{(\eta_\eps(d_1)d_2) A (\eta_\eps(d_1)d_2)}$$
%$$\mathrm{dist}_{2}(u^*_2du_2, D_1)< \eps/2,\quad \mathrm{dist}_{2}(u_2du^*_2, D_1) < \eps/2,\quad d \in D_1, $$
%and
%$$\norm{u_2u_2^* - 1}_{2}, \norm{u_2^*u_2 - 1}_{2} < \eps/8.$$
%
%Note that
%$$\xi_\eps(d_1) (\eta_\eps(d_1)d_2) = (\eta_\eps(d_1)d_2),$$
%and
%$$\norm{d_1 - \xi_\eps(d_1)} < \eps/2,$$
%where
%$$
%\xi_\eps = \left\{
%\begin{array}{ll}
%0, & t = 0, \\
%\mathrm{linear}, & t \in [0, 1-\eps/2], \\
%1, & t \in[1-\eps/2, 1].
%\end{array}
%\right.
%$$
%One then has
%$$(u_1^*\psi_1 u_1) u_2 \phi_2 u_2^* \approx^{\norm{\cdot}_2}_{\eps'}  d_1 u_2 \phi_2 u_2^* \approx_{\eps/2} \xi_\eps(d_1) u_2 \phi_2 u_2^* = u_2 \phi_2 u_2^*;$$
%that is,
%$$ \norm{ (u_1u_2)^*\psi_1 (u_1u_2) \phi_2 - \phi_2}_2 < \eps. $$
%Since $u_2 \in \overline{\phi_1A\phi_1}+\Comp 1$, $$(u_1u_2)^*\psi_1 (u_1u_2) \in \overline{\phi_1A\phi_1}.$$
%Therefore, $u = u_1u_2$ is the desired unitary.
%\end{proof}
%



%\begin{thm}\label{sandwich}
%Let $A = \mathrm{M}_n(\mathrm{C}(X))$, and let $D$ be the diagonal subalgebra. Let $N \in \mathbb N$ and $\eps>0$.
%
%Assume $f, g, h\in D$ are positive contractions such that $$gh = h,$$ 
%$$ \mathrm{rank}(g(x)) > \mathrm{rank}(f(x)) > \mathrm{Tr}(\eta_\delta(f)(x)) > \mathrm{rank}(h(x)),\quad x\in X$$
%for some $\delta>0$.
%%and
%%$$\mathrm{tr}(\chi_{1-\delta}(f)(x)) > \mathrm{rank}(h(x)),\quad x \in X.$$ 
%Then, with $n$ sufficiently large (depends on $\delta$), there is a unitary $u \in A$ such that 
%$$u^*f u \in \overline{gAg} \quad \mathrm{and} \quad \norm{(u^*f u)h - h}_{2, \tr(A)} < 2\delta,$$
%and
%$$ \mathrm{dist}_{2, \tr(A)}(u^*d u, D) < 4/n,\quad d \in (D)_1. $$
%
%%a positive contraction $a \in \overline{gAg}$ such that
%%$$\mathrm{dist}_2(a, D) = o(\frac{1}{n}),$$
%%$$ \norm{ah - h}_2 < \eps $$
%%and
%%$$ \abs{\mathrm{tr}(a^n(x)) - \mathrm{tr}(f^n(x))} < \eps,\quad i=1, 2, ..., N,\ x\in X. $$
%\end{thm}
%
%
%
%This property can be reduced to the following property:
%
%\begin{lem}\label{lem-comp}
%Let $A = \mathrm{M}_n(\mathrm{C}(X))$, and let $D$ be the diagonal subalgebra. Let $g \in D$ be a positive contraction. Assume that $f, h \in \overline{gDg}$ such that
%$$\mathrm{rank}(h)(x) \leq \mathrm{rank}(h)(x),\quad x \in X,$$
%then there is a unitary $u \in \overline{gAg} + \Comp 1_A$ such that
%$$u^* h u \in \overline{fAf}.$$
%\end{lem}
%
%\begin{proof}[Proof that \ref{lem-comp} implies \ref{sandwich}]
%
%Let $f, g, h$ be given in \ref{sandwich}. Since $$\mathrm{rank}f(x) < \mathrm{rank}g(x),\quad x\in X, $$ by Lemma \ref{lem-comp}, there is a unitary $w_1 \in A$ such that
%$$w_1^*fw_1 \in \overline{gAg}.$$
%
%Choose $ \widetilde{w_1^*fw_1} \in \overline{gDg} $ such that $$ \norm{ \widetilde{w_1^*fw_1} - w_1^*fw_1 }_2 < \frac{3}{n}.$$ With $n$ sufficiently large, one has
%$$ \abs{\tau(\eta_\delta( \widetilde{w_1^*fw_1} )) - \tau(\eta_\delta(w_1^*fw_1 ))} <  \tau(\eta_\delta(f)) - \mathrm{d}_\tau(h),\quad \tau \in \tr(A), $$
%and then
%$$ \mathrm{rank}(h) < \tau(\eta_\delta(f)) < \mathrm{rank}(\eta_\delta( \widetilde{w_1^*fw_1} ))  $$
%
%Therefore, by Lemma \ref{lem-comp} again, there is a unitary $w_2 \in \overline{gAg}+\Comp 1_A$ such that 
%$$w_2 h w_2 \subseteq \overline{(\eta_{\delta}(\widetilde{w_1^*fw_1}) A \eta_{\delta}(\widetilde{w_1^*fw_1}))}.$$
%In particular, 
%$$\widetilde{w^*_1fw_1}(w_2hw_2^*) \approx_\delta w_2hw_2^*,$$
%and hence
%$$\norm{{w^*_1fw_1}(w_2hw_2^*) - w_2hw_2^*}_2 < \delta + {3}/{n} < 2\delta. $$
%That is,
%$$\norm{(w_1w_2)^*f(w_1w_2) - h}_2 < 2\delta.$$
%Since $w_2 \in \overline{gAg}+\Comp 1_A$, one still has $$ (w_1w_2)^*f(w_1w_2)  \in \overline{gAg},$$
%as desired.
%\end{proof}


%\section{The Small Boundary Property}


\section{Property (C)}


\begin{defn}\label{defn-C}
Let $A$ be a unital C*-algebra and let $D$ be a unital commutative sub-C*-algebra. Then the pair $(D, A)$ is said to have Property (C) if for any positive contractions $f, g, h \in D$ satisfying $f, g \in \overline{hDh}$, and $$\mathrm{d}_\tau(f) < \mathrm{d}_\tau(g),\quad \tau \in \tr(A),$$ and for any $\eps>0$, there is a contraction $u \in \overline{hAh} + \Comp 1_A$ such that
$$ u f u^* \in_\eps^{\norm{\cdot}_2} \overline{gAg},$$
$$\mathrm{dist}_{2, \tr(A)}(udu^*, (D)_1) < \eps,\quad \mathrm{dist}_{2, \tr(A)}(u^*du, (D)_1) < \eps, \quad d \in (D)_1, $$ 
and
$$\norm{uu^* - 1}_{2, \tr(A)},\  \norm{u^*u - 1}_{2, \tr(A)} < \eps.$$
\end{defn}

\begin{rem}
Comparing to Property (COS), the approximations in Property (C) are with respect to the uniform trace norm, but on the other hand, the comparison in Property (C) is implemented by an almost unitary which is also an almost normalizer (with respect to the uniform trace norm).
\end{rem}


\begin{thm}\label{SBP=>C}
Let $(X, \Gamma)$ be a minimal and free dynamical system, where $\Gamma$ is a countable discrete amenable group.  Assume that $(X, \Gamma)$ has the (SBP), then $(D, A)$ has the Property (C).
\end{thm}

\begin{proof}
Let $\eps>0$ be arbitrary, and let $f, g, h \in D$ be positive contractions satisfying $f, g \in \overline{hDh}$, and $$\mathrm{d}_\tau(f) < \mathrm{d}_\tau(g),\quad \tau \in \tr(A).$$

Define the open sets
$$E_{(f-\eps)_+} = f^{-1}(\eps, +\infty),\quad E_g = g^{-1}(0, +\infty), \quad \mathrm{and} \quad E_h = h^{-1}(0, +\infty). $$

Then there are finite subset $ K \subseteq \Gamma$ and $\delta>0$ such that if a finite set $\Gamma_0 \subseteq \Gamma$ is $(K, \delta)$-invariant, then
$$ 
\frac{ \abs{x\Gamma_0 \cap E_{(f-\eps)_+}}}{\abs{\Gamma_0}} < \frac{ \abs{x\Gamma_0 \cap E_g}}{\abs{\Gamma_0}}, \quad x \in X.
$$

Since $(X, \Gamma)$ has the (SBP), by \cite{KS-comparison}, it has the (URP), and therefore there are mutually disjoint closed Rokhlin towers 
$$(B_1, \Gamma_1), ..., (B_S, \Gamma_S)$$
such that each $\Gamma_s$, $s=1, ..., S$, is $(K, \delta)$-invariant, and $$\mu(X \setminus \bigsqcup_{s=1}^S \bigsqcup_{\gamma\in \Gamma_s} B_s\gamma) < \eps,\quad \mu \in \mathcal M_1(X, \Gamma).$$

For each $x \in B_s$, $s=1, ..., S$,  since $\Gamma_s$ is $(K, \delta)$-invariant, one has
$$ \abs{x\Gamma_s \cap E_{(f-\eps)_+}} < \abs{x \Gamma_s \cap E_g}.$$ 
Note that $x\Gamma_s \cap E_{(f-\eps)_+}, x\Gamma_s \cap E_{g} \subseteq x\Gamma_s \cap E_{h}$. Therefore, there is a permutation $\sigma_x$ on $x\Gamma_s$ such that its restriction to $x\Gamma_s \cup (E_h)^c$ is identity and
$$\sigma_x(x\Gamma_s \cap E_{(f-\eps)_+}) \subseteq x\Gamma_s \cap E_{g}.$$ 
Since $E_g \subseteq X$ is open, there is a neighborhood $U_x \subseteq X$ such that 
%$$ \gamma \in \Gamma_0\  \textrm{and}\  x\gamma \in E_{(f- \eps)_+}  \quad \Rightarrow \quad U_x \gamma \in E_{(f- \eps)_+}, $$
%and
$$ \gamma \in \Gamma_s\  \textrm{and}\  x\gamma \in E_{g}  \quad \Rightarrow \quad U_x \gamma \in E_{g},$$
and hence
$$\sigma_x(\{\gamma: \gamma \in \Gamma_s,\ x\gamma \in E_{(f-\eps)_+}\}) \subseteq \{\gamma \in \Gamma_s: U_x \gamma \subseteq E_{g}  \}$$
One also assume that $U_x$ is small enough so $(U_x, \Gamma_s)$ is an open tower.

Thus, the permutation unitary matrix $$u_{\sigma_x} \in \mathrm{M}_{\abs{\Gamma_s}}(\Comp) \subseteq (\mathrm{M}_{\abs{\Gamma_s}}(\mathrm{C_0}(U_x))) + \Comp 1$$ has the property that if $\phi_{x}$ is a continuous function supported in $U_x$, then 
\begin{equation}\label{into-u}
u_{\sigma_x} ( \sum_{\gamma \in \Gamma_s} (f - \eps)_+(x\gamma) \gamma(\phi_x) ) u^*_{\sigma_x} \subseteq \overline{gAg},
\end{equation}
where $\sum_{\gamma \in \Gamma_s} (f - \eps)_+(x\gamma) \gamma(\phi_x)$ is regarded as an element of $\mathrm{M}_{\abs{\Gamma_s}}(\mathrm{C_0}(U_x)))$, the C*-algebra of the open tower $(U_x, \Gamma_s)$.
%$g|_{B_s\Gamma_s}$ is the restriction of $g$ to the tower $(B_s, \Gamma_s)$, and regard it as a (diagonal) element of $\mathrm{M}_{\abs{\Gamma_s}}(\mathrm{C}_0(U_x))$, the C*-algebra of generated by the tower $(B_s, \Gamma_s)$. 
Since $u_x$ is a permutation matrix, it preserves the diagonal subalgebra of $\mathrm{M}_{\abs{\Gamma_s}}(\mathrm{C}_0(U_x))$:
$$u_{\sigma_x} (\sum_{\gamma \in \Gamma_s} h \gamma(\phi_x) )u_{\sigma_x} \in \mathrm{C}(X),\quad h \in \mathrm{C}(X). $$

One also assume that $U_x$ is sufficiently small so $$\abs{(f - \eps)_+(y_1) - (f - \eps)_+(y_2)} < \eps,\quad y_1, y_2 \in U_x.$$

Since $B_s$ is compact, there are $x_{s, 1}, x_{s, 2}, ..., x_{s, n_s}$ and open sets $U_{s, 1} \ni x_{s, 1}$, ..., $U_{s, n_s} \ni x_{s, n_s}$ such that $\{U_{s, 1}, ..., U_{s, n_s}\}$ is an open cover of $B_s$. Since $(X, \Gamma)$ has (SBP), there are closed sets $x_{s, 1} \in V_{s, 1} \subseteq U_{s, 1}$, ..., $x_{s, n_s} \in V_{s, n_s} \subseteq U_{s, n_s}$ such that $V_{s, 1}, ..., V_{s, n}$ are mutually disjoint, and
$$\mu(B_s \setminus (V_1 \sqcup \cdots \sqcup V_n)) < \frac{\eps}{\abs{\Gamma_s}},\quad \mu \in \mathcal M_1(X, \Gamma).$$

Note that the tower $(B_s, \Gamma_s)$ splits into thiner towers $(V_{s, 1}, \Gamma_s)$, ..., $(V_{s, n_s}, \Gamma_s)$.

Choose continuous functions $$\phi_{x_i}: X \to [0, 1],\quad i=1, ..., n$$
such that
$$\phi_{x_{s, i}}|_{V_{s, i}} = 1 \quad \mathrm{and} \quad \phi_{x_{s, i}} |_{U_{s, i}^c} = 0,\quad i=1, ..., n,$$
and
$$ \gamma(\phi_{x_{s, i}}),\quad s=1, ..., S,\ i=1, ..., n,\ \gamma \in \Gamma_s$$
are mutually orthogonal.

Then
\begin{eqnarray*}
(f - \eps)_+ & \approx^{\norm{\cdot}_2}_{2 \eps} & (f- \eps)_+ \sum_{s=1}^S\sum_{i=1}^n \sum_{\gamma \in \Gamma_s} \gamma(\phi_{x_{s, i}}) \\
& \approx_\eps & \sum_{s = 1}^S \sum_{i=1}^n \sum_{\gamma \in \Gamma_s}  (f - \eps)_+(x_{s, i}\gamma) \gamma(\phi_{x_{s, i}}). 
\end{eqnarray*}

On each tower $(V_i, \Gamma_s)$, consider the permutation unitary $u_{\sigma_{x_i}} \in \mathrm{M}_{\abs{\Gamma_s}}(\Comp) $, and extend it to $(\tilde{V}_i, \Gamma_s)$, where each $\tilde{V}_i$, $i=1, ..., n$, is an open neighborhood of $V_i$ such that $$(\tilde{V}_i, \Gamma_s),\quad i=1, ..., n_s,\ s=1, ..., S$$ are mutually disjoint (open) towers.

Extend each permutation unitary $u_{x_{s, i}}$, $i=1, ..., n_s$, $s=1, ..., S$, to a contraction of $\mathrm{M}_{\abs{\Gamma_s}}(\mathrm{C}_0(\tilde{V}_{s, i}))$, the C*-algebra of the open tower $(\tilde{V}_{s, i}, \Gamma_s)$, which is regarded as a sub-C*-algebras of $\mathrm{C}(X) \rtimes \Gamma$. Note that, since the towers $(\tilde{V}_i, \Gamma_s),\quad i=1, ..., n_s,\ s=1, ..., S$, are mutually disjoint, the contractions $u_{x_{s, i}}$, $i=1, ..., n_s$, $s=1, ..., S$, are mutually orthogonal. Hence
$$u:= \sum_{s=1}^S\sum_{i=1}^{n_s} u_{x_{s, i}}$$
is a contraction.

Let us show that $u$ has the desired properties: By \eqref{into-u}, 
\begin{eqnarray*}
u(f - \eps)_+u^* & \approx_{2\eps}^{\norm{\cdot}_2} & u \sum_{s=1}^S\sum_{i=1}^{n_s}\sum_{\gamma \in \Gamma_s}(f - \eps)_+(x_{s, i} \gamma) \gamma(\phi_{s, i}) u^* \\
& = & \sum_{s=1}^S \sum_{i=1}^{n_s} u_{x_{s, i}}  (\sum_{\gamma \in \Gamma_s}(f - \eps)_+(x_{s, i} \gamma) \gamma(\phi_{s, i})) u^*_{x_{s, i}} \in \overline{gAg}.
\end{eqnarray*}

The element $u$ is also an almost normalizer: Indeed, for any contraction $c \in \mathrm{C}(X)$, one has
\begin{eqnarray*}
u c u^* & \approx_{2\eps}^{\norm{\cdot}_2} & u \sum_{s=1}^S\sum_{i=1}^{n_s} \sum_{\gamma \in \Gamma_s} c \gamma(\phi_{s, i}) u^* \\
& = & \sum_{s=1}^S \sum_{i=1}^{n_s} u_{x_{x_{s, i}}} ( \sum_{\gamma \in \Gamma_s}  c \gamma(\phi_{s, i})) u_{x_{s, i}}^* 
 \in  \mathrm{C}(X).
\end{eqnarray*}
Applying the equation above with $c = 1$ shows that
$$uu^* \approx_{4\eps}^{\norm{\cdot}_2} 1, $$
The similar argument shows that
$$u^*c u \in_{2\eps}^{\norm{\cdot}_2} \mathrm{C}(X), \quad c\in \mathrm{C}(X),\ \norm{h}\leq 1$$
and
$$u^*u \approx_{4\eps}^{\norm{\cdot}_2} 1,$$
as desired.
\end{proof}


\section{Property (S)}

Motivated by Theorem \ref{SBP-2-norm}, let us introduce the following property of a C*-algebra.
\begin{defn}\label{defn-S}
Let $A$ be a unital C*-algebra, and let $\Delta \subseteq \tr(A)$ be a closed set of tracial states. The pair $(A, \Delta)$ will be said to have Property (S) if for any self-adjoint contraction $f$ and any $\eps>0$, there is a self-adjoint contraction $g \in A$ such that 
\begin{enumerate}
\item $\norm{f - g }_{2, \Delta} < \eps$, and
\item there is $\delta>0$ such that $\tau(\chi_\delta(g)) < \eps$, $\tau \in \Delta$,
where 
\begin{equation}\label{chi-defn}
\chi_\delta(t) = \left\{ \begin{array}{ll}
1, & \abs{t} < \delta, \\
2-\abs{t}/\delta, & \delta\leq \abs{t} < 2\delta, \\
0, & \mathrm{otherwise}.
\end{array}\right.
\end{equation}
\end{enumerate}

In the case that $\Delta = \tr(A)$, we shall just say that $A$ has Property (S) if $(A, \tr(A))$ has Property (S).
\end{defn}

Compared to Theorem \ref{SBP-2-norm}, Property (S) can be regarded as a weaker version of the small boundary property, without referring to a commutative subalgebra $D$. 

In general, it is weaker than the actual small boundary property, as shown in Example \ref{weak-S} later in this section. However, it will be shown  (Theorem  \ref{prop-S}) that, if some other properties (Properties (C) and (E)) are provided, then Property (S) for $A$ implies the small boundary property of the pair $(D, A)$.



Property (S) is closely related to the real rank of the sequence algebra:
\begin{defn}\label{qRR-defn}
Let $A$ be a C*-algebra, and let $\Delta\subseteq\tr(A)$. Let us say that $\mathrm{qRR}(l^\infty(A)/J_{2, \omega, \Delta}) = 0$ if the class of the constant sequence of any self-adjoint contraction of $A$ can be approximated by invertible self-adjoint contractions of $l^\infty(A)/J_{2, \omega, \Delta}$ with respect to the uniform limit trace norm $\norm{\cdot}_{2, \omega, \Delta}$. It is clear that if $\mathrm{RR}(A) = 0$ or if $\mathrm{RR}(l^\infty(A)/J_{2, \omega, \Delta}) = 0$, then $\mathrm{qRR}(l^\infty(A)/J_{2, \omega, \Delta}) = 0$.

\end{defn}


\begin{prop}\label{character-S}
Let $A$ be a unital C*-algebra, and let $\Delta\subseteq\tr(A)$ be closed. The following conditions are equivalent:
\begin{enumerate}
\item[(1)] $\mathrm{qRR}(l^\infty(A)/J_{2, \omega, \Delta}) = 0$;
\item[(2)] $(A, \Delta)$ has Property (S);
\item[(3)] $\mathrm{RR}(l^\infty(A)/J_{2, \omega, \Delta}) = 0$.
\end{enumerate}
\end{prop}

\begin{proof}
$(1) \Rightarrow (2)$: Assume $\mathrm{qRR}(l^\infty(A)/J_{2, \omega, \Delta}) = 0$.  Let $(f, \eps)$ be given, where $f \in A$ is a self-adjoint contraction and $\eps>0$. Consider the constant sequence $(f)$, and then there is a sequence of self-adjoint contractions $(g_k) \in l^\infty(A)$ such that 
$$\norm{(f - g_k)}_{2, \omega, \Delta} = \lim_{k \to \omega}\norm{f - g_k}_{2, \Delta} < \eps$$
and $\overline{(g_k)}$ is invertible in $l^\infty(A)/J_{2, \omega, \Delta}$. Then there is $\delta>0$ such that
$$\overline{(\chi_\delta(g_k))} = \chi_\delta(\overline{(g_k)}) = 0.$$ 
In other words,
$$ (\chi_\delta(g_k)) \in J_{2, \omega, \Delta}.$$
Then, with some sufficiently large $k$, one has
\begin{itemize}
\item $\norm{f - g_k}_{2, \Delta} < \eps$, and
\item $\tau(\chi_\delta(g_k)) < \eps$, $\tau\in \Delta$,
\end{itemize}
as desired.

$(2) \Rightarrow (3)$: Assume that $(A, \Delta)$ has Property (S), and let us show that $l^\infty(A)/ J_{2, \omega, \Delta}$ has real rank zero. Let $\overline{(a_1, a_2, ...)} \in l^\infty(A)/ J_{2, \omega, \Delta}$ be a self-adjoint contraction (without loss of generality, one may assume that each $a_n$, $n=1, 2, ...$, is a contraction), and let $\eps>0$ be arbitrary. By Property (S), for each $n=1, 2, ...$, there is a self-adjoint contraction $b_n$ such that
\begin{itemize}
\item $\tau((a_n - b_n)^2) < 1/n$ for all $\tau\in\Delta$, and
\item there is $\delta_n>0$ such that $\tau(\chi_{\delta_n}(b_n)) < 1/n$ for all $\tau\in\Delta$.
\end{itemize}
Without loss of generality, one may assume that $\delta_n<\eps$.
Note that $$\overline{(a_1, a_2, ...)} = \overline{(b_1, b_2, ...)}.$$

Consider $b_n' = g_n(b_n)$ and $c_n = h_n(b_n)$, where
$$
g_n(t) = \left\{ \begin{array}{ll}
\eps t/\delta_n, & \abs{t} < \delta_n, \\
t + \mathrm{sign}(t)(\eps-\delta_n), & \mathrm{otherwise}
\end{array}\right.
$$
and
$$
h_n(t) = \left\{ \begin{array}{ll}
t/\eps\delta_n, & \abs{t} < \delta_n, \\
1/(t + \mathrm{sign}(t)(\eps-\delta_n)), & \mathrm{otherwise}.
\end{array}\right.
$$
Then 
$$\norm{b_n' - b_n} < \eps,\quad \norm{c_n} < 1/\eps,\quad \mathrm{and}\quad \tau((b_n'c_n - 1)^2) <\tau(\chi_{\delta_n}(b_n))< 1/n,\quad \tau \in \Delta.$$ In particular, the element $\overline{(b'_1, b'_2, ...)}$ is invertible in $l^\infty(A)/J_{\omega, \Delta}$ with the inverse $\overline{(c_1, c_2, ...)}$. Moreover, $$\norm{\overline{(a_1, a_2, ...)} - \overline{(b'_1, b'_2, ...)}} = \norm{\overline{(b_1, b_2, ...)} - \overline{(b'_1, b'_2, ...)}}< \eps.$$ This shows that $l^\infty(A)/ J_{2, \omega, \Delta}$ has real rank zero.

$(3) \Rightarrow (1)$: Trivial.
\end{proof}


Uniform property $\Gamma$ implies Property (S):
\begin{prop}\label{Property-S-Prop}
If a unital C*-algebra $A$ has uniform property $\Gamma$, then $A$ has Property (S).
\end{prop}

\begin{proof}
The proof is similar to the proof of Theorem 3.5 of \cite{EN-SBP-I}:

Let $(f, \eps)$ be given, where $f \in A$ is a self-adjoint contraction and $\eps>0$. By Corollary 3.9 of \cite{EN-SBP-I}, 
for the given $\eps$, there exist $n\in\mathbb N$ and self-adjoint contractions $$f_1, f_2, ..., f_n \in A$$  such that
\begin{equation}\label{mult-pert-eq-1}
\norm{f - f_i} < \eps,\quad i=1, 2, ..., n, 
\end{equation}
and 
there is $\delta>0$ such that 
\begin{equation*}
\frac{1}{n}(\tau((f_1)_\delta) + \cdots + \tau((f_n)_\delta)) < \eps,\quad \tau\in\mathrm{T}(A),
\end{equation*} 
where $(f)_\delta = \chi_\delta(f)$ (see \eqref{chi-defn}). 
In particular, regarding $(f_1)_\delta, ..., (f_n)_\delta$ as constant sequences in $A$, we have
\begin{equation}\label{mult-pert-eq-2}
\frac{1}{n}(\tau((f_1)_\delta) + \cdots + \tau((f_n)_\delta)) < \eps,\quad \tau\in\tr(A)_\omega.
\end{equation} 

Since $A$ has uniform property $\Gamma$, there is a partition of unity $$p_1, p_2, ..., p_n \in ( l^\infty(A)/J_{2, \omega, \Delta}) \cap A' $$
such that
\begin{equation}\label{r-gamma-cut-eq}
\tau(p_iap_i) = \frac{1}{n} \tau(a),\quad a\in A,\ \tau\in \tr(A)_\omega.
\end{equation}

Consider the contraction $$g:=p_1\overline{(f_1)}p_1 + \cdots + p_n\overline{(f_n)}p_n \in l^\infty(A)/J_{2, \omega, \Delta}. $$
By \eqref{mult-pert-eq-1},
\begin{equation}\label{close-cond-1}
\norm{f-g} = \norm{p_1\overline{(f-f_1)}p_1 + \cdots + p_n\overline{(f-f_n)}p_n}  < \eps.
\end{equation}
Note that, for each $\tau \in \tr(A)_\omega$, by \eqref{r-gamma-cut-eq},
$$\tau(p_i\overline{((f_i)_\delta)} p_i) = \frac{1}{n}\tau((f_i)_\delta),\quad i=1, ..., n, \ \tau\in\tr(A)_\omega,$$
and hence, together with \eqref{mult-pert-eq-2},
\begin{eqnarray}\label{close-cond-2}
\tau((g)_\delta) & = & \tau(p_1\overline{((f_1)_\delta)}p_1) + \cdots + \tau(p_n\overline{((f_n)_\delta)}p_n) \\
& = & \frac{1}{n}(\tau((f_1)_\delta) + \cdots + \tau((f_n)_\delta)  ) \nonumber \\
& < & \eps.  \nonumber
\end{eqnarray}
Pick a representative sequence $g = \overline{(g_k)}$ with $g_k$, $k=1, 2, ...$, self-adjoint contractions of $A$. By \eqref{close-cond-1} and \eqref{close-cond-2}, with some sufficiently large $k$, the function $g_k$ satisfies
\begin{enumerate}
\item $\norm{f - g_k}_{2, \tr(A)} < \eps$, and
\item $\tau((g_k)_\delta) < \eps$, $\tau\in \tr(A)$,
\end{enumerate}
as desired.
\end{proof}

But Property (S) in general is a weaker property than the actual small boundary property:
\begin{example}\label{weak-S}
Let
$$
\xymatrix{
A_1 \ar[r] & A_2 \ar[r] & \cdots \ar[r] & A}
$$
be a unital AH algebra such that $A_s = \mathrm{M}_{n_s}(\mathrm{C}([0, 1]^{d_s}))$, and $$ \lim_{s \to\infty}n_s = \infty \quad \mathrm{and}\quad  \limsup_{s \to \infty}\frac{d_s}{n_s} = \gamma < \infty.$$ Such C*-algebras include Villadsen algebras which are not $\mathcal Z$-absorbing (so the standard commutative sub-C*-algebra does not have the small boundary property).

Note that $\mathrm{M}_n(\Comp)^{s.a.}$ is a manifold with of dimension $n^2$ and $\mathrm{U}_n(\Comp)$ is a compact Lie group of dimension $n^2$.

Denote by $V_{n, k}$ the set of self-adjoint $n \times n$ matrices with at least $k$ eigenvalues which are zero. It can be written as
$$ V_{n, k}=\{u\mathrm{diag}\{\underbrace{0, ..., 0}_k, \lambda_1, ..., \lambda_{n-k}\}u^*: u \in \mathrm{U}_n(\Comp),\ \lambda_1, ..., \lambda_{n-k} \in \Real\},$$
which is the $\mathrm{U}_n(\Comp)$-orbit of the following $(n-k)$ dimensional submanifold: $$\{\mathrm{diag}\{\underbrace{0, ..., 0}_k, \lambda_1, ..., \lambda_{n-k}\},\quad \lambda_1, ..., \lambda_{n-k} \in \Real\} \subseteq \mathrm{M}_n(\Comp)^{s.a.}.$$ 
For each diagonal matrix above, its stabilizer contains the subgroup
$$
\left\{
\left( 
\begin{array}{cc}
U & 0 \\
0 & I_{n-k}
\end{array}
\right): U \in \mathrm{U}_k(\Comp)
\right\} \subseteq \mathrm{U}_n(\Comp),
$$
which has dimension $k^2$.
The set $V_{n, k}$ is a union of finitely many submanifolds of $ \mathrm{M}_n(\Comp)^{s.a.}$ with dimension at most $$ (n - k) + n^2 - k^2.$$

Now, let $\eps>0$ be arbitrary, and let $a \in A_1$ be a self-adjoint element.
Choose $s$ large enough that if $A_s$ is written as $\mathrm{M}_{n}(\mathrm{C}([0, 1]^{d}))$, then
$$ \frac{d}{n} \leq \gamma+1$$
and
$$
(1 - \eps^2) n^2 + 2\eps n +  (1-\eps)n + (\gamma+1) n + 1 \leq n^2.
$$
Choose a natural number $k$ such that
$$
\frac{k}{n} \leq \eps < \frac{k+1}{n}.
$$

Note that the set $V_{n, k}$ has dimension at most
$$(n-k) + n^2 - k^2 \leq (n - \eps n +1) + n^2 - (\eps n-1)^2.$$
Then 
\begin{eqnarray*}
\mathrm{dim}(V_{n, k}) + d + 1 & \leq & (n - \eps n +1) + (n^2 - (\eps n-1)^2) + (\gamma + 1) n + 1 \\
& = & (1-\eps^2)n^2 +2\eps n  + (1-\eps)n + (\gamma+ 1)n+ 1 \\
&\leq & n^2 = \mathrm{dim}(\mathrm{M}_n(\Comp)^{s.a.}).
\end{eqnarray*}

Hence, the self-adjoint element $a$, regarded as a continuous map from $[0, 1]^d$ to $\mathrm{M}_n(\Comp)^{s.a.}$, can be approximated by continuous maps $b$ from $[0, 1]^d$ to $\mathrm{M}_n(\Comp)^{s.a.} \setminus V_{n, k}$, and so at most $k$ of the eigenvalues of $b(x)$ can be zero at every $x \in [0, 1]^d$, and hence there is $\delta_x>0$ such that at most $k$ of the eigenvalues of $b(x)$ are in the $\delta_x$-neighbourhood of $0$. Then, by  the continuity of the function $b$, a compactness argument shows that there is $\delta>0$ such that for every $x\in [0, 1]^d$, there at most $k$ of the eigenvalues of $b(x)$ are in the $\delta$-neighbourhood of $0$, and hence
$$ \tau(\chi_{\delta/2}(b)) < \frac{k}{n} < \eps,\quad \tau \in \tr(A_s).$$
So the C*-algebra $A$ has Property (S).


\end{example}

In the recent work \cite{Vaccaro:2026aa}, the example above is generalized to $\mathrm{C}(X)\rtimes \Gamma$:
\begin{thm}[\cite{Vaccaro:2026aa}]\label{S-holds}
Let $A \mathrm{C}(X) \rtimes \Gamma $, where $(X, \Gamma)$ is a free and minimal dynamical system with the (URP), or let $A$ be an AH algebra with diagonal maps. Then $A$ has the Property (S).
\end{thm}

\section{An existence property}








%\subsection{Property (C)}

%Let us consider Property (C) now.

%\begin{lem}
%Let $(X, \Gamma)$ be a free and minimal dynamical system. Let $E \subseteq X$ be a closed set and let $F \subseteq X$ be an open set such that 
%$$ \mu(E) < \mu(F),\quad \mu \in \mathcal M_1(X, \Gamma).$$ 
%Then, for any $\eps>0$, there exist open sets $U_1, ..., U_n$, and $\gamma_1, ..., \gamma_n \in \Gamma$ such that
%\begin{enumerate}
%\item $U_i \gamma_i \subseteq F$, $i=1, ..., n$,
%\item $\mu(E \setminus \bigcup_{i=1}^n U_i) < \eps$, $\mu \in \mathcal M_1(X, \Gamma)$, and
%\item $ \mu(\bigcup_{i=1}^n (U_i \setminus \bigcup_{j\neq i} U_j)) < \eps$, $\mu \in \mathcal M_1(X, \Gamma)$.
%\end{enumerate}
%\end{lem}
%
%\begin{proof}
%List $$\Gamma = \{\gamma_1, \gamma_2, ...\}.$$ Choose an open set $$V_1 \supseteq E$$ such that $$\mu(\overline{V_1}) - \mu(E) < \eps/2. $$
%
%Set
%$$E_1 = E \setminus (V_1 \cap F\gamma_1^{-1}) = E\setminus F\gamma^{-1} \quad \mathrm{and} \quad F_1 = F \setminus \overline{V_1}\gamma_1. $$
%
%
%Choose an open set
%$$V_2 \supseteq E_1$$
%such that
%$$ \mu(\overline{V_2}) - \mu(E_1) < \eps/2^2,\quad \mu \in \mathcal M_1(X, \Gamma).$$
%Set
%
%$$E_2 = E_1 \setminus (V_2 \cap F_1 \gamma^{-1}_2) = E_1 \setminus F_1\gamma_2^{-1} \quad \mathrm{and} \quad F_2 = F_1 \setminus \overline{V_2} \gamma_2. $$
%
%In general, with $E_k, F_k$ constructed,
%choose an open set $$V_{k+1} \supseteq E_k$$
%such that
%$$ \mu(V_{k+1}) - \mu(E_k) < \eps/2^{k+1},\quad \mu \in \mathcal M_1(X, \Gamma).$$
%Then, define
%$$E_{k+1} = E_k \setminus (V_{k+1} \cap F_k\gamma^{-1}_{k+1}) = E_k \setminus F_k \gamma_k^{-1} \quad \mathrm{and} \quad F_{k+1} = F_k \setminus \overline{V_{k+1}}\gamma_{k+1}. $$
%
%Consider $$E_\infty = \bigcap_{k=1}^\infty E_k \quad \mathrm{and}\quad  F_\infty = \bigcap_{k=1}^\infty F_k. $$
%Then $$\mu(E_\infty) = 0,\quad \mu \in \mathcal M_1(X, \Gamma).$$
%
%By Dini's Theorem, there is $n$ such that
%$$ \mu(E_n) < \eps,\quad \mu \in \mathcal M_1(X, \Gamma).$$
%
%Consider the open sets
%$$ V_1 \cap F\gamma^{-1}_1,\ V_2 \cap F_1 \gamma_2^{-1},\ ..., V_n \cap F_{n-1} \gamma_n^{-1}, $$
%which cover $$E \setminus E_n. $$
%
%%Let $$x \in (V_1 \cap F\gamma_1^{-1}) \cap (V_2 \cap F_1\gamma_2^{-1}).$$ Since $$x \in V_2 \cap F_1\gamma_2^{-1},$$
%%one has
%%$$x \in V_2 \supseteq E_1.$$ If $x \in E_1 = E \setminus (V_1 \cap F\gamma_1^{-1})$, then $x \notin V_1 \cap F\gamma_1^{-1}$, which is not possible. So $x \in V_2 \setminus E_1$.
%
%For each $i < j$, 
%$$ (V_i \cap F_{i-1} \gamma_i^{-1}) \cap (V_j \cap F_{j-1} \gamma_j^{-1}) \subseteq  V_j \setminus E_{j-1}.    $$
%Indeed, let $$x \in V_j \cap F_{j-1} \gamma_j^{-1}. $$ 
%If 
%\begin{eqnarray*}
%x \in E_{j-1} & = & E_{j-2} \setminus (V_{j-1} \cap F_{j-2} \gamma_{j-1}^{-1}) \\
%& = & (E_{j-3} \setminus (V_{j-2} \cap F_{j-3}\gamma_{j-2}^{-1})) \setminus (V_{j-1} \cap F_{j-1} \gamma_{j-1}^{-1})  \\
%& = & (\cdots (E_{i-1} \setminus (V_i \cap F_{i-1} \gamma_i^{-1})) \setminus \cdots )\setminus (V_{j-1} \cap F_{j-1} \gamma_{j-1}^{-1}),
%\end{eqnarray*}
%then $x \notin V_i \cap F_{i-1} \gamma_i^{-1})$, which is a contradiction. So
%$$x \in V_{j} \setminus E_{j-1}.$$
%Therefore,
%$$ \mu(\bigcup_{i \neq j} (V_i \cap F_{i-1} \gamma_i^{-1}) \cap (V_j \cap F_{j-1} \gamma_j^{-1})) \subseteq \eps + \eps/2 + \cdots = 2\eps,\quad \mu \in \mathcal M_1(X, \Gamma).$$
%Then
%$$ U_1:=V_1 \cap F\gamma^{-1}_1,\ U_2:=V_2 \cap F_1 \gamma_2^{-1},\ ..., U_n:=V_n \cap F_{n-1} \gamma_n^{-1}, $$
%have the desired property of the lemma.
%\end{proof}
%
%\begin{thm}
%Let $(X, \Gamma)$ be a free and minimal dynamical system, where $\Gamma$ is a countable discrete group. Then $(D, A)$ has Property (C).
%\end{thm}
%
%\begin{proof}
%Let $f, g, h \in \mathrm{C}(X)$ be positive elements such that $f, g \in \mathrm{Her}(h)$ and
%$$\mathrm{d}_\tau(f) < \mathrm{d}_\tau(g),\quad \tau \in \mathcal{M}_1(X, \Gamma).$$ Let $\eps>0$ be arbitrary.
%
%Set $$E = f^{-1}((0, +\infty)), \quad F = g^{-1}((0, +\infty)),\quad \mathrm{and} \quad K = h^{-1}((0, +\infty)). $$ Then, $$ E, F \subseteq K. $$
%\end{proof}



%\begin{thm}\label{sandwich-global}
%Let $A$ be a simple AH algebra with diagonal maps, or let $A = \mathrm{C}(X) \rtimes \Gamma$, where $(X, \Gamma)$ is free minimal with the (URP). 
%
%Assume $f, g, h\in D$ are positive contractions such that $$gh = h,$$ 
%$$ \mathrm{d}_\tau(g) > \mathrm{d}_\tau(f) > \tau(\eta_\delta(f)) > \mathrm{d}_\tau(h),\quad \tau \in \tr(A)$$
%for some $\delta>0$, and
%$$\gamma:=\inf\{\tau(\eta_\delta(f)) - d_\tau(h): \tau \in \tr(A)\} >0.$$
%
%Then, for any $\eps>0$, there is a contraction $u \in A$ such that
%$$ \mathrm{dist}_{2, \tr(A)}(u^*fu, \overline{gAg})<\eps, \quad \norm{(u^*fu) h - h}_{2, \tr(A)} < 2\delta,$$
%$$\mathrm{dist}_{2, \tr(A)}(u^*du, D_1) < \eps,\quad d \in (D)_1,$$
%and
%$$\norm{uu^* - 1}_{2,\tr(A)}, \norm{u^*u - 1}_{2,\tr(A)} < \eps.$$
%\end{thm}
%
%\begin{proof}[Proof of \ref{sandwich-global} assuming \ref{lem-comp-global}]
%
%Let $f, g, h$ be given in \ref{sandwich-global}, and let $\eps$ be arbitrary. Pick $\eps' >0$ such that if $a, b$ are positive contractions satisfying $$\norm{a-b}_{2, \tr(A)} < \eps',$$ then 
%$$\abs{\tau(\eta_\delta(a)) - \tau(\eta_\delta(b))} < \gamma,\quad \tau \in \tr(A).$$ Without loss of generality, one may assume that $\eps' < \min\{\eps/2, \delta\}$.
%
%
%Since $$\mathrm{d}_\tau(f) < \mathrm{d}_\tau(g),\quad \tau\in\tr(A), $$ by Lemma \ref{lem-comp-global}, there is a unitary $w_1 \in A$ such that
%$$\mathrm{dist}_{2, \tr(A)}(w_1^*fw_1, \overline{gAg}) <\eps'/4<\eps,$$
%and
%$$\mathrm{dist}_{2, \tr(A)}(w_1^*dw_1, D) < \eps'/4,\quad d \in (D)_1.$$
%
%Choose $n$ large enough so that $$\norm{g^{\frac{1}{n}} (w_1^*fw_1) - (w_1^*fw_1)} < \eps'/4$$
%and choose $d \in D_1$ such that $$\norm{d - w^*_1fw_1}_{2, \tr(A)} < \eps'/4.$$ 
%Then
%$$g^{\frac{1}{n}}dg^{\frac{1}{n}} \approx_{\eps'/4} g^{\frac{1}{n}} (w^*_1fw_1) g^{\frac{1}{n}} \approx_{\eps'/2} w^*_1fw_1,$$
%where the approximations are in $\norm{\cdot}_{2, \tr(A)}$.
%
%Denote by $\widetilde{w_1^*fw_1}:= g^{\frac{1}{n}}dg^{\frac{1}{n}} \in \overline{gDg}$, one then has
%$$ \norm{ \widetilde{w_1^*fw_1} - w_1^*fw_1 }_{2, \tr(A)} < \eps'.$$ By the choice of $\eps'$, one has
%$$ \abs{\tau(\eta_\delta( \widetilde{w_1^*fw_1} )) - \tau(\eta_\delta(w_1^*fw_1 ))} <  \tau(\eta_\delta(f)) - \mathrm{d}_\tau(h),\quad \tau \in \tr(A), $$
%and then
%$$ \mathrm{d}_\tau(h) < \mathrm{d}_\tau(\eta_\delta( \widetilde{w_1^*fw_1} )),\quad \tau\in \tr(A).$$
%
%Therefore, by Lemma \ref{lem-comp-global} again, there is a unitary $w_2 \in \overline{gAg}+\Comp 1_A$ such that 
%$$w_2 h w_2 \subseteq \overline{(\eta_{\delta}(\widetilde{w_1^*fw_1}) A \eta_{\delta}(\widetilde{w_1^*fw_1}))},$$
%and 
%$$\mathrm{dist}_{2, \tr(A)}(w_2dw_2^*, D) < \eps/2,\quad d \in (D)_1.$$
%In particular, 
%$$\widetilde{w^*_1fw_1}(w_2hw_2^*) \approx_\delta w_2hw_2^*,$$
%and hence
%$$\norm{{w^*_1fw_1}(w_2hw_2^*) - w_2hw_2^*}_2 < \eps'+\delta < 2\delta. $$
%That is,
%$$\norm{(w_1w_2)^*f(w_1w_2) h - h}_2 < 2\delta.$$
%Since $w_2 \in \overline{gAg}+\Comp 1_A$, one has
%$$ \mathrm{dist}_{2,\tr(A)}((w_1w_2)^*f(w_1w_2), \overline{gAg}) <\eps'+ \mathrm{dist}(w^*_2(\widetilde{w^*_1fw_1})w_2, \overline{gAg}) = \eps'<\eps. $$
%%still has $$ (w_1w_2)^*f(w_1w_2)  \in \overline{gAg}.$$
%Also note that
%$$\mathrm{dist}_{2, \tr(A)}((w_1w_2)^*d(w_1w_2), D) < \eps'+\eps/2<\eps,\quad d \in D,$$
%as desired.
%\end{proof}

%If a pair $(D, A)$ has Property (C), it then 


The following definition is an existence property for a pair of C*-algebras $D \subseteq A$.

\begin{defn}\label{defn-D}
Let $A$ be a unital C*-algebra and let $D$ be a unital commutative subalgebra. 

The pair $(D, A)$ is said to have Property (E) if for any positive contraction $a \in A$, its the trace spectral distribution can be approximated by the trace spectral distribution of the positive contractions of $D$ in the sense that for any finite subset $\mathcal F \subseteq \mathrm{C}([0, 1])$, and any $\eps>0$, there is a positive contraction $b \in D$ such that
$$ \abs{ \tau(f(a)) - \tau(f(b))} < \eps,\quad f \in \mathcal F,\ \tau \in \tr(A).$$
\end{defn}


This property holds for all AH algebras with diagonal maps and all crossed products $\mathrm{C}(X) \rtimes \Gamma$ where $(X, \Gamma)$ is a minimal and free dynamical system with the (URP). Recall the following theorem which is due to Thomsen and Li (\cite{Thomsen-AI-Tr} and \cite{Li-interval}).
\begin{thm}\label{Thomsen-Li}
Suppose that $X$ is a path connected metrizable compact space. For any finite subset $\mathcal F \subseteq \mathrm{Aff}\mathrm{T}(\mathrm{C}(X))$ and $\eps>0$, there is $N > 0$ with the following property:

For any unital positive linear map $\xi: \mathrm{Aff}\mathrm{T}(\mathrm{C}(X)) \to \mathrm{Aff}\mathrm{T}(\mathrm{C}(Y))$, where $Y$ is an arbitrary metrizable compact space, and any $n\geq N$, there are homomorphisms 
$$ \phi_1, ..., \phi_n: \mathrm{C}(X) \to  \mathrm{C}(Y) $$
such that
$$\abs{\xi(f)(\tau) -  \frac{1}{n} \sum_{i=1}^n \tau(\phi_i(f)) } < \eps,\quad f \in \mathcal F,\ \tau \in \mathrm{T}(\mathrm{C}(Y)). $$
\end{thm}

With Property (URP), the crossed product C*-algebra $\mathrm{C}(X) \rtimes \Gamma$ has the following tracial approximation property:
\begin{lem}[cf.~Theorem 3.9 of \cite{Niu-MD-Z}]\label{pre-approx}
Let $(X, \Gamma)$ be a free and minimal dynamical system with the (URP). Then, for any finite set $ \mathcal F \subseteq \mathrm{C}(X)$ and any $\eps>0$, there exist a positive element $p \in \mathrm{C}(X)$ with $\norm{p} = 1$ and a sub-C*-algebra $C \subseteq \mathrm{C}(X) \rtimes\Gamma$ such that
\begin{enumerate}
\item $ \norm{ [p, f] } < \eps$, $ f \in \mathcal F$,
\item $pfp \in_\eps C $, $f \in \mathcal F$, 
\item $pdp \in C$, $d \in \mathrm{C}(X)$,  
\item $C \cong \bigoplus_{i=1}^S \mathrm{M}_{n_i}(\mathrm{C}_0(Z_i))$, where $Z_i \subseteq X$, $i=1, ..., S$, are mutually disjoint, and under this isomorphism, the elements $pdp$ are  diagonal elements of $C$ for all $d \in \mathrm{C}(X)$,
\item under the isomorphism above, all diagonal elements of $\bigoplus_{i=1}^S \mathrm{M}_{n_i}(\mathrm{C}_0(Z_i))$ are in $C \cap \mathrm{C}(X)$,
\item $\mathrm{d}_\tau(1-p) < \eps$, $\tau \in \tr(A)$, 
\item there is a closed subset $[Z_{i}] \subseteq Z_i$ for each $i=1, ..., S$ such that if $a \in C$, $\norm{a} \leq 1$, and $a$ is supported in $\bigsqcup_{i=1}^S (Z_i \setminus [Z_{i}])$ under the isomorphism $C \cong \bigoplus_{i=1}^S\mathrm{M}_{n_i}(\mathrm{C}_0(Z_i))$, then $$ \tau(a) < \eps, \quad \tau \in \tr(A),$$
\item for any $d \in \mathrm{C}(X)$, there are $d_1, d_2 \in \mathrm{C}(X)$ such that  $$d = d_1 + d_2,\quad \norm{d_1}_{2, \tr(A)} < \eps,\quad d_2 \in C \cap \mathrm{C}(X), $$ and the support of $d_2$ is inside $\bigsqcup_{i=1}^S [Z_i]$ under the isomorphism $C \cong \bigoplus_{i=1}^S \mathrm{M}_{n_i}(\mathrm{C}_0(Z_i))$.
\end{enumerate}
\end{lem}
\begin{proof}
This follows from the same argument as Theorem 3.9 of \cite{Niu-MD-Z} (or, actually a simpler one).
\end{proof}


\begin{thm}\label{Property-D}
Let $A$ be a simple AH algebra with diagonal maps, or let $A = \mathrm{C}(X) \rtimes \Gamma$, where $(X, \Gamma)$ is a free and minimal dynamical system with the (URP). Then the pair $(D, A)$ has Property (E).
\end{thm}

\begin{proof}
Let $(\mathcal F, \eps)$ be given. Without loss of generality, one may assume that each element of $\mathcal F$ has norm $1$ and is real valued,  so $\mathcal F$ can be regarded as a subset of $\mathrm{Aff}(\tr(\mathrm{C}([0, 1])))$.  Applying the Thomsen-Li Theorem above to $(\mathcal F, \eps)$ (where $X = [0, 1]$), one obtains $N$.

Let us first consider the case of AH algebras. Let $a \in A$ be a positive element with norm $1$; to prove the theorem, without loss of generality, one may assume that $a \in \mathrm{M}_n(\mathrm{C}(X)) \subseteq A$ for some $n \geq N$. Consider the homomorphism $\phi: \mathrm{C}[0, 1] \to \mathrm{M}_n(\mathrm{C}(X))$ induced by $a$, and consider the induced unital positive linear map $\phi^*: \mathrm{Aff}( \tr(\mathrm{C}[0, 1])) \to \mathrm{Aff}(\tr(\mathrm{C}(X)))$. By the Thomsen-Li Theorem, there are continuous maps $\lambda_1, ..., \lambda_n: X \to [0, 1]$ such that 
$$ \abs{\phi^*(f)(\tau) - \frac{1}{n}(\tau(f \circ \lambda_1) + \cdots + \tau(f \circ \lambda_n))} < \eps,\quad f\in \mathcal F,\  \tau \in \tr(\mathrm{C}(X)).$$
Then, with $$b = \mathrm{diag}\{(\lambda_1)_*(\mathrm{id}) , ..., (\lambda_n)_*(\mathrm{id})\} \in D_n,$$
(and since $\phi(f) = f(a)$,) one has 
$$ \abs{\tau(f(a)) - \tau(f(b))} < \eps,\quad  f\in \mathcal F,\ \tau \in \tr(\mathrm{M}_n(\mathrm{C}(X))).$$

Let us now consider the case that $A = \mathrm{C}(X) \rtimes \Gamma$, where $(X, \Gamma)$ is free and minimal, and has the (URP). 

Let $a \in A$ be a positive element with norm $1$. Choose $\delta>0$ such that if $\norm{x - y}_{2, \tr(A)} < \delta$, where $x, y$ are positive contractions, then $\norm{f(x) - f(y)}_{2, \tr(A)} < \eps$  for all $f \in \mathcal F$.

By Lemma \ref{pre-approx}, there exist a sub-C*-algebra $C \subseteq A$ and a positive contraction $p \in C$ such that
\begin{enumerate}
\item $C \cong \bigoplus_{i=1}^S \mathrm{M}_{n_i}(\mathrm{C}_0(Z_i))$, where $n_i>N$, $i=1, ..., S$,
\item $pap \in_{\delta/2} C$,
\item $\mathrm{d}_\tau(1-p) < \delta/2$, $\tau \in \tr(A)$,
\item\label{E-condition-4} there is a closed subset $[Z_{i}] \subseteq Z_i$ for each $i=1, ..., S$ such that if $a \in C$, $\norm{a} \leq 1$, and $a$ is supported in $\bigsqcup_{i=1}^S (Z_i \setminus [Z_{i}])$ under the isomorphism $C \cong \bigoplus_{i=1}^S\mathrm{M}_{n_i}(\mathrm{C}_0(Z_i))$, then $$ \tau(a) < \eps, \quad \tau \in \tr(A).$$
\end{enumerate}
Choose $\tilde{a} \in C$ such that $\norm{\tilde{a} - pap} < \delta/2$; hence
$$ \norm{a - \tilde{a}}_{2, \tr(A)} < \delta. $$ 
and so, by the choice of $\delta$, 
\begin{equation}\label{E-equation-1}
 \norm{f(a) - f(\tilde{a})}_{2, \tr(A)} < \eps,\quad f \in \mathcal F. 
 \end{equation}
Consider the homomorphism
$$ \phi: \mathrm{C}([0, 1]) \ni f \mapsto \pi_{[Z]}(f(\tilde{a})) \in \pi_{[Z]}(C) \cong  \bigoplus_{i=1}^S \mathrm{M}_{n_i}(\mathrm{C}([Z_i])),$$
where $[Z] = \bigsqcup_{i=1}^S [Z_i]$, and consider the unital positive linear map
$\phi^*:  \mathrm{Aff}( \tr(\mathrm{C}[0, 1])) \to \mathrm{Aff}(\tr(\mathrm{C}([Z]))).$ Then, by Theorem \ref{Thomsen-Li}, there are continuous maps $\lambda_{i, 1}, ..., \lambda_{i, n_i}: [Z_i] \to [0, 1]$, $i=1, ..., S$, such that
$$\abs{\phi^*(f)(\tau) - \frac{1}{n_i}(\tau(f \circ \lambda_{i, 1}) + \cdots + \tau(f \circ \lambda_{i, n_i}) )} < \eps,\quad f \in \mathcal F,\ \tau \in \tr(\mathrm{C}([Z_i])).$$
Define $$b_i = \mathrm{diag}\{ \mathrm{id} \circ \lambda_{i, 1}, ...,  \mathrm{id} \circ \lambda_{i, n} \} \in \mathrm{M}_{n_i}(\mathrm{C}([Z_i])),$$
where $\mathrm{id}$ is the identity function on $[0, 1]$,
and define 
$$b = b_1 \oplus \cdots \oplus b_{n_S}. $$ 
Then
\begin{equation}\label{E-equation-2}
\abs{\tau(\pi_{Z_0}(f(\tilde{a}))) - \tau(f(b))} < \eps,\quad f \in \mathcal F,\   \tau \in \tr(\bigoplus_{i=1}^S\mathrm{M}_{n_i}(\mathrm{C}([Z_i]))).
\end{equation}
Note that $b$ is a diagonal element. Extend $b$ to a positive diagonal contraction $\hat{b} \in C$. Note that, by Lemma \ref{pre-approx}(5), the element $\hat{b}$ is also in $D$. %such that $\norm{\hat{b}} \leq 1$. 
By \eqref{E-condition-4}, a calculation shows 
\begin{equation}\label{E-equation-3}
 \abs{\tau(f(\hat{b})) - \tau|_{[Z]}(f(b))} < 4 \eps,\quad f\in \mathcal F,\  \tau \in \tr(A),
 \end{equation}
where $\tau|_{[Z]}$ is the tracial state of $\bigoplus_{i=1}^S \mathrm{M}_{n_i}(\mathrm{C}([Z_i]))$ %, as the quotient of $C$, 
induced by $\tau|_C$:
$$\tau|_{[Z]}(x) = \lim_{n \to\infty}\frac{1}{\tau(e^2_n)} \tau(e_n \tilde{x} e_n), 
$$
where $x \in \bigoplus_{i=1}^S \mathrm{M}_{n_i}(\mathrm{C}([Z_i]))$, $\tilde{x} \in C \cong  \bigoplus_{i=1}^S \mathrm{M}_{n_i}(\mathrm{C}_0(Z_i))$ is an extension of $x$, and $(e_n)$ is a decreasing sequence of positive contractions of $C$ which converges (pointwise) to $\mathbf{1}_{Z}$. Also note that, by \eqref{E-condition-4},
\begin{equation}\label{E-equation-4}
\abs{\tau(f(\tilde{a})) - \tau|_{[Z]}(\pi_{[Z]}(f(\tilde{a})))} < 2\eps,\quad f \in \mathcal F,\  \tau \in \tr(A).
\end{equation}
Then, for all $f \in \mathcal F$ and $\tau \in \tr(A)$, by \eqref{E-equation-1}, \eqref{E-equation-4}, \eqref{E-equation-2}, and \eqref{E-equation-3}, 
$$ \tau(f(a)) \approx_\eps \tau(f(\tilde{a})) \approx_{2\eps} \tau|_{[Z]}(\pi_{Z_0}(f(\tilde{a}))) \approx_{\eps} \tau|_{[Z]}(f(b)) \approx_{4\eps} \tau(f(\hat{b})),$$
as desired. 
\end{proof}





\section{Some approximation lemmas}
In the section, let us prepare some approximation lemmas for the next section.

\begin{lem}\label{tr-to-norm-red}
Let $X$ be a metrizable compact space, and let $\Delta$ be a compact set of probability Borel measures of $X$. Then, for any $\eps \in (0, 1)$, there is $\delta>0$ such that if $a_0, a_1, d: X \to [0, 1]$ are continuous functions satisfying
\begin{enumerate}
\item $a_0a_1 = a_1$,
\item $\norm{a_0 d - d}_{2, \Delta} < \delta$, and
\item $\norm{da_1 - a_1}_{2, \Delta} < \delta$,
\end{enumerate}
then there is a continuous function $\tilde{d}: X \to [0, 1]$ such that
\begin{enumerate}
\item $\norm{d - \tilde{d}}_{2, \Delta} < \eps$,
\item $\norm{a_0 \tilde{d} - \tilde{d}} < \eps$, and
\item $\norm{\tilde{d} a_1 - a_1} < \eps$.
\end{enumerate}
\end{lem}

\begin{proof}
With the given $\eps$, choose $\eps' \in (0, 1)$ such that $\eps' < \eps$ and $2\sqrt{\eps'}< \eps$.
Then $$\delta= \eps'/\sqrt{2/(\eps')^3} 
$$
has the property of the lemma.

Indeed, let $a_0, a_1, d$ satisfy the conditions of the statement. Define sets 
$$A_{0, \leq 1-\eps'} = a_0^{-1}([0, 1-\eps']) \quad \mathrm{and}\quad A_{1, \geq \eps'} = a_1^{-1}([\eps', 1]).$$ Note that $A_{0, \leq 1-\eps'} \cap A_{1, \geq \eps'} = \O$.

On $A_{0, \leq 1-\eps'}$, consider the set
$$W_0 = \{x \in A_{0, \leq 1-\eps'}: \abs{d(x)}^2 \geq \eps'\}.$$ Then, for any $\mu \in \Delta$,
$$
(\eps')^2 \eps \mu(W_0) \leq (\eps')^2 \int_{A_{0, \leq 1-\eps'}} d^2(x) \mathrm{d}\mu(x) \leq \int_{A_{0, \leq 1-\eps'}} (a_0(x) - 1)^2d^2(x) \mathrm{d}\mu(x) < \delta^2,
$$
and hence
$$\mu(W_0) < \delta^2/(\eps')^3.$$

On $A_{1, \geq \eps'}$, consider the set
$$W_1 = \{x \in A_{1, \geq \eps'}: \abs{d(x) - 1}^2 \geq \eps'\}.$$ Then, for any $\mu \in \Delta$,
$$
(\eps')^2(\eps' \mu(W_1)) \leq (\eps')^2 \int_{A_{1, \geq \eps'}}\abs{d(x) - 1}^2 \mathrm{d}\mu(x) \leq \int_{A_{1, \geq \eps'}}\abs{d(x) - 1}^2 a_1^2(x) \mathrm{d}\mu(x) < \delta^2,
$$
and hence
$$\mu(W_1) < \delta^2/(\eps')^3.$$

Choose disjoint open sets $U_0 \supseteq W_0$ and $U_1 \supseteq$. Since $\Delta$ is compact, $U_0$ and $U_1$ can be chosen such that
$$ \mu(U_0) < \delta^2/(\eps')^3 \quad \mathrm{and} \quad \mu(U_1) < \delta^2/(\eps')^3,\quad \mu \in \Delta.$$ Choose continuous functions $r_0, r_1: X \to [0, 1]$ such that
$$r_i|_{W_i} = 1 \quad \mathrm{and} \quad r_i|_{X \setminus U_i} = 0,\quad i=0, 1.$$


Define 
$$
\tilde{d}(x) = \left\{
\begin{array}{ll}
0, & x \in W_0,\\
(1-r_0(x))d(x), & x \in U_0, \\
d(x), & x \in X \setminus (U_0 \cup U_1), \\
(1-r_1(x))d(x) + r_1(x), & x \in U_1, \\
1, & x \in W_1.
\end{array}
\right. 
$$

Then 
$$\int_X(d(x) - \tilde{d}(x))^2 \mathrm{d}\mu(x) \leq \mu(U_0) + \mu(U_1) < 2\delta^2/(\eps')^3,\quad \mu \in \Delta.$$ That is,
$$ \norm{d - \tilde{d}}_{2, \Delta} < \delta \sqrt{2/(\eps')^3} = \eps' < \eps.$$

Note that
$$\tilde{d}(x) < \sqrt{\eps'},\quad x \in A_{0, \leq 1-\eps'}.$$
So
$$\norm{a_0 \tilde{d} - \tilde{d}} < \max\{\eps', 2\sqrt{\eps'}\} < 2\sqrt{\eps'} < \eps.$$

Also note that
$$1-\tilde{d}(x) < \sqrt{\eps'}, \quad x \in A_{1, \geq \eps'}.$$
So
$$\norm{\tilde{d} a_1 - a_1} < \max\{\sqrt{\eps'}, 2\eps'\} < 2\sqrt{\eps'} < \eps,$$
as desired.
\end{proof}

The following lemma certainly is well known:
\begin{lem}\label{pert-almost-ortho}
Let $A$ be a C*-algebra. Let $N \in \mathbb N$ and $\eps>0$. Then there is $\delta>0$ such that if $c_1, ..., c_N$ are self-adjoint contractions such that $$ \norm{c_i c_j} < \delta,\quad i, j =1, ..., N,\ i\neq j,$$
then
$$\norm{c_1 + \cdots + c_N} < \max\{\norm{c_i}: i=1, ..., N\} + \eps.$$
\end{lem}

\begin{proof}
For the given $(N, \eps)$, there is $\delta>0$ such that if $c_1, ..., c_N$ are self-adjoint contractions such that $$ \norm{c_i c_j} < \delta,\quad i, j =1, ..., N,\ i\neq j,$$
then there are self-adjoint elements $\tilde{c}_1, ..., \tilde{c}_N \in A$ such that
$$ \norm{\tilde{c}_i - c_i} < \eps/2N, \quad \textrm{and} \quad \tilde{c}_i \perp \tilde{c}_j,\quad i, j = 1, ..., N,\  i\neq j.$$
Then, this $\delta$ has the property of the lemma, as
\begin{eqnarray*}
\norm{c_1 + \cdots + c_N} & \approx_{\eps/2} & \norm{\tilde{c}_1 + \cdots + \tilde{c}_N} \\
& = & \max\{\norm{\tilde{c}_i}: i=1, ..., N\} \\
& \approx_{\eps/2} & \max\{\norm{c_i}: i=1, ..., N\},
\end{eqnarray*}
as desired.
\end{proof}

\begin{lem}[cf.~\cite{CE-str1}]\label{hereditary-sets-norm-app}
Let $A$ be a C*-algebra. For any $\eps>0$, there are $N \in \mathbb N$ and $\delta>0$ with the following property:

Let $a \in A$ be a positive element with norm at most $1$. Define $$a_i = \chi_i(a),\quad  i=1, ..., N, $$ where $$\chi_i(t) = \left\{ 
\begin{array}{ll}
0, & t \leq \frac{i-1}{N}, \\
\mathrm{linear}, & t \in [\frac{i-1}{N}, \frac{i}{N}], \\
1, & t \geq \frac{i}{N}.
\end{array} 
\right.
$$
Assume there are positive elements $d_1, ..., d_N\in A$  with norm at most $1$ such that
         \begin{enumerate}
         \item\label{commut-cond} $[a, d_i] = 0$, $i=1, ..., N$,
         \item\label{here-cond-1-a} $ \norm{a_{i} d_{i+1} - d_{i+1}} < \delta$, $i=1, ..., N-1$,
         \item\label{here-cond-1-b} $ \norm{d_i a_{i+1} - a_{i+1}} < \delta$, $i=1, ..., N-1$.
         \end{enumerate} 
        Then
        $$\norm{ a  - \frac{1}{N}(d_1 + \cdots + d_N)} <  \eps.$$
\end{lem}

\begin{proof}
Choose $N > {32}/{\eps}.$ Applying Lemma \ref{pert-almost-ortho} to $N$ and $\eps/16$, one obtains $\delta_1$.
Choose $\delta>0$ such that
$$16\delta< \delta_1 \quad \mathrm{and} \quad \frac{\eps}{8} + \delta + 4((4\delta+\frac{\eps}{8}) + \frac{\eps}{16}) < \eps. $$
 Then the pair $(N, \delta)$ has the property of the lemma.

Since $a_i a_j = a_j$, $i < j$, by \eqref{here-cond-1-a},
$$a_i d_j \approx_\delta a_i a_{j-1} d_j = a_{j-1} d_j \approx_\delta d_j,\quad i< j. $$ Hence, together with \eqref{here-cond-1-a},
\begin{equation}\label{here-cond-1-a-eq}
\norm{a_id_j - d_j} < 2\delta,\quad  i < j.
\end{equation}
A similar argument applied to \eqref{here-cond-1-b} shows
\begin{equation}\label{here-cond-1-b-eq}
\norm{d_ia_j - a_j} < 2\delta, \quad  i < j.
\end{equation}

Also note that, if $i < j-1$, then
$$d_i d_j \approx_\delta d_i a_{j-1} d_j \approx_{2\delta} a_{j-1} d_j \approx_\delta d_j,$$
and therefore, a straightforward calculation shows that 
\begin{equation}\label{orth-cond}
\norm{(d_i-d_{i+1})(d_j - d_{j+1})} <16\delta < \delta_1 ,\quad i < j-2.
\end{equation}

Note that $$a \approx_{\frac{2}{N}} aa_2,$$
and then, together with \eqref{here-cond-1-b}, one has
\begin{equation}\label{first-bump}
a \approx _{\frac{2}{N}} aa_2 \approx_{\delta } a a_2d_1 \approx_{\frac{2}{N}} ad_1.
\end{equation}
Noting that $$a = \frac{1}{N}(a_1 + \cdots + a_N),$$ together with \eqref{here-cond-1-a-eq} and  \eqref{here-cond-1-b-eq}, 
one also has that, for $i=1, ..., N-1$, 
\begin{eqnarray*}
& & a(d_i - d_{i+1})  \\
& = & \frac{1}{N}(a_1 + \cdots + a_N)(d_i - d_{i+1}) \nonumber \\
& = & \frac{1}{N}((a_1 d_i + \cdots + a_Nd_i) - (a_1 d_{i+1} + \cdots + a_Nd_{i+1})) \nonumber \\
& \approx_{4\delta} & \frac{1}{N}((\underbrace{d_i + \cdots + d_{i}}_{i-1} + a_id_i + a_{i+1} + \cdots + a_N) \\
&&  - (\underbrace{d_{i+1} + \cdots + d_{i+1}}_{i} +  a_{i+1}d_{i+1} + a_{i+2}+ \cdots + a_N)) \nonumber \\
& = & \frac{i}{N}(d_i - d_{i+1}) + \frac{1}{N}(a_{i+1}-d_{i} + a_id_i - a_{i+1}d_{i+1}) \nonumber \\
&\approx_{\frac{4}{N}}&  \frac{i}{N}(d_i - d_{i+1}). \nonumber
\end{eqnarray*}
That is, 
\begin{equation}\label{local-bump}
\norm{a(d_i - d_{i+1}) -  \frac{i}{N}(d_i - d_{i+1})} < 4\delta + \frac{4}{N} < 4\delta + \frac{\eps}{8},\quad i=1, 2, ..., N-1.
\end{equation}
A similar argument also shows
\begin{equation}\label{local-bump-1}
\norm{ad_N - d_N} < 2\delta + \frac{2}{N} < 4\delta + \frac{\eps}{8}.
\end{equation}

Then, on applying \eqref{first-bump}, \eqref{local-bump}, \eqref{local-bump-1}, \eqref{orth-cond}, and Lemma \ref{pert-almost-ortho} (note that $a, d_1, ..., d_N$ commute),
\begin{eqnarray*}
a & \approx_{\frac{4}{N} + \delta} & ad_1\\
& = & a(d_1 - d_2 + d_2 - d_3 + \cdots + d_{N-1} - d_{N} + d_{N}) \\
& = & a(d_1 - d_2) + a(d_2 - d_3) + \cdots + a(d_{N-1} - d_{N}) + ad_{N} \\
& \approx_{4((4\delta + \frac{\eps}{8}) + \frac{\eps}{16})} & \frac{1}{N}(d_1 - d_2) + \frac{2}{N}(d_2 - d_3)  + \cdots + \frac{N-1}{N}(d_{N-1} - d_{N}) + d_{N} \\
& = & \frac{1}{N}(d_1 + d_2 + \cdots + d_{N}),
\end{eqnarray*}
as desired.
\end{proof}

\begin{lem}\label{hereditary-sets-tr-app}
Let $(D, A)$ be a pair of unital C*-algebras, where $D$ is commutative. For any $\eps>0$, there are $\delta>0$ and $N \in \mathbb N$ with the following property:  

Let $a \in D \subseteq A$ be a positive element with norm at most $1$, and set $$a_i = \chi_i(a),\quad  i=1, ..., N, $$ where $$\chi_i(t) = \left\{ 
\begin{array}{ll}
0, & t \leq \frac{i-1}{N}, \\
\mathrm{linear}, & t \in [\frac{i-1}{N}, \frac{i}{N}], \\
1, & t \geq \frac{i}{N}.
\end{array} 
\right.
$$
Let $d_1, ..., d_N \in A$ be positive elements with norm at most $1$ such that
         \begin{enumerate}
         \item $\mathrm{dist}_{2, \tr(A)}(d_i, (D)^+_1) < \delta$, $i=1, 2, ..., N$,
         \item\label{here-cond-1-a-tr} $ \norm{a_{i} d_{i+1} - d_{i+1}}_{2, \tr(A)} < \delta$, $i=1, ..., N-1$, and 
         \item\label{here-cond-1-b-tr} $ \norm{d_i a_{i+1} - a_{i+1}}_{2, \tr(A)} < \delta$, $i=1, ..., N-1$.
         \end{enumerate} 
        Then
        $$\norm{ a  - \frac{1}{N}(d_1 + \cdots + d_N)}_{2, \tr(A)} < \eps.$$
\end{lem}

\begin{proof}
Applying Lemma \ref{hereditary-sets-norm-app} to $\eps/2$, one obtains, say, the pair $(N, \delta_2)$. Applying Lemma \ref{tr-to-norm-red} to $\min\{\delta_2, \eps/4\}$, one obtains $\delta_1$. Set $\delta = \min\{\delta_1/4, \eps/4\}$. Then $(N, \delta)$ has the property of the lemma.

Indeed, let $d_1, ..., d_N$ be given. Since $\mathrm{dist}_{2, \tr(A)}(d_i, (D)^+_1) < \delta$, $i=1, 2, ..., N$, there are positive contractions $\tilde{d}_1, ... \tilde{d}_N \in D$ such that
\begin{equation}\label{pert-lem-cond-1}
\norm{d_i - \tilde{d}_i}_{2, \tr(A)} < \min\{\delta_1/4, \eps/4\},\quad i=1, ..., N.
\end{equation}
Then, it follows from \eqref{here-cond-1-a-tr} and \eqref{here-cond-1-b-tr} that
$$  \norm{a_{i} \tilde{d}_{i+1} - \tilde{d}_{i+1}}_{2, \tr(A)} < \delta_1 \quad \mathrm{and} \quad \norm{\tilde{d}_i a_{i+1} - a_{i+1}}_{2, \tr(A)} < \delta_1,\quad i=1, 2, ..., N-1. $$ 
Applying Lemma \ref{tr-to-norm-red} to each element $\tilde{d}_i$, $i=1, ..., N$, one obtains a positive contraction $\tilde{\tilde{d}}_i \in D$ such that
\begin{equation}\label{pert-lem-cond-2}  
 \norm{\tilde{d}_i - \tilde{\tilde{d}}_i}_{2, \tr(A)} < \eps/4,\quad i=1, ..., N 
 \end{equation}
and
$$  \norm{a_{i} \tilde{\tilde{d}}_{i+1} - \tilde{\tilde{d}}_{i+1}} < \delta_2 \quad \mathrm{and} \quad \norm{\tilde{\tilde{d}}_i a_{i+1} - a_{i+1}} < \delta_2,\quad i=1, 2, ..., N-1. $$ 
Then, by Lemma \ref{hereditary-sets-norm-app}, one has
$$\norm{ a  - \frac{1}{N}(\tilde{\tilde{d}}_1 + \cdots + \tilde{\tilde{d}}_N)} <  \frac{\eps}{2},$$
and hence, together with \eqref{pert-lem-cond-1} and \eqref{pert-lem-cond-2}, one has
$$\norm{ a  - \frac{1}{N}(d_1 + \cdots + d_N)}_{2, \tr(A)} < \frac{\eps}{2} + \frac{\eps}{2} = \eps,$$
as desired.
\end{proof}

\begin{lem}\label{hereditary-sets-preserve-tr}
Let $A$ be a unital C*-algebra. For any $\eps>0$, any $N \in \mathbb N$, and any $\chi\in \mathrm{C}([0, 1])^+$, there are $\delta>0$ and $M \in \mathbb N$ such that if $b_1, b_2, ..., b_N \in A$ and $d_1, d_2, ..., d_N \in A$ are positive elements with norm at most $1$ such that
\begin{enumerate}
\item $ \norm{d_id_{i+1} - d_{i+1}}_{2, \tr(A)} < \delta$, $i=1, ..., N-1$,  
\item $b_i b_{i+1} = b_{i+1}$, $i=1, 2, ..., N$, and
\item $\abs{\tau(b_i^j) - \tau(d_i^j)} < \delta$, $i=1, 2, ..., N$, $j=1, ..., M$, $\tau \in \tr(A)$,
\end{enumerate}
then
$$\abs{\tau(\chi(\frac{1}{N}(b_1 + \cdots + b_N))) - \tau(\chi(\frac{1}{N}(d_1 + \cdots + d_N))) } < \eps,\quad \tau \in \tr(A).$$
\end{lem}

\begin{proof}
It is enough to prove the statement for $\chi$ a monomial, i.e., $\chi(t) = t^n$. Note that there are positive numbers $\alpha_{i, j}$, $i=1, ..., N$, $j=1, ..., n$, such that
\begin{eqnarray*}
b^n & = & (\frac{1}{N}(b_1 + \cdots + b_N))^n \\
& = & \frac{1}{N^n}\sum_{i_1+\cdots + i_N = n} b_1^{i_1}\cdots b_N^{i_N} \\
& = & \sum_{i=1}^N \sum_{j=1}^n \alpha_{i, j} b_i^j.
\end{eqnarray*}
Hence,
$$\tau(b^n) = \sum_{i=1}^N \sum_{j=1}^n \alpha_{i, j} \tau(b_i^j).$$

Then there is $\delta>0$ such that if $$ \norm{d_id_{i+1} - d_{i+1}}_{2, \tr(A)} < \delta,\quad i=1, ..., N-1,$$ then
$$ \norm{ (\frac{1}{N}(d_1 + \cdots + d_N))^n - \sum_{i=1}^N \sum_{j=1}^n \alpha_{i, j} d_i^j }_{2, \tr(A)} < \eps/2.$$ 
In particular,
$$\tau((\frac{1}{N}(d_1+\cdots + d_N))^n) \approx_{\eps/2} \tau( \sum_{i=1}^N \sum_{j=1}^n \alpha_{i, j} d_i^j) = \sum_{i=1}^N \sum_{j=1}^n \alpha_{i, j} \tau(d_i^j). $$

Moreover, one may assume that $\delta>0$ is sufficiently small such that if 
$$\abs{\tau(b_i^j) - \tau(d_i^j)} < \delta,\quad i=1, 2, ..., N,\ j=1, ..., n,$$
then
$$ \abs{ \sum_{i=1}^N \sum_{j=1}^n \alpha_{i, j} \tau(b_i^j) - \sum_{i=1}^N \sum_{j=1}^n \alpha_{i, j} \tau(d_i^j) } < \eps/2.$$
Then this $\delta$ and $M:=n$ have the desired property.
\end{proof}

%
%%%%%%%%%%%%% The original approximation lemma. %%%%%%%%%%%%%%%%%%
%\begin{lem}[cf.~\cite{CE-str1}]\label{hereditary-sets}
%Let $A$ be a unital C*-algebra. 
%\begin{enumerate}
%
%\item[(1)] Let $\eps>0$, and let $a \in A$ be a positive element with norm at most $1$. Define $$a_i = \chi_i(a),\quad  i=1, ..., N, $$ where $$\chi_i(t) = \left\{ 
%\begin{array}{ll}
%0, & t \leq \frac{i-1}{N}, \\
%\mathrm{linear}, & t \in [\frac{i-1}{N}, \frac{i}{N}], \\
%1, & t \geq \frac{i}{N}.
%\end{array} 
%\right.
%$$
%Assume there are positive elements $d_1, d_2, ..., d_{N-1}, (d_N= 0) \in A$  with norm at most $1$ such that
%         \begin{enumerate}
%         \item\label{here-cond-1-a} $ \norm{a_{i} d_{i+1} - d_{i+1}} < \eps$, $i=1, ..., N-1$,
%         \item\label{here-cond-1-b} $ \norm{d_i a_{i+1} - a_{i+1}} < \eps$, $i=1, ..., N-1$.
%         \end{enumerate} 
%        Then
%        $$\norm{ a  - \frac{1}{N}(d_1 + \cdots + d_N)} < (4N^2+8N+1)\eps + \frac{8}{N}.$$
%
%\item[(1')] For any $\eps>0$, there are $\delta>0$ (the $\eps''$ in the main theorem) and $N \in \mathbb N$ with the following property:  
%
%Let $a \in D \subseteq A$ be a positive element with norm at most $1$, and denote by $$a_i = \chi_i(a),\quad  i=1, ..., N, $$ where $$\chi_i(t) = \left\{ 
%\begin{array}{ll}
%0, & t \leq \frac{i-1}{N}, \\
%\mathrm{linear}, & t \in [\frac{i-1}{N}, \frac{i}{N}], \\
%1, & t \geq \frac{i}{N}.
%\end{array} 
%\right.
%$$
%Let $d_1, d_2, ..., d_{N-1}, (d_N= 0) \in A$ be positive elements with norm at most $1$ such that
%         \begin{enumerate}
%         \item $\mathrm{dist}_{2, \tr(A)}(d_i, D) < \delta$, $i=1, 2, ..., N$,
%         \item\label{here-cond-1-a-tr} $ \norm{a_{i} d_{i+1} - d_{i+1}}_{2, \tr(A)} < \delta$, $i=1, ..., N-1$, and 
%         \item\label{here-cond-1-b-tr} $ \norm{d_i a_{i+1} - a_{i+1}}_{2, \tr(A)} < \delta$, $i=1, ..., N-1$.
%         \end{enumerate} 
%        Then
%        $$\norm{ a  - \frac{1}{N}(d_1 + \cdots + d_N)}_{2, \tr(A)} < \eps.$$
%
%        
%\item[(2)] For any $\eps>0$ and any $\chi\in \mathrm{C}([0, 1])^+$, there are $\delta>0$ and $M \in \mathbb N$ such that if $b_1, b_2, ..., b_N \in A$ and $d_1, d_2, ..., d_N \in A$ are positive elements with norm at most $1$ such that
%\begin{enumerate}
%\item $ \norm{d_id_{i+1} - d_{i+1}}_{2, \tr(A)} < \delta$, $i=1, ..., N-1$,  %and 
%%\item there positive elements $b_1, ..., b_N \in A$ with norm at most $1$ such that 
%%\begin{enumerate}
%\item $b_i b_{i+1} = b_{i+1}$, $i=1, 2, ..., N$, and
%\item $\abs{\tau(b_i^j) - \tau(d_i^j)} < \delta$, $i=1, 2, ..., N$, $j=1, ..., M$, $\tau \in \tr(A)$.
%%\end{enumerate}
%\end{enumerate}
%Then
%$$\abs{\tau(\chi(\frac{1}{N}(b_1 + \cdots + b_N))) - \tau(\chi(\frac{1}{N}(d_1 + \cdots + d_N))) } < \eps,\quad \tau \in \tr(A).$$
%%where $$b:=\frac{1}{N}(b_1 + \cdots + b_N). $$
%
%
%\end{enumerate}
%
%\end{lem}
%
%\begin{proof}
%Let us prove (1). 
%Since $a_i a_j = a_j$, $i < j$, by \eqref{here-cond-1-a},
%$$a_i d_j \approx_\eps a_i a_{j-1} d_j = a_{j-1} d_j \approx_\eps d_j,\quad i< j. $$ Hence, together with \eqref{here-cond-1-a},
%\begin{equation}\label{here-cond-1-a-eq}
%\norm{a_id_j - d_j} < 2\eps,\quad  i < j.
%\end{equation}
%A similar argument applying to \eqref{here-cond-1-b} shows
%\begin{equation}\label{here-cond-1-b-eq}
%\norm{d_ia_j - a_j} < 2\eps,\quad  i < j.
%\end{equation}
%
%Also note that, if $i < j-1$, then
%$$d_i d_j \approx_\eps d_i a_{j-1} d_j \approx_{2\eps} a_{j-1} d_j \approx_\eps d_j,$$
%and therefore, a straightforward calculation shows that 
%\begin{equation}\label{orth-cond}
%\norm{(d_i-d_{i+1})(d_j - d_{j+1})} <16\eps ,\quad i < j-2.
%\end{equation}
%
%Note that $$a \approx_{\frac{2}{N}} aa_2,$$
%and then, together with \eqref{here-cond-1-b}, one has
%\begin{equation}\label{first-bump}
%a \approx _{\frac{2}{N}} aa_2 \approx_{\eps } a a_2d_1 \approx_{\frac{2}{N}} ad_1.
%\end{equation}
%Noting that $$a = \frac{1}{N}(a_1 + \cdots + a_N),$$ together with \eqref{here-cond-1-a-eq} and  \eqref{here-cond-1-b-eq}, 
%one also has that, for $i=1, ..., N-1$, 
%\begin{eqnarray*}
%& & a(d_i - d_{i+1})  \\
%& = & \frac{1}{N}(a_1 + \cdots a_N)(d_i - d_{i+1}) \nonumber \\
%& = & \frac{1}{N}((a_1 d_i + \cdots + a_Nd_i) - (a_1 d_{i+1} + \cdots + a_Nd_{i+1})) \nonumber \\
%& \approx_{4\eps} & \frac{1}{N}((\underbrace{d_i + \cdots + d_{i}}_{i-1} + a_id_i + a_{i+1} + \cdots + a_N) \\
%&&  - (\underbrace{d_{i+1} + \cdots + d_{i+1}}_{i} +  a_{i+1}d_{i+1} + a_{i+2}+ \cdots + a_N)) \nonumber \\
%& = & \frac{i}{N}(d_i - d_{i+1}) + \frac{1}{N}(a_{i+1}-d_{i} + a_id_i - a_{i+1}d_{i+1}) \nonumber \\
%&\approx_{\frac{4}{N}}&  \frac{i}{N}(d_i - d_{i+1}). \nonumber
%\end{eqnarray*}
%That is, 
%\begin{equation}\label{local-bump}
%\norm{a(d_i - d_{i+1}) -  \frac{i}{N}(d_i - d_{i+1})} < 4\eps + \frac{2}{N},\quad i=1, 2, ..., N-1.
%\end{equation}
%
%Then, applying \eqref{first-bump}, \eqref{local-bump}, and \eqref{orth-cond},
%\begin{eqnarray*}
%a & \approx_{\frac{4}{N} + \eps} & ad_1\\
%& = & a(d_1 - d_2 + d_2 - d_3 + \cdots + d_{N-1} - d_{N} + d_{N}) \\
%& = & a(d_1 - d_2) + a(d_2 - d_3) + \cdots + a(d_{N-1} - d_{N}) + ad_{N} \\
%& \approx_{2(4N\eps + \frac{2}{N}) + (\frac{N}{2})^2 16\eps} & \frac{1}{N}(d_1 - d_2) + \frac{2}{N}(d_2 - d_3)  + \cdots + \frac{N-1}{N}(d_{N-1} - d_{N}) + d_{N} \\
%& = & \frac{1}{N}(d_1 + d_2 + \cdots + d_{N-1}).
%\end{eqnarray*}
%This proves (1).
%
%As for (2), it is enough to prove the statement of $\chi$ to be a monomial, i.e, $\chi(t) = t^n$. Note that there are positive numbers $\alpha_{i, j}$, $i=1, ..., N$, $j=1, ..., n$, such that
%\begin{eqnarray*}
%b^n & = & (\frac{1}{N}(b_1 + \cdots + b_N))^n \\
%& = & \frac{1}{N^n}\sum_{i_1+\cdots + i_N = n} b_1^{i_1}\cdots b_N^{i_N} \\
%& = & \sum_{i=1}^N \sum_{j=1}^n \alpha_{i, j} b_i^j.
%\end{eqnarray*}
%Hence
%$$\tau(b^n) = \sum_{i=1}^N \sum_{j=1}^n \alpha_{i, j} \tau(b_i^j).$$
%
%Then there is $\delta>0$ such that if $$ \norm{d_id_{i+1} - d_{i+1}}_2 < \delta,\quad i=1, ..., N-1,$$ then
%$$ \norm{ (\frac{1}{N}(d_1 + \cdots + d_N))^n - \sum_{i=1}^N \sum_{j=1}^n \alpha_{i, j} d_i^j }_2 < \eps/2.$$ 
%In particular,
%$$\tau((\frac{1}{N}(d_1+\cdots + d_N))^n) \approx_{\eps/2} \tau( \sum_{i=1}^N \sum_{j=1}^n \alpha_{i, j} d_i^j) = \sum_{i=1}^N \sum_{j=1}^n \alpha_{i, j} \tau(d_i^j). $$
%
%Moreover, one may assume that $\delta>0$ is sufficiently small such that if 
%$$\abs{\tau(b_i^j) - \tau(d_i^j)} < \delta,\quad i=1, 2, ..., N,\ j=1, ..., n,$$
%then
%$$ \abs{ \sum_{i=1}^N \sum_{j=1}^n \alpha_{i, j} \tau(b_i^j) - \sum_{i=1}^N \sum_{j=1}^n \alpha_{i, j} \tau(d_i^j) } < \eps/2.$$
%Then this $\delta$ together with $M:=n$ have the desired property.
%\end{proof}
%
%%%%%%%%%%%%%%%%%%%%%%%%%%%%%%%%%%%%%%%%%%%%




\section{The small boundary property}

%Let us first introduce the following relative comparison property of a commutative C*-algebra inside an ambient C*-algebra:
%\begin{defn}\label{defn-C}
%Let $A$ be a unital C*-algebra and let $D$ be a unital commutative sub-C*-algebra. Then the pair $(D, A)$ is said to have Property (C) if for any positive contractions $f, g, h \in D$ satisfying $f, g \in \overline{hDh}$, and $$\mathrm{d}_\tau(f) < \mathrm{d}_\tau(g),\quad \tau \in \tr(A),$$ and for any $\eps>0$, there is a contraction $u \in \overline{hAh} + \Comp 1_A$ such that
%$$ u f u^* \in_\eps^{\norm{\cdot}_2} \overline{gAg},$$
%$$\mathrm{dist}_{2, \tr(A)}(udu^*, (D)_1) < \eps,\quad \mathrm{dist}_{2, \tr(A)}(u^*du, (D)_1) < \eps, \quad d \in (D)_1, $$ 
%and
%$$\norm{uu^* - 1}_{2, \tr(A)},\  \norm{u^*u - 1}_{2, \tr(A)} < \eps.$$
%\end{defn}
%
%\begin{rem}
%Comparing to Property (COS), the approximations in Property (C) are with respect to the uniform trace norm, but on the other hand, the comparison in Property (C) is implemented by an almost unitary which is also an almost normalizer (with respect to the uniform trace norm).
%\end{rem}
%
%\begin{rem}
%If the free and minimal dynamical system $(X, \Gamma)$ has the (SBP), then an argument with almost finiteness of \cite{KS-comparison} implies that $\mathrm{C}(X) \rtimes \Gamma$ has Property (C).
%\end{rem}


Let us show that Property (S) (Definition \ref{defn-S}) for the ambient C*-algebra $A$ indeed implies the (SBP) for the subalgebra $D$ when Properties (C) and (E) (Definitions \ref{defn-C} and \ref{defn-D} ) are present (Theorem \ref{prop-S}). Since the crossed product C*-algebra $\mathrm{C}(X) \rtimes \Gamma$ always has Property (E) and Property (S) (Theorem \ref{S-holds}) if $(X, \Gamma)$ has the (URP), together with Theorem \ref{SBP=>C}, this implies that Property (C) is equivalent to the (SBP) in this setting.

For each $\eps>0$, define 
\begin{equation}\label{defn-eta}
\eta_\eps(t) = \left\{
\begin{array}{ll}
0, & t \leq 1-\eps, \\
\mathrm{linear}, & t \in [1-\eps, 1-\eps/2], \\
1, & t \geq 1-\eps/2.
\end{array}
\right.
\end{equation}

In the proof of the following lemma, we use $O(\eps)$ to denote a quantity which converges to $0$ when $\eps$ approaches $0$.
\begin{lem}\label{D-sandwich}
Let $A$ be a unital C*-algebra, and let $D \subseteq A$ be a unital commutative subalgebra such that the pair $(D, A)$ has Property (C). Let $\phi_1, \phi_2, \phi_3, \psi_1, \psi_2 \in (D)_1^+$ and $\eps_0>0$ have the following properties:
\begin{enumerate}
\item\label{cond-s-1} $\phi_2, \phi_3 \in \overline{\phi_1 D\phi_1}$,
\item\label{cond-s-2} $\psi_1\psi_2 = \psi_2$, $\phi_2 \phi_3 = \phi_3$,
\item\label{cond-s-3} $\mathrm{d}_\tau(\psi_1) < \mathrm{d}_\tau(\phi_1)$ for all $\tau\in\tr(A)$,
\item\label{cond-s-4} $\inf\{\tau(\psi_2) - \mathrm{d}_\tau(\phi_3): \tau \in \tr(A)\} >0 $, and
\item\label{cond-s-5} $\inf\{\mathrm{d}_\tau(\phi_2) - \tau(\eta_{\eps_0}(\psi_1)): \tau \in \tr(A)\} >0 $.  \end{enumerate}
Then, for any $\eps>0$, there is a contraction $u\in A$ such that
$$ u^*\psi_1u \in_\eps^{\norm{\cdot}_2}  \overline{\phi_1A\phi_1},\quad u^*\eta_{\eps_0/2}(\psi_1)u \in^{\norm{\cdot}_2}_{\eps} \overline{\phi_2A\phi_2}, \quad  \eta_{\eps}(u^*\psi_1u) \phi_3 \approx_{\eps}^{\norm{\cdot}_2} \phi_3,$$
and
$$\mathrm{dist}_{2, \tr(A)}(udu^*, (D)_1) < \eps,\quad \mathrm{dist}_{2, \tr(A)}(u^*du, (D)_1) < \eps,\quad d \in (D)_1,$$
and
$$\norm{uu^* - 1}_{2, \tr(A)}, \quad \norm{u^*u - 1}_{2, \tr(A)} < \eps.$$
\end{lem}

\begin{proof}
Let $\eps>0$ be given. Choose $\delta \in (0, \min\{\eps_0, \eps\}/4)$ such that if $x, y$ are positive contractions of a unital C*-algebra $A$ such that $\norm{x - y}_{2, \tr(A)} < \delta$, then
\begin{equation}\label{pre-pert-delta}
\norm{\eta_\eps(x) - \eta_\eps(y)}_{2, \tr(A)} < \eps/2.
\end{equation}


Define
$$ \delta_1 := \inf\{\tau(\psi_2) - \mathrm{d}_\tau(\phi_3): \tau \in \tr(A)\} >0$$
and
$$\delta_2 := \inf\{\mathrm{d}_\tau(\phi_2) - \tau(\eta_{\eps_0}(\psi_1)): \tau \in \tr(A)\} > 0. $$

By Property (C) and Assumption (\ref{cond-s-3}), for any $\eps'>0$ (to be determined later), there is a contraction $u_1\in A$ such that
\begin{equation}\label{approx-C-1-h}
u^*_1 \psi_1 u_1 \in^{\norm{\cdot}_2}_{\eps'} \overline{\phi_1 A \phi_1}, 
\end{equation}
\begin{equation}\label{approx-C-1-n}
\mathrm{dist}_{2, \tr(A)}(u_1du^*_1, (D)_1) < \eps', \quad \mathrm{dist}_{2, \tr(A)}(u_1^*du_1, (D)_1) < \eps',\quad d \in (D)_1, 
\end{equation}
and
\begin{equation}\label{approx-C-1-u}
\norm{u_1u_1^* - 1}_{2, \tr(A)},\quad  \norm{u_1^*u_1 - 1}_{2, \tr(A)} < \eps'.
\end{equation}

By \eqref{approx-C-1-h}, there is $n$ large enough that $$ \norm{(\phi_1^{\frac{1}{n}}) (u_1^*\psi_1u_1) - u_1^*\psi_1u_1}_{2, \tr(A)} < \eps'. 
$$
By \eqref{approx-C-1-n}, there is a positive contraction $d_1 \in D$ such that
\begin{equation*}
\norm{d_1- u_1^*\psi_1u_1}_{2, \tr(A)}< \eps', 
\end{equation*}
and then
$$\norm{\phi_1^{\frac{1}{n}}d_1 - u_1^*\psi_1u_1}_{2, \tr(A)} < 2\eps'. 
$$
Consider $\phi_1^{\frac{1}{n}}d_1$, and still denote this element by $d_1$. One has $$d_1 \in \overline{\phi_1 D \phi_1}$$
and
\begin{equation}\label{approx-C-1}
\norm{d_1 - u_1^*\psi_1u_1}_{2, \tr(A)} < 2\eps'. 
\end{equation}


Consider  the contraction
$$\eta_{\eps_0}(u^*_1\psi_1 u_1),$$
and note that (by \eqref{approx-C-1}, \eqref{approx-C-1-u} and Condition \eqref{cond-s-5})
$$ \mathrm{d}_\tau( \eta_{\eps_0/2}(d_1) ) \leq \tau(\eta_{\eps_0}(d_1))  \approx_{O(\eps')}^{\norm{\cdot}_2}  \tau( \eta_{\eps_0}(u^*_1\psi_1 u_1) ) \approx_{O(\eps')}^{\norm{\cdot}_2} \tau( \eta_{\eps_0}(\psi_1) ) < \mathrm{d}_\tau(\phi_2) - \delta_2/2,\quad \tau \in \tr(A).$$
With $\eps'$ sufficiently small, one has 
$$\mathrm{d}_\tau(\eta_{\eps_0/2}(d_1)) < \mathrm{d}_\tau(\phi_2),\quad \tau \in \tr(A).$$ 
By Property (C), for any $\eps''>0$ (to be determined later), there is $u_2 \in \overline{\phi_1 A \phi_1}+\Comp 1$ such that 
\begin{equation}\label{approx-C-2-here}
u^*_2 \eta_{\eps_0/2}(d_1) u_2 \in^{\norm{\cdot}_2}_{\eps''} \overline{\phi_2 A \phi_2},
\end{equation}
\begin{equation}\label{approx-C-2-n}
\mathrm{dist}_{2, \tr(A)}(u_2 d u^*_2, D_1) < \eps'', \quad \mathrm{dist}_{2, \tr(A)}(u_2^*du_2, D_1) < \eps'',\quad d \in D_1,
\end{equation}
and
\begin{equation}\label{approx-C-2-u}
\norm{u_2u_2^* - 1}_{2, \tr(A)},\quad  \norm{u_2^*u_2 - 1}_{2, \tr(A)} < \eps''.
\end{equation}


Since (as $\delta < \eps_0/4$)
$$\eta_{\eps_0/2}(d_1) \eta_\delta(d_1) = \eta_\delta(d_1),$$
it follows from \eqref{approx-C-2-here} that
\begin{equation}\label{close-hereditary}
u_2^* \eta_\delta(d_1) u_2 \in^{\norm{\cdot}_2}_{2\eps''} \overline{\phi_2 A \phi_2}. 
\end{equation}

By \eqref{approx-C-2-n}, there are positive contractions $$ d_{1, 1},\  d_{1, \delta} \in  D $$ such that
\begin{equation}\label{pert-C-step-2}
 \norm{d_{1, 1} - u_2^* d_1 u_2}_{2, \tr(A)} < \eps'',\quad  \norm{d_{1, \delta} - u_2^*\eta_{\delta}(d_1)u_2}_{2, \tr(A)} < \eps''. 
\end{equation}  
 By \eqref{close-hereditary}, $$d_{1, \delta} \in_{3\eps''}^{\norm{\cdot}_2} \in \overline{\phi_2A\phi_2}.$$
 With the same argument as above for $d_1$, there is a positive contraction, still denoted by $ d_{1, \delta}$, such that 
 \begin{equation}\label{u-2-2-approx-C-1}
 d_{1, \delta} \in \overline{\phi_2 D \phi_2}\quad \mathrm{and}\quad \norm{d_{1, \delta} - u_2^*\eta_{\delta}(d_1)u_2}_{2, \tr(A)} < 5\eps''.
 \end{equation}
Also note that, by \eqref{approx-C-2-u}, $$ \norm{ \eta_\delta(d_{1, 1}) - d_{1, \delta} }_{2, \tr(A)} = O(\eps''). $$

Define
$$\tilde{d}_{1, 1} = f_{\delta}(d_{1, 1}) \quad \mathrm{and} \quad \tilde{d}_{1, \delta} = \eta_{\delta}(d_{1, 1}),$$
where
$$
f_{\delta}(t) = \left\{ 
\begin{array}{ll}
0, & t \leq 0, \\
\mathrm{linear}, & t\in [0, 1-\delta], \\
1, & t \geq 1-\delta.
\end{array}
\right.
$$ 
Then
\begin{equation}\label{pert-fur-hereditary}
 \tilde{d}_{1, 1} \tilde{d}_{1, \delta} = \tilde{d}_{1, \delta} \quad \mathrm{and} \quad \norm{\tilde{d}_{1, 1} - d_{1, 1}} < \delta,
 \end{equation}
and by \eqref{pert-C-step-2},
$$ \norm{\tilde{d}_{1, \delta} - d_{1, \delta}}_{2, \tr(A)} \approx_{O(\eps'')} \norm{\eta_\delta(u_2^* d_1 u_2) -  u^*_2\eta_\delta(d_1)u_2} = O(\eps'').$$
Then, with $\eps''$ sufficiently small, there is $n \in \mathbb N$ such that $$ \norm{(\tilde{d}_{1, \delta})^{\frac{1}{n}} d_{1, \delta} - d_{1, \delta}}_{2, \tr(A)} < \delta_1/4, $$
and hence, for all $\tau \in \tr(A)$,
\begin{eqnarray*}
\mathrm{d}_\tau((\tilde{d}_{1, \delta})^{\frac{1}{n}} d_{1, \delta}) + \delta_1/4 & \geq & \tau((\tilde{d}_{1, \delta})^{\frac{1}{n}} d_{1, \delta}) + \delta_1/4   \\
%& > & \tau(d_{1, \delta}) \\
%& > & \mathrm{d}_\tau(\tilde{d}_{1, \delta})   >  \tau(\tilde{d}_{1, \delta}) \\
& > & \tau(d_{1, \delta}) \approx_{5\eps''} \tau(u_2^* \eta_{\delta}(d_1)u_2) \quad \quad (\eqref{u-2-2-approx-C-1}) \\
& \approx_{O(\eps')} &  \tau(u^*_2\eta_{\delta}(u_1^* \psi_1 u_1 )u_2) \quad\quad (\eqref{approx-C-1}) \\
& \approx_{O(\eps' + \eps'')} & \tau(\eta_{\delta}(\psi_1 ))  > \tau(\psi_2) \\ 
& > & \mathrm{d}_\tau(\phi_3) + \delta_1/2.
\end{eqnarray*}
With $\eps'$ and $\eps''$ sufficiently small, one has 
$$ \mathrm{d}_\tau((\tilde{d}_{1, \delta})^{\frac{1}{n}} d_{1, \delta}) > \mathrm{d}_\tau(\phi_3),\quad \tau \in \tr(A). $$
  

Note that $(\tilde{d}_{1, \delta})^{\frac{1}{n}} d_{1, \delta} \in \overline{\phi_2D\phi_2}$ (both $(\tilde{d}_{1, \delta})^{\frac{1}{n}}$ and  $d_{1, \delta}$ belong to $D$, so they commute). By Property (C), for any $\eps'''>0$ (to be determined later), there is $u_3 \in \overline{\phi_2 A \phi_2} + \Comp 1$ such that 
\begin{equation}\label{approx-C-3-h}
u_3\phi_3u_3^* \in _{\eps'''}^{\norm{\cdot}_2} \overline{(\tilde{d}_{1, \delta})^{\frac{1}{n}} d_{1, \delta} D (\tilde{d}_{1, \delta})^{\frac{1}{n}} d_{1, \delta}},
\end{equation}
$$\mathrm{dist}_{2, \tr(A)}(u_3 d u^*_3, D_1) < \eps''', \quad \mathrm{dist}_{2, \tr(A)}(u_3^*du_3, D_1) < \eps''',\quad d \in D_1, $$
and
\begin{equation}\label{approx-C-3-u}
\norm{u_3u_3^* - 1}_{2, \tr(A)},\quad  \norm{u_3^*u_3 - 1}_{2, \tr(A)} < \eps'''.
\end{equation}

Note that, by \eqref{pert-fur-hereditary},
$$\tilde{d}_{1, 1}c = c,\quad c \in \overline{(\tilde{d}_{1, \delta})^{\frac{1}{n}} d_{1, \delta} A (\tilde{d}_{1, \delta})^{\frac{1}{n}} d_{1, \delta} }.$$
By \eqref{approx-C-3-h}, there is $c \in \overline{(\tilde{d}_{1, \delta})^{\frac{1}{n}} d_{1, \delta}D (\tilde{d}_{1, \delta})^{\frac{1}{n}} d_{1, \delta}}$ such that $\norm{c} \leq 1$ and
$$ \norm{c - u_3\phi_3 u_3^*}_{2, \tr(A)} < \eps'''.$$ 
Then 
\begin{eqnarray}\label{comp-approx}
 ( u^*_2\eta_{\eps}(u_1^* \psi_1 u_1)u_2) (u_3 \phi_3 u_3^*) & \approx_{O(\eps')}^{\norm{\cdot}_2} &  ( u^*_2\eta_{\eps}(d_1)u_2) (u_3 \phi_3 u_3^*) \quad\quad ( \eqref{approx-C-1} )   \\
 & \approx_{O(\eps'')}^{\norm{\cdot}_2} & \eta_{\eps}(u_2^*d_1 u_2) (u_3 \phi_3 u_3^*) \quad\quad (\eqref{approx-C-2-u})   \nonumber \\
 & \approx_{O(\eps'')}^{\norm{\cdot}_2} & \eta_{\eps}(d_{1, 1}) (u_3 \phi_3 u_3^*) \quad\quad (\eqref{pert-C-step-2})   \nonumber \\
 & \approx_{\eps/2+O(\eps'')}^{\norm{\cdot}_2} &  \eta_{\eps}(\tilde{d}_{1, 1}) (u_3 \phi_3 u_3^*) \quad \quad (\eqref{pert-fur-hereditary} \eqref{pre-pert-delta}) \nonumber \\
 & \approx_{\eps'''}^{\norm{\cdot}_2} &  \eta_{\eps}(\tilde{d}_{1, 1})  c \nonumber \\
 &=&c \approx_{\eps'''}^{\norm{\cdot}_2} u_3 \phi_3 u_3^*. \nonumber
\end{eqnarray} 


Then, consider the contraction $$ u = u_1u_2u_3.$$
Since 
$$u_2 \in \overline{\phi_1 A \phi_1}+\Comp 1\quad \mathrm{and}\quad u_3 \in \overline{\phi_2 A \phi_2} + \Comp 1, $$ one has
$$(u_1u_2u_3)^*\psi_1(u_1u_2u_3) \in \overline{\phi_1 A \phi_1}.$$
Moreover,
\begin{eqnarray*}
(u_1u_2u_3)^*\eta_{\eps_0/2}(\psi_1)(u_1u_2u_3)  & \approx_{O(\eps')}^{\norm{\cdot}_2} & u^*_3u_2^*\eta_{\eps_0/2}(u_1^*\psi_1 u_1)u_2 u_3 \quad\quad (\eqref{approx-C-1-u}) \\
& \approx_{O(\eps')}^{\norm{\cdot}_2} & u^*_3(u_2^*\eta_{\eps_0/2}(d_1)u_2) u_3 \quad\quad (\eqref{approx-C-1}) \\
& \in_{\eps''}^{\norm{\cdot}_2} & \overline{\phi_2 A \phi_2}, \quad\quad (\eqref{approx-C-2-here})
\end{eqnarray*}
and 
\begin{eqnarray*}
\eta_{\eps}((u_1u_2u_3)^*\psi_1(u_1u_2u_3)) \phi_3 & \approx_{O(\eps' + \eps'')}^{\norm{\cdot}_2} & u_3^*u_2^*\eta_\eps(u_1^*\psi_1 u_1)u_2u_3 \phi_3 \quad\quad (\eqref{approx-C-1-u}, \eqref{approx-C-2-u}) \\
 & \approx_{\eps'''}^{\norm{\cdot}_2} & u^*_3 (u^*_2  \eta_{\eps}(u_1^* \psi_1 u_1)u_2)(u_3 \phi_3 u_3^*) u_3 \quad\quad (\eqref{approx-C-3-u}) \\
 & \approx_{\eps + O(\eps'+\eps''+\eps''')} & u_3^* (u_3\phi_3 u_3^*) u_3= \phi_3. \quad\quad(\eqref{comp-approx})
 \end{eqnarray*}
With $\eps', \eps'', \eps'''$ sufficiently small, the contraction $u$ has the desired property.
\end{proof}

For technical reasons, we also need the following lemma which, very roughly, asserts that, after a perturbation with respect to the uniform trace norm, the spectrum of a positive element of the subalgebra $D$ is, in a strong sense, dense.
\begin{lem}\label{dense-sp-pert}
Let $A$ be a unital C*-algebra and let $D = \mathrm{C}(X) \subseteq A$ be a unital commutative sub-C*-algebra. Assume $A$ is simple, $X$ has no isolated points, and the pair $(D, A)$ has Property (E). Then, for any positive contraction $g \subseteq D$, any finite set $\{x_1, ..., x_n\} \subseteq (0, 1]$, and any $\eps>0$, there is a positive contraction $\tilde{g} \in D$ such that 
\begin{enumerate}
\item $\norm{\tilde{g} - g}_{2, \tr(A)} < \eps,$ and 
\item each point $x_i$, $i=1, ..., n$, is in $\mathrm{sp}(\tilde{g})$, and is not isolated from the left inside $\mathrm{sp}(\tilde{g})$ (i.e., $(s, x_i) \cap \mathrm{sp}(\tilde{g}) \neq \O$ for all $s < x_i$).
\end{enumerate}
\end{lem}
\begin{proof}
Since $A$ is simple (and non-elementary), there are mutually orthogonal positive elements $a_1, ..., a_n \in A$ such that $$\norm{a_i} = 1\quad  \mathrm{and} \quad \mathrm{d}_\tau(a_i) < \eps/n^2,\quad i=1, ..., n,\ \tau \in\tr(A). $$ (See, for instance, Lemma 4.7 of \cite{Niu-MD-Z}.)

Consider the contraction $$a= \frac{1}{n} a_1 + \cdots + \frac{n}{n} a_n,$$ and note that it has the property 
$$\quad 0 < \tau(h_i(a)) < \eps/n,\quad i=1, ..., n,\ \tau \in\tr(A),$$ where 
$h_i: [0, 1] \to [0, 1]$ is the continuous function taking value $1$ at $[(4i-1)/4n, (4i+1)/4n]$, $0$ on $[0, (2i-1)/2n]$ and $[(2i+1)/2n, 1]$, and linear between. By Property (E), there is a positive contraction $d \in D$ such that
$$\quad 0 < \tau(h_i(d)) < \eps/n,\quad i=1, ..., n,\ \tau \in\tr(A),$$
and there are mutually orthogonal positive elements $b_1, ..., b_n \in D$ such that
$$ \norm{b_i} = 1 \quad \mathrm{and} \quad \mathrm{d}_\tau(b_i)<\eps/n,\quad i=1, ..., n,\ \tau \in\tr(A). $$

Consider the sets  $$U_i = b^{-1}_i((0, 1]) \quad \mathrm{and} \quad V_i = b_i^{-1}((1/2, 1]),\quad i=1, ..., n. $$ Then
$$\overline{V_i} \subseteq U_i \quad\mathrm{and} \quad \mu_\tau(U_i) < \eps/n,\quad i=1, ..., n,\ \tau \in\tr(A).$$
For each $V_i$, $i=1, ..., n$, since $X$ has no isolated points, there is a continuous function $g_i: X \to [0, 1]$ such that $ g_i|_{V_i^c} = 0$ and $1$ is not isolated from the left in $g_i(X)$ (i.e., $(s, 1) \cap g_i(X) \neq \O$ for all $s<1$). Also pick a continuous function $r: X \to [0, 1]$ such that $$r|_{X \setminus (\bigcup_{i=1}^n U_i)} = 1\quad \mathrm{and}\quad r|_{\bigcup_{i=1}^n \overline{V_i}} = 0.$$
Then the function 
$$ \tilde{g}: = g r + (x_1g_1 + x_2g_2 + \cdots + x_ng_n)$$
has the desired property.
\end{proof}


We are now ready to prove the main theorem of the paper, which states that Property (S) of $A$ implies the (SBP) of $(D, \tr(A))$ if Properties (C) and (E) are present.


\begin{thm}\label{prop-S}
Let $A$ be a unital simple C*-algebra, and let $D = \mathrm{C}(X) \subseteq A$ be a unital commutative sub-C*-algebra such that $X$ has no isolated points. Assume that the pair $(D, A)$ has Properties (C) and (E). Then, if the C*-algebra $A$ has Property (S), the pair $(D, \tr(A))$ has the (SBP).
\end{thm}

\begin{proof}
By Theorem \ref{SBP-2-norm}, it is enough to show that for any self-adjoint contraction $f \in D$ and any $\eps>0$,  there is a self-adjoint element $g \in D$ such that
\begin{enumerate}
\item $\norm{f - g }_{2, \tr(A)} < \eps$, and
\item there is $\delta>0$ such that $\tau(\chi_\delta(g)) < \eps$, $\tau \in \tr(A)$.
\end{enumerate}


To show this statement, it is enough to prove it for $f$ such that $\mathrm{sp}(f) = [-1, 1]$. Indeed, set 
$$t_- = \sup\{t < 0: t \notin \mathrm{sp}(h)\} \quad \mathrm{and} \quad  t_+ = \inf\{t> 0: t \notin \mathrm{sp}(h)\}.$$

If $t_- = 0$ or $t_+ = 0$, then it is straightforward to perturb $f$ to produce $g$ (with $\chi_\delta(g) = 0$). Assume neither of $t_-$ and $t_+$ is zero. Choose $s_-,\  s_+ \notin \mathrm{sp}(f)$ such that $$0\leq t_- - s_- < \min\{\eps, -t_-\} \quad \mathrm{and} \quad 0\leq s_+ - t_+ < \min\{\eps, t_+\} $$ and consider the self-adjoint element $h(f)$ where
$$ h = \left\{
\begin{array}{ll}
0, & t < t_-, \\
t_- - s_-, & t \in [t_-, s_-], \\
t, & t \in [t_-, t+-], \\
s_+ - s_-, & t \in [t_+, s_+],\\
0 & t > s_+.
 \end{array}
\right. $$
Then $\mathrm{sp}(h(f)) = [t_-, t_+]$, and $$\norm{f - (f^-_{s_-} + h(f) +f^+_{s_+} )} < \eps,$$
where $f^-_{s_-} (t)= t $ if $t < s_-$ and $f^-_{s_-} (t)= 0$ otherwise, and $ f^+_{s_+} $ is defined similarly. Then, applying the statement to the self-adjoint element $h(f)$, one obtains the desired approximation $g$. 

Now, let us assume that $\mathrm{sp}(f) = [-1, 1]$. Identifying $[-1, 1]$ with $[0, 1]$, let us show the following (equivalent) statement: 

Let $f \in D$ be a positive contraction with $\mathrm{sp}(f) = [0, 1]$, and let $\eps>0$. Then there is a positive contraction $g \in D$ such that
\begin{equation}\label{positive-s1}
\norm{f - g }_{2, \tr(A)} < \eps,
\end{equation} 
and there is $\delta>0$ such that
\begin{equation}\label{positive-s2}
\tau(\chi_{\frac{1}{2}, \delta}(g)) < \eps,\quad \tau \in \tr(A),
\end{equation}
where 
$$ \chi_{\frac{1}{2}, \delta}(t) = \left\{
\begin{array}{ll}
0, & t < \frac{1}{2} - \delta, \\
\mathrm{linear}, & t \in [\frac{1}{2} - \delta, \frac{1}{2} - \frac{\delta}{2}],\\
1, & t \in [\frac{1}{2} - \frac{\delta}{2}, \frac{1}{2} + \frac{\delta}{2}],\\
\mathrm{linear}, & t \in [\frac{1}{2} + \frac{\delta}{2}, \frac{1}{2} + {\delta}],\\
0, & t > \frac{1}{2} + \delta.
\end{array}
\right. $$


Let $(f, \eps)$ be given. Applying Lemma \ref{hereditary-sets-tr-app} to $\eps/2$ (in place of $\eps$), we obtain $N$ and  $\eps_0$ (in place of $\delta$). Choose $\eps_1>0$ such that $$3\eps_1< 1/N,\quad 8N\eps_1 < \eps_0, \quad \mathrm{and} \quad  \eps_1 < \eps/4.$$



For each $i=1, 2, ..., N,$ consider the following functions
$$\chi_i(t) = \left\{ 
\begin{array}{ll}
0, & t \leq \frac{i-1}{N}, \\
\mathrm{linear}, & t \in [\frac{i-1}{N}, \frac{i}{N}], \\
1, & t \geq \frac{i}{N},
\end{array} 
\right.
\quad
\chi_{i, \eps_1}(t) = \left\{ 
\begin{array}{ll}
0, & t \leq \frac{i-1}{N} + \eps_1, \\
\mathrm{linear}, & t \in [\frac{i-1}{N} + \eps_1, \frac{i}{N}+\eps_1], \\
1, & t \geq \frac{i}{N}+\eps_1,
\end{array} 
\right.
$$

$$\kappa_{i, \eps_1}(t) = 
\left\{ 
\begin{array}{ll}
0, & t \leq \frac{i-1}{N} + \frac{\eps_1}{2}, \\
\mathrm{linear}, & t \in [\frac{i-1}{N} +\frac{\eps_1}{2}, \frac{i-1}{N} +\eps_1], \\
1, & t \geq \frac{i-1}{N} + \eps_1,
\end{array}
\right.
\quad
\theta_{i, \eps_1} = 
\left\{
\begin{array}{ll}
0, & t \leq \frac{i}{N} + \eps_1, \\
\textrm{linear}, & t \in [\frac{i}{N} + \eps_1, \frac{i}{N} + 2\eps_1], \\
1, & t \geq \frac{i}{N} + 2 \eps_1,
\end{array}
\right.
$$

$$
\xi_{i, \eps_1}^{+} = 
\left\{
\begin{array}{ll}
0, & t \leq \frac{i}{N}, \\
\textrm{linear}, & t \in [\frac{i}{N}, \frac{i}{N} + \frac{\eps_1}{2}], \\
1, & t \geq \frac{i}{N} + \frac{\eps_1}{2},
\end{array}
\right.
\quad 
\mathrm{and}
\quad
\xi_{i, \eps_1}^{-} = 
\left\{
\begin{array}{ll}
0, & t \leq \frac{i}{N} + 2\eps_1, \\
\textrm{linear}, & t \in [\frac{i}{N} + 2\eps_1, \frac{i}{N} + 3\eps_1], \\
1, & t \geq \frac{i}{N} + 3\eps_1.
\end{array}
\right.
$$



Consider the finite set of functions 
\begin{equation}
\mathcal H = \{\chi_i,\  \chi_{i, \eps_1},\ \kappa_{i, \eps_1},\ \theta_{i, \eps_1},\ \eta_{\eps_1/2}\circ \chi_{i, \eps_1}:\ i=1, 2, ..., N\}.
\end{equation}
Note that $$\eta_{\eps_1/2}(\chi_{i, \eps_1}(t)) = 0,\quad t \leq i/N+\eps_1/2,\ i=1, ..., N,$$
where, recall (\eqref{defn-eta}),
$$
\eta_\eps(t) = \left\{
\begin{array}{ll}
0, & t \leq 1-\eps, \\
\mathrm{linear}, & t \in [1-\eps, 1-\eps/2], \\
1, & t \geq 1-\eps/2.
\end{array}
\right.
$$


Since $A$ is simple and $\mathrm{sp}(f) = [0, 1]$, there is $\gamma > 0$ such that
\begin{eqnarray}\label{small-gap}
\gamma & < & \frac{1}{4}\min\{d_\tau(\chi_i(f)) - \tau(\kappa_{i, \eps_1}(f)), \ \tau(\theta_{i, \eps_1}(f)) - \mathrm{d}_\tau(\xi_{i, \eps_1}^-(f)),\\
& & \quad \quad \quad  \tau(\xi^+_{i, \eps_1}(f)) - \tau(\eta_{\eps_1/2}(\chi_{i, \eps_1}(f))): i=1, ..., N,\ \tau\in\tr(A)\}. \nonumber 
\end{eqnarray}
Without loss of generality, one may assume that $\gamma < \eps/4$. 


Since $A$ has Property (S), there is a positive contraction $\tilde{g} \in A$ such that $\norm{f - \tilde{g}}_{2, \tr(A)}$ is sufficiently small
that
\begin{equation}\label{1-pert}
\abs{ \tau(\chi(f) - \tau(\chi(\tilde{g}))} < \gamma, \quad \chi \in \mathcal H \subseteq \mathrm{C}([0, 1]),\ \tau \in \tr(A),
\end{equation}
and there is $\delta>0$ such that
\begin{equation}\label{1-small-bd}
\tau(\chi_{\frac{1}{2}, \delta}(\tilde{g})) < \eps/4,\quad \tau \in \tr(A).
\end{equation}
(Note that the choice of $\delta$ depends on $\tilde{g}$.) 



Applying Lemma \ref{hereditary-sets-preserve-tr} to $\eps/4$, $N$, and $\chi_{\frac{1}{2}+\eps_1, \delta}$ (in place of $\eps$, $N$, and $\chi$, respectively), one obtains $\delta_0$ (in place of $\delta$) and $M$ which have the property specified in Lemma \ref{hereditary-sets-preserve-tr} with respect to $\eps/2$, $N$, and $\chi_{\frac{1}{2}+\eps_1, \delta}$. Choose $\delta_1>0$ such that
$$4\delta_1 < \eps_0,\quad 6\delta_1 < \delta_0, \quad \mathrm{and} \quad 3^N\delta_1 < \min\{\eps_1, \gamma/2\} . $$

Also choose $\delta_2>0$ such that if $a, b$ are positive contractions of a C*-algebra $A$, then
\begin{equation}\label{choice-delta-2}
 \norm{a - b}_{2, \tr(A)} < \delta_2 \quad \Longrightarrow \quad \norm{\chi_{\frac{1}{2} + \eps_1, \delta}(a)) - \chi_{\frac{1}{2} + \eps_1, \delta}(b)}_{2, \tr(A)} < \eps/4.
\end{equation} 



Since $(D, A)$ has Property (E), there is a positive contraction $\tilde{\tilde{g}} \in D $ such that
\begin{equation}\label{2-pert}
\abs{\tau(\chi(\tilde{g})) - \tau(\chi(\tilde{\tilde{g}}))} < \gamma, \quad  \chi \in \mathcal H \cup \{ \chi_{\frac{1}{2}, \delta}\},\ \tau \in \tr(A).
\end{equation}
In particular, together with \eqref{1-pert} and \eqref{1-small-bd}, one has 
\begin{equation}\label{2-small-bd}
\abs{ \tau(\chi(f)) - \tau(\chi(\tilde{\tilde{g}})) } < 2 \gamma,\quad \chi \in \mathcal H,\ \tau \in \tr(A),
\end{equation}
and
\begin{equation}
\tau(\chi_{\frac{1}{2}, \delta}(\tilde{\tilde{g}})) < \eps/4+ \gamma < \eps/2.
\end{equation}
So
\begin{equation}\label{small-strech}
 \tau(\chi_{\frac{1}{2}+\eps_1, \delta}((\tilde{\tilde{g}} - \eps_1)_+))  < \eps/2. 
 \end{equation}

By Lemma \ref{dense-sp-pert}, after a small perturbation (with respect to $\norm{\cdot}_{2, \tr(A)}$), without loss of generality, one may assume that the numbers 
$$i/N + \eps_1,\quad i=1, 2, ..., N-1,$$
are in $\mathrm{sp}(\tilde{\tilde{g}})$, and are  not isolated from the left.



Now, let us consider the elements
$$\chi_1(f),\ \chi_2(f),\ ...,\ \chi_N \in D$$
and
$$\chi_{1, \eps_1}(\tilde{\tilde{g}}),\ \chi_{1, \eps_1}(\tilde{\tilde{g}}),\ ...,\ \chi_{1, \eps_1}(\tilde{\tilde{g}}) \in D.$$



By \eqref{2-pert}, one has
\begin{equation}
\mathrm{d}_\tau(\chi_{1, \eps_1}(\tilde{\tilde{g}})) < \tau(\kappa_{1, \eps_1}(\tilde{\tilde{g}})) \approx_{2 \gamma} \tau(\kappa_{1, \eps_1}(f)) < \mathrm{d}_\tau(\chi_1(f)),\quad \tau\in \tr(A).
\end{equation}

By the choice of $\gamma$ (\eqref{small-gap}), one has
\begin{equation}
\mathrm{d}_\tau(\chi_{1, \eps_1}(\tilde{\tilde{g}})) < \mathrm{d}_\tau(\chi_1(f)),\quad \tau\in\tr(A).
\end{equation}

Note that, by the construction of $\xi_{1, \eps_1}^+$, $\xi_{1, \eps_1}^-$, and $\theta_{1, \eps_1}$, we have 
$$   \xi_{1, \eps_1}^+(f),\ \xi_{1, \eps_1}^-(f) \in \overline{\chi_1(f) D \chi_1(f) },  $$

$$ \chi_{1, \eps_1}(\tilde{\tilde{g}}) \theta_{1, \eps_1}(\tilde{\tilde{g}}) = \theta_{1, \eps_1}(\tilde{\tilde{g}}),\quad \xi_{1, \eps_1}^+(f) \xi_{1, \eps_1}^-(f) = \xi_{1, \eps_1}^-(f),  $$

$$ \tau(\theta_{1, \eps_1}(\tilde{\tilde{g}})) \approx_{2\gamma} \tau(\theta_{1, \eps_1}(f))  >   \mathrm{d}_\tau(\xi_{1, \eps_1}^-( f )),\quad \tau \in \tr(A), $$
and
$$ \tau(\eta_{\eps_1/2}(\chi_{1, \eps_1}(\tilde{\tilde{g}}))) \approx_{2 \gamma} \tau(\eta_{\eps_1/2}(\chi_{1, \eps_1}(f))) < \tau( \xi_{1, \eps_1}^+(f)) < \mathrm{d}_\tau( \xi_{1, \eps_1}^+(f)),\quad \tau\in \tr(A). $$

By Lemma \ref{D-sandwich}, for any $\delta''>0$ (to be fixed later), there is a contraction $u_1 \in A$ such that
\begin{equation}\label{u-1-cond-1}
 u_1^*\chi_{1, \eps_1}(\tilde{\tilde{g}})u_1 \in^{\norm{\cdot}_2}_{\delta''}   \overline{\chi_1(f)A\chi_1(f)},
 \end{equation}
 %
\begin{equation}\label{u-1-cond-2}
\eta_{\eps_1/4}(u^*_1 \chi_{1, \eps_1}(\tilde{\tilde{g}}) u_1) \in^{\norm{\cdot}_2}_{\delta''} \overline{\xi_{1, \eps_1}^+(f) A \xi_{1, \eps_1}^+(f)},
\end{equation}
%
\begin{equation}\label{u-1-cond-3}
 \eta_{\delta''}(u_1^*\chi_{1, \eps_1}(\tilde{\tilde{g}})u_1) \xi_{1, \eps_1}^-(f)  \approx_{\delta''}^{\norm{\cdot}_2} \xi_{1, \eps_1}^-(f),
 \end{equation}
%and
\begin{equation}\label{u-1-cond-4}
\mathrm{dist}_{2, \tr(A)}(u_1du_1^*, (D)_1) < \delta'',\quad \mathrm{dist}_{2, \tr(A)}(u_1^*du_1, (D)_1) < \delta'',\quad d \in (D)_1,
\end{equation}
and
\begin{equation}\label{u-1-cond-5}
\norm{u_1u_1^* - 1}_{2, \tr(A)},\  \norm{u_1^*u_1 - 1}_{2, \tr(A)} < \delta''.
\end{equation}

With $\delta''$ sufficiently small, one has  
\begin{equation}\label{u-1-match-tr-p}
\abs{\tau((u_1^* x u_1)^j) - \tau(x^j)} < \delta_1,\quad j=1, ..., M,\  x\in (A)_1,\ \tau \in \tr(A),
\end{equation}
and (by \eqref{u-1-cond-4} and \eqref{u-1-cond-5})
\begin{equation}\label{u-1-close-D}
\mathrm{dist}_{2, \tr(A)}(u_1^*\chi_{1, \eps_1}(\tilde{\tilde{g}}) u_1, (D)^+_1) < 3\delta''< \min\{\eps_0, \delta_2, \eps/4\}.
\end{equation}

Set 
$$ \rho_1:=\min\{\tau(\rho_{1, \delta_1}(\tilde{\tilde{g}})): \tau \in \tr(A)\},$$
where
$$\rho_{1, \delta_1} = \left\{
\begin{array}{ll}
0, & t\leq \frac{1}{N}+\eps_1-\delta_1, \\
\mathrm{linear}, & \frac{1}{N} + \eps_1-\delta_1 \leq t \leq \frac{1}{N} + \eps_1-\delta_1/2, \\
1, & t = \frac{1}{N} + \eps_1-\delta_1/2, \\
\mathrm{linear}, & \frac{1}{N} + \eps_1-\delta_1/2 \leq t \leq \frac{1}{N} + \eps_1, \\
0, & t \geq \frac{1}{N}+\eps_1.
\end{array}
\right. 
$$
Since $1/N + \eps_1$ is not isolated from the left in $\mathrm{sp}(\tilde{\tilde{g}})$ and $A$ is simple, we have $\rho_1>0$.

By \eqref{u-1-close-D}, there is a positive contraction
$$ [u^*_1\chi_{1, \eps_1}(\tilde{\tilde{g}}) u_1] \in D $$
such that
\begin{equation}\label{down-to-D-1}
 \norm{ [u^*_1\chi_{1, \eps_1}(\tilde{\tilde{g}}) u_1] - u^*_1\chi_{1, \eps_1}(\tilde{\tilde{g}}) u_1 }_{2, \tr(A)} < 3\delta''<\delta_1. 
\end{equation}
With $\delta''$ sufficiently small, one has
$$\norm{ \eta_{\delta_1}([u^*_1\chi_{1, \eps_1}(\tilde{\tilde{g}}) u_1]) - u^*_1\eta_{\delta_1}(\chi_{1, \eps_1}(\tilde{\tilde{g}})) u_1 }_{2, \tr(A)} < \min\{\rho_1/8, \delta_1 \}.$$
Define
\begin{equation}\label{defn-down-to-D-m}
[u^*_1\eta_{\delta_1}(\chi_{1, \eps_1}(\tilde{\tilde{g}})) u_1] := \eta_{\delta_1}([u^*_1\chi_{1, \eps_1}(\tilde{\tilde{g}}) u_1]) \in D.
\end{equation}
Then
\begin{equation}\label{down-to-D-2}
 \norm{ [u^*_1\eta_{\delta_1}(\chi_{1, \eps_1}(\tilde{\tilde{g}})) u_1] - u^*_1\eta_{\delta_1}(\chi_{1, \eps_1}(\tilde{\tilde{g}})) u_1 }_{2, \tr(A)} < \min\{ \rho_1/8, \delta_1\}.
 \end{equation}
 Moreover (by \eqref{defn-down-to-D-m}), 
 \begin{equation}\label{down-to-D-3}
 \eta_{2\delta_1}([u^*_1\chi_{1, \eps_1}(\tilde{\tilde{g}}) u_1]) [u^*_1\eta_{\delta_1}(\chi_{1, \eps_1}(\tilde{\tilde{g}})) u_1]  = [u^*_1\eta_{\delta_1}(\chi_{1, \eps_1}(\tilde{\tilde{g}})) u_1]. 
 \end{equation}
 
 One should assume $\delta''$ is sufficiently small  that
 \begin{equation}\label{down-to-D-4} 
 \norm{ \eta_{\eps_1/4}([u_1^*\chi_{1, \eps_1}(\tilde{\tilde{g}})u_1]) - \eta_{\eps_1/4}(u_1^*\chi_{1, \eps_1}(\tilde{\tilde{g}})u_1) }_{2, \tr(A)} < \delta_1,
 \end{equation}
and
 \begin{equation}\label{down-to-D-5} 
 \norm{ \eta_{2\delta_1}([u_1^*\chi_{1, \eps_1}(\tilde{\tilde{g}})u_1]) - \eta_{2\delta_1}(u_1^*\chi_{1, \eps_1}(\tilde{\tilde{g}})u_1) }_{2, \tr(A)} < \delta_1. 
 \end{equation}
 
Note that (use \eqref{u-1-cond-3} in the third and fifth steps), 
\begin{eqnarray*}
 (u_1^* \chi_{1, \eps_1}(\tilde{\tilde{g}}) u_1)\chi_2(f)
 & \approx_{3N\eps_1} & (u_1^* \chi_{1, \eps_1}(\tilde{\tilde{g}}) u_1)\chi_{2, 3\eps_1}(f) \\
 &=&  (u_1^* \chi_{1, \eps_1}(\tilde{\tilde{g}}) u_1) \xi_{1, \eps_1}^-(f)\chi_{2, 3\eps_1}(f) \\
 & \approx_{\delta''}^{\norm{\cdot}_2} & (u_1^* \chi_{1, \eps_1}(\tilde{\tilde{g}}) u_1) \eta_{\delta''}(u_1^*\chi_{1, \eps_1}(\tilde{\tilde{g}})u_1) \xi_{1, \eps_1}^-(f)\chi_{2, 3\eps_1}(f) \\
& \approx_{\delta''} & \eta_{\delta''}(u_1^*\chi_{1, \eps_1}(\tilde{\tilde{g}})u_1) \xi_{1, \eps_1}^-(f)\chi_{2, 3\eps_1}(f)\\
 & \approx_{\delta''}^{\norm{\cdot}_2} & \xi_{1, \eps_1}^-(f)\chi_{2, 3\eps_1}(f) \\
 & = & \chi_{2, 3\eps_1}(f) \approx_{3N\eps_1} \chi_2(f).
 \end{eqnarray*}
With $\delta''$ sufficiently small, one has 
\begin{equation}\label{id-down-1}
[u^*_1\chi_{1, \eps_1}(\tilde{\tilde{g}}) u_1]\chi_2(f) \approx^{\norm{\cdot}_2}_{7N\eps_1} \chi_2(f). 
\end{equation}


%Consider 
%$$  \eta_{\delta'}(u_1^*\chi_{1, \eps''}(\tilde{\tilde{g}})u_1) \quad [u_1^* \eta_{\delta'}( \chi_{1, \eps''}(\tilde{\tilde{g}})) u_1]$$
%and its hereditary subalgebra.

Note that, by \eqref{down-to-D-2}, 
\begin{equation}\label{left-nbhd-1}
 \abs{\tau( [u_1^* \eta_{\delta_1}( \chi_{1, \eps_1}(\tilde{\tilde{g}})) u_1] ) - \tau( u_1^* \eta_{\delta_1}( \chi_{1, \eps_1}(\tilde{\tilde{g}})) u_1 )} < \rho_1/8,\quad \tau \in\tr(A).
 \end{equation}

With $\delta''$ sufficiently small, one has
\begin{equation}\label{left-nbhd-2}
 \tau(\eta_{\delta_1}(\chi_{1, \eps_1}(\tilde{\tilde{g}}))) \approx_{\rho_1/8} \tau(u^*_1\eta_{\delta_1}(\chi_{1, \eps_1}(\tilde{\tilde{g}}))u_1),\quad \tau \in \tr(A), 
 \end{equation}
and
\begin{equation}\label{swap-u-1}
\norm{ \eta_{\delta_1}(u_1^* \chi_{1, \eps_1}(\tilde{\tilde{g}})u_1) - u_1^* \eta_{\delta_1}(\chi_{1, \eps_1}(\tilde{\tilde{g}}))u_1}_2 < \delta_1/4.
\end{equation}

Hence, by \eqref{left-nbhd-1} and \eqref{left-nbhd-2},
\begin{eqnarray*}
\mathrm{d}_\tau(\chi_{2, \eps_1}(\tilde{\tilde{g}})) & \leq &  \tau(\eta_{\delta_1}(\chi_{1, \eps_1}(\tilde{\tilde{g}}))) -\rho_1/2 \\
& \approx_{\rho_1/8} & \tau(u^*_1\eta_{\delta_1}(\chi_{1, \eps_1}(\tilde{\tilde{g}}))u_1) - \rho_1/2 \\
& \approx_{\rho_1/8} & \tau([u^*_1\eta_{\delta_1}(\chi_{1, \eps_1}(\tilde{\tilde{g}}))u_1]) - \rho_1/2  \\
& < & \mathrm{d}_\tau([u^*_1\eta_{\delta_1}(\chi_{1, \eps_1}(\tilde{\tilde{g}}))u_1]) - \rho_1/2,
\end{eqnarray*}
for all $\tau \in \tr(A)$,
and therefore
\begin{equation}\label{lem-cond-1}
\mathrm{d}_\tau(\chi_{2, \eps_1}(\tilde{\tilde{g}})) < \mathrm{d}_\tau([u^*_1\eta_{\delta_1}(\chi_{1, \eps_1}(\tilde{\tilde{g}}))u_1]),\quad \tau \in\tr(A). 
\end{equation}





By \eqref{u-1-cond-3} (and \eqref{down-to-D-2}, \eqref{swap-u-1}), (note that $3\eps_1 < 1/N$ and $\delta''\ll \delta_1$)
\begin{eqnarray}\label{left-nbhd-3}
  [u^*_1\eta_{\delta_1}(\chi_{1, \eps_1}(\tilde{\tilde{g}}))u_1] (\xi^+_{2, \eps_1}(f)) & = & [u^*_1\eta_{\delta_1}(\chi_{1, \eps_1}(\tilde{\tilde{g}}))u_1]  (\xi^-_{1, \eps_1}(f))   (\xi^+_{2, \eps_1}(f))  \\
  & \approx_{\delta_1}^{\norm{\cdot}_2}& (u^*_1\eta_{\delta_1}(\chi_{1, \eps_1}(\tilde{\tilde{g}}))u_1)  (\xi^-_{1, \eps_1}(f))   (\xi^+_{2, \eps_1}(f)) \nonumber \\ 
  & \approx_{\delta_1}^{\norm{\cdot}_2}& \eta_{\delta_1}((u^*_1(\chi_{1, \eps_1}(\tilde{\tilde{g}}))u_1))  (\xi^-_{1, \eps_1}(f))   (\xi^+_{2, \eps_1}(f))  \nonumber \\ 
  & \approx_{\delta''}^{\norm{\cdot}_2} & (\xi^-_{1, \eps_1}(f))   (\xi^+_{2, \eps_1}(f)) \nonumber  \\
  & = &   \xi^+_{2, \eps_1}(f), \nonumber
  \end{eqnarray}
and therefore, one also has 
\begin{equation}\label{left-nbhd-4}
  [u^*_1\eta_{\delta_1}(\chi_{1, \eps_1}(\tilde{\tilde{g}}))u_1] (\xi^-_{2, \eps_1}(f)) \approx^{\norm{\cdot}_2}_{3\delta_1} \xi^-_{2, \eps_1}(f). 
\end{equation}  

Now, fixing $\delta''$, we have the contraction $u_1 \in A$.

Let us inductively assume that contractions $u_1, ..., u_k$, where $k \leq N-1$, have been constructed such that (note that \eqref{k-approx-bak} and \eqref{k-match-tr} are void if $k=1$)
\begin{equation}\label{k-approx-close-D}
\mathrm{dist}_{2, \tr(A)}(u_i^*\chi_{i, \eps_1}(\tilde{\tilde{g}}) u_i, (D)^+_1) < \min\{\eps_0, \delta_2, \eps/4\},\quad i=1, ..., k, \quad \quad \textrm{(\eqref{u-1-close-D} when $k=1$)}
\end{equation}
\begin{equation}\label{k-approx-bak}
 \norm{ \chi_{i-1}(f) (u_{i}^* \chi_{i, \eps_1}(\tilde{\tilde{g}}) u_{i}) - (u_{i}^* \chi_{i, \eps_1}(\tilde{\tilde{g}}) u_{i}) }_{2, \tr(A)} < 4 \delta_1 < \eps_1,\quad i=1, ..., k,
\end{equation}
\begin{equation}\label{k-approx-for}
\norm{ (u_{i}^* \chi_{i, \eps_1}(\tilde{\tilde{g}}) u_{i}) \chi_{i+1}(f) - \chi_{i+1}(f) }_{2, \tr(A)} < 7 N \eps_1 + 3^i\delta_1,\quad i=1, ..., k, \quad \quad \textrm{(\eqref{id-down-1} when $k=1$)}
\end{equation}
and
\begin{equation}\label{k-match-tr}
\norm{ (u_{i-1}^* \chi_{i, \eps_1}(\tilde{\tilde{g}}) u_{i-1})(u_{i}^* \chi_{i, \eps_1}(\tilde{\tilde{g}}) u_{i}) - (u_{i}^* \chi_{i, \eps_1}(\tilde{\tilde{g}}) u_{i}) }_{2, \tr(A)} < 6\delta_1,\quad i=1, ..., k, 
\end{equation}
\begin{equation}\label{k-match-tr-p}
\abs{\tau((u_i^* x u_i)^j) - \tau(x^j)} < \delta_1,\quad i =1, ..., k,\ j=1, ..., M,\  x\in (A)_1,\ \tau \in \tr(A). \quad \quad \textrm{(\eqref{u-1-match-tr-p} when $k=1$)} 
\end{equation}

Moreover, the contraction $u_k$ satisfies  
\begin{equation}\label{k-u-1-cond-2}
\eta_{\eps_1/4}(u^*_k \chi_{k, \eps_1}(\tilde{\tilde{g}}) u_k) \in^{\norm{\cdot}_2}_{\delta_1} \overline{\xi_{k, \eps_1}^+(f) A \xi_{k, \eps_1}^+(f)}, \quad \quad \textrm{(\eqref{u-1-cond-2} when $k=1$; note that $\delta''<\delta$)}
\end{equation}
and there are positive contractions $[u^*_k\chi_{k, \eps_1}(\tilde{\tilde{g}}) u_k], [u^*_k\eta_{\delta_1}(\chi_{k, \eps_1}(\tilde{\tilde{g}})) u_k] \in D$ such that
\begin{equation}\label{k-down-to-D-3}
 \eta_{2\delta_1}([u^*_k\chi_{k, \eps_1}(\tilde{\tilde{g}}) u_k]) [u^*_k\eta_{\delta_1}(\chi_{k, \eps_1}(\tilde{\tilde{g}})) u_k]  = [u^*_k\eta_{\delta_1}(\chi_{k, \eps_1}(\tilde{\tilde{g}})) u_k], \quad \quad \textrm{(\eqref{down-to-D-3} when $k=1$)}
 \end{equation}
 \begin{equation}\label{k-down-to-D-4} 
 \norm{ \eta_{\eps_1/4}([u_k^*\chi_{k, \eps_1}(\tilde{\tilde{g}})u_k]) - \eta_{\eps_1/4}(u_k^*\chi_{k, \eps_1}(\tilde{\tilde{g}})u_k) }_{2, \tr(A)} < \delta_1, \quad \quad \textrm{(\eqref{down-to-D-4} when $k=1$)}
 \end{equation}
 \begin{equation}\label{k-down-to-D-5} 
 \norm{ \eta_{2\delta_1}([u_k^*\chi_{k, \eps_1}(\tilde{\tilde{g}})u_k]) - \eta_{2\delta_1}(u_k^*\chi_{k, \eps_1}(\tilde{\tilde{g}})u_k) }_{2, \tr(A)} < \delta_1, \quad \quad \textrm{(\eqref{down-to-D-5} when $k=1$)}
 \end{equation}
\begin{equation}\label{k-lem-cond-1}
\mathrm{d}_\tau(\chi_{k+1, \eps_1}(\tilde{\tilde{g}})) < \mathrm{d}_\tau([u^*_k\eta_{\delta_1}(\chi_{k, \eps_1}(\tilde{\tilde{g}}))u_k]), \quad \tau \in\tr(A), \quad \quad \textrm{(\eqref{lem-cond-1} when $k=1$)}
\end{equation}
\begin{equation}\label{k-left-nbhd-3}
  [u^*_k\eta_{\delta_1}(\chi_{k, \eps_1}(\tilde{\tilde{g}}))u_k] (\xi^+_{k+1, \eps_1}(f)) \approx^{\norm{\cdot}_2}_{3^k\delta_1} \xi^+_{k+1, \eps_1}(f), \quad \quad \textrm{(\eqref{left-nbhd-3} when $k=1$)} 
\end{equation} 
and
\begin{equation}\label{k-left-nbhd-4}
  [u^*_k\eta_{\delta_1}(\chi_{k, \eps_1}(\tilde{\tilde{g}}))u_k] (\xi^-_{k+1, \eps_1}(f)) \approx^{\norm{\cdot}_2}_{3^k\delta_1} \xi^-_{k+1, \eps_1}(f). \quad \quad \textrm{(\eqref{left-nbhd-4} when $k=1$)}
\end{equation} 

Let us construct $u_{k+1}$.
Define
$$\left< \xi^-_{k+1, \eps_1}(f) \right> := [u^*_k\eta_{\delta_1}(\chi_{k, \eps_1}(\tilde{\tilde{g}}))u_k] (\xi^-_{k+1, \eps_1}(f)) \in \mathrm{Her}([u_k^* \eta_{\delta_1}(\chi_{k, \eps_1}(\tilde{\tilde{g}}))u_k] ) \cap D.$$
It follows from \eqref{k-left-nbhd-4} that
\begin{equation}\label{k-Xi-cut}
 \norm{\left< \xi^-_{k+1, \eps_1}(f) \right> - \xi^-_{k+1, \eps_1}(f)}_{2, \tr(A)} < 3^k\delta_1.  
 \end{equation}
Then, for all $\tau \in \tr(A)$, 
$$ \tau(\theta_{k+1, \eps_1}(f)) > d_\tau(\xi^-_{k+1, \eps_1}(f)) + 2 \gamma \geq d_\tau([u^*_k\eta_{\delta_1}(\chi_{k, \eps_1}(\tilde{\tilde{g}}))u_k]\xi^-_{k+1, \eps_1}(f)) + 2 \gamma  = \mathrm{d}_\tau(\left< \xi^-_{k+1, \eps_1}(f) \right>) + 2\gamma,$$
and hence 
$$\tau(\theta_{k+1, \eps_1}(\tilde{\tilde{g}})) \approx_{\gamma} \tau(\theta_{k+1, \eps_1}(f)) > \mathrm{d}_\tau( \left< \xi_{k+1, \eps_1}^-(f) \right>) + 2 \gamma. $$ 
In particular,
\begin{equation}\label{k-lem-cond-2}
\tau(\theta_{k+1, \eps_1}(\tilde{\tilde{g}})) >  \mathrm{d}_\tau( \left< \xi_{k+1, \eps_1}^-(f) \right> ),\quad \tau \in \tr(A). 
\end{equation}

Also define 
\begin{equation}\label{k+1-Xi-+}
\left< \xi^+_{k+1, \eps_1}(f) \right> := [u_k^*(\eta_{\delta_1}(\chi_{k, \eps_1}(\tilde{\tilde{g}}))u_k] (\xi^+_{k+1, \eps_1}(f)) \in \mathrm{Her}([u_k^*(\eta_{\delta_1}((\chi_{k, \eps_1}(\tilde{\tilde{g}})))u_k] ) \cap D.
\end{equation}
Then,
\begin{eqnarray*}
&& \tau(\eta_{\eps_1/2}(\chi_{k+1, \eps_1}( f )) ) \\
& < & \tau(\xi_{k+1, \eps_1}^+(f)) - 3 \gamma \quad \quad (\eqref{small-gap})\\
& \approx_{3^k\delta_1} & \tau( [u_k^*(\eta_{\delta_1}(\chi_{k, \eps_1}(\tilde{\tilde{g}}))u_k] (\xi^+_{k+1, \eps_1}(f)) ) - 3 \gamma \quad\quad (\eqref{k-left-nbhd-3}) \\
&= & \tau(  \left< \xi^+_{k+1, \eps_1}(f) \right> ) - 3 \gamma \\
& < & \mathrm{d}_\tau( \left< \xi^+_{k+1, \eps_1}(f) \right> ) - 3 \gamma, 
\end{eqnarray*}
and therefore
$$ \tau(\eta_{\eps_1/2}(\chi_{k+1, \eps_1}(\tilde{\tilde{g}}))) \approx_{2\gamma} \tau(\eta_{\eps_1/2}(\chi_{k+1, \eps_1}( f )) )  < d_\tau( [\xi^+_{k+1, \eps_1}(f)]) - 3 \gamma+ 3^k\delta_1.$$
In particular  (note that $3^N\delta_1 < \gamma /2$), 
\begin{equation}\label{k-lem-cond-3}
\tau(\eta_{\eps_1/2}(\chi_{k+1, \eps_1}(\tilde{\tilde{g}}))) < d_\tau( [\xi^+_{k+1, \eps_1}(f)]),\quad \tau \in\tr(A).
\end{equation}

With \eqref{k-lem-cond-1}, \eqref{k-lem-cond-2}, and \eqref{k-lem-cond-3}, by Lemma \ref{D-sandwich}, for any $\delta''>0$ (to be fixed later), there is a contraction $u_{k+1} \in A$ such that
\begin{equation}\label{k-u-2-cond-1}
 u_{k+1}^*\chi_{k+1, \eps_1}(\tilde{\tilde{g}})u_{k+1} \in^{\norm{\cdot}_{2}}_{\delta''}   \overline{[u_k^*\eta_{\delta'}(\chi_{k, \eps_1}(\tilde{\tilde{g}}))u_k] A [u_k^*\eta_{\delta'}(\chi_{k, \eps_1}(\tilde{\tilde{g}}))u_k] },
 \end{equation}
 %
\begin{equation}\label{k-u-2-cond-2}
\eta_{\eps_1/4}(u^*_{k+1} \chi_{k+1, \eps_1}(\tilde{\tilde{g}}) u_{k+1}) \in^{\norm{\cdot}_2}_{\delta''} \overline{\left< \xi_{k+1, \eps_1}^+(f) \right> A \left< \xi_{k+1, \eps_1}^+(f) \right>},
\end{equation}
%
\begin{equation}\label{k-u-2-cond-3}
 \eta_{\delta''}(u_{k+1}^*\chi_{k+1, \eps_1}(\tilde{\tilde{g}})u_{k+1}) \left< \xi_{k+1, \eps_1}^-(f) \right>  \approx_{\delta''}^{\norm{\cdot}_2} \left< \xi_{k+1, \eps_1}^-(f) \right>,
 \end{equation}
%
\begin{equation}\label{u-2-cond-4}
\mathrm{dist}_{2, \tr(A)}(u_{k+1}du_{k+1}^*, D_1) < \delta'',\quad \mathrm{dist}_{2, \tr(A)}(u_{k+1}^*du_{k+1}, D_1) < \delta'',\quad d \in D_1,
\end{equation}
and
\begin{equation}\label{u-2-cond-5}
\norm{u_{k+1}u_{k+1}^* - 1}_{2, \tr(A)},\  \norm{u_{k+1}^*u_{k+1} - 1}_{2, \tr(A)} < \delta''.
\end{equation}
 
By \eqref{u-2-cond-4} and \eqref{u-2-cond-5}, with $\delta''$ sufficiently small,
\begin{equation}\label{k+1-approx-close-D}
\mathrm{dist}_{2, \tr(A)}(u_{k+1}^*\chi_{k+1, \eps_1}(\tilde{\tilde{g}}) u_{k+1}, (D)^+_1) < 3\delta'' < \min\{\eps_0, \delta_2, \eps/2\}.
\end{equation}
This verifies Assumption \eqref{k-approx-close-D} for $k+1$.

 Also note 
 \begin{eqnarray*}
 & & u_{k+1}^*( \chi_{k+1, \eps_1}(\tilde{\tilde{g}})) u_{k+1} \\
 & \approx_{\delta''}^{\norm{\cdot}_2} & (u_{k+1}^*( \chi_{k+1, \eps_1}(\tilde{\tilde{g}})) u_{k+1}) ( \eta_{2\delta_1}([u_k^* \chi_{k, \eps_1}(\tilde{\tilde{g}}) u_k] )) \quad \quad (\eqref{k-u-2-cond-1}, \eqref{k-down-to-D-3}) \\
 & \approx_{\delta_1}^{\norm{\cdot}_2} & (u_{k+1}^*( \chi_{k+1, \eps_1}(\tilde{\tilde{g}})) u_{k+1}) ( \eta_{2\delta_1}(u_k^* \chi_{k, \eps_1}(\tilde{\tilde{g}}) u_k )) \quad \quad (\eqref{k-down-to-D-5}) \\
 & \approx_{2\delta_1}^{\norm{\cdot}_2} & (u_{k+1}^*( \chi_{k+1, \eps_1}(\tilde{\tilde{g}})) u_{k+1}) ( \eta_{2\delta_1}(u_k^* \chi_{k, \eps_1}(\tilde{\tilde{g}}) u_k )) (u_k\chi_{k, \eps_1}(\tilde{\tilde{g}}) u_k) \\
 &\approx_{\delta_1+\delta''}& (u_{k+1}^*( \chi_{k+1, \eps_1}(\tilde{\tilde{g}})) u_{k+1}) (u^*_k\chi_{k, \eps_1}(\tilde{\tilde{g}}) u_k). \quad \quad (\eqref{k-u-2-cond-1}, \eqref{k-down-to-D-3}, \eqref{k-down-to-D-5})
\end{eqnarray*}
In other words,
\begin{equation*}
 (u_k^* \chi_{k, \eps_1}(\tilde{\tilde{g}}) u_k)(u_{k+1}^*(\chi_{k+1, \eps_1}(\tilde{\tilde{g}}) )u_{k+1}) \approx_{6\delta_1}^{\norm{\cdot}_2} u_{k+1}^*(\chi_{k+1, \eps_1}(\tilde{\tilde{g}}) )u_{k+1}.
 \end{equation*}
 This verifies Assumption \eqref{k-match-tr} for $k+1$.
 
 By \eqref{k-u-1-cond-2} and \eqref{k-u-2-cond-1}, \eqref{k-down-to-D-3}, \eqref{k-down-to-D-4},
\begin{eqnarray*}
\chi_k(f) (u_{k+1}^*(\chi_{k+1, \eps_1}(\tilde{\tilde{g}}) )u_{k+1}) & \approx_{\delta''} & \chi_k(f) \eta_{2\delta_1}([u_k^*\chi_{k, \eps_1}(\tilde{\tilde{g}})u_k]) (u_{k+1}^*(\chi_{k+1, \eps_1}(\tilde{\tilde{g}}) )u_{k+1}) \quad \quad (\eqref{k-down-to-D-3}) \\
& = & \chi_k(f) \eta_{\eps_1/4}([u_k^*\chi_{k, \eps_1}(\tilde{\tilde{g}})u_k]) (u_{k+1}^*(\chi_{k+1, \eps_1}(\tilde{\tilde{g}}) )u_{k+1})\\
& \approx^{\norm{\cdot}_2}_{\delta_1} & (\chi_k(f) \eta_{\eps_1/4}(u_k^*\chi_{k, \eps_1}(\tilde{\tilde{g}})u_k)) (u_{k+1}^*(\chi_{k+1, \eps_1}(\tilde{\tilde{g}}) )u_{k+1}) \quad \quad (\eqref{k-down-to-D-4}) \\
& \approx^{\norm{\cdot}_2}_{\delta_1} & \eta_{\eps_1/4}(u_k^*\chi_{k, \eps_1}(\tilde{\tilde{g}})u_k) (u_{k+1}^*(\chi_{k+1, \eps_1}(\tilde{\tilde{g}}) )u_{k+1}) \quad \quad (\eqref{k-u-1-cond-2}) \\
& \approx^{\norm{\cdot}_2}_{\delta_1} & \eta_{\eps_1/4}([u_k^*\chi_{k, \eps_1}(\tilde{\tilde{g}})u_k]) (u_{k+1}^*(\chi_{k+1, \eps_1}(\tilde{\tilde{g}}) )u_{k+1}) \quad \quad (\eqref{k-down-to-D-4}) \\
& = & u_{k+1}^*(\chi_{k+1, \eps_1}(\tilde{\tilde{g}}) )u_{k+1}.
\end{eqnarray*}
So,
$$ \chi_k(f) (u_{k+1}^*(\chi_{k+1, \eps_1}(\tilde{\tilde{g}}) )u_{k+1}) \approx^{\norm{\cdot}_2}_{4\delta_1} u_{k+1}^*(\chi_{k+1, \eps_1}(\tilde{\tilde{g}}) )u_{k+1}.  $$
This verifies Assumption \eqref{k-approx-bak} for $k+1$.

If $k+1 \leq N-1$, with the same argument as for \eqref{id-down-1}, 
\begin{eqnarray*}
 && (u_{k+1}^* \chi_{{k+1}, \eps_1}(\tilde{\tilde{g}}) u_{k+1})\chi_{k+2}(f) \\
 & \approx_{3N\eps_1} & (u_{k+1}^* \chi_{k+1, \eps_1}(\tilde{\tilde{g}}) u_{k+1})\chi_{k+2, 3\eps_1}(f) \\
 &=&  (u_{k+1}^* \chi_{k+1, \eps_1}(\tilde{\tilde{g}}) u_{k+1}) \xi_{k+1, \eps_1}^-(f)\chi_{k+2, 3\eps_1}(f) \\
 &\approx^{\norm{\cdot}_2}_{3^k\delta_1}&  (u_{k+1}^* \chi_{k+1, \eps_1}(\tilde{\tilde{g}}) u_{k+1}) \left< \xi_{k+1, \eps_1}^-(f) \right> \chi_{k+2, 3\eps_1}(f) \quad \quad (\eqref{k-Xi-cut}) \\
 & \approx_{\delta''}^{\norm{\cdot}_2} & (u_{k+1}^* \chi_{k+1, \eps_1}(\tilde{\tilde{g}}) u_{k+1}) \eta_{\delta''}(u_{k+1}^*\chi_{k+1, \eps_1}(\tilde{\tilde{g}})u_{k+1}) \left< \xi_{k+1, \eps_1}^-(f) \right> \chi_{k+2, 3\eps_1}(f) \quad \quad (\eqref{k-u-2-cond-3}) \\
& \approx_{\delta''} & \eta_{\delta''}(u_{k+1}^*\chi_{k+1, \eps_1}(\tilde{\tilde{g}})u_{k+1}) \left< \xi_{k+1, \eps_1}^-(f)\right> \chi_{k+2, 3\eps_1}(f)\\
 & \approx_{\delta''}^{\norm{\cdot}_2} &  \left<  \xi_{k+1, \eps_1}^-(f) \right> \chi_{k+2, 3\eps_1}(f) \quad \quad ( \eqref{k-u-2-cond-3}) \\
 & \approx^{\norm{\cdot}_2}_{3^k\delta_1} & \chi_{k+2, 3\eps_1}(f) \approx_{3N\eps_1} \chi_{k+2}(f).  \quad \quad (\eqref{k-Xi-cut})
 \end{eqnarray*}
Thus, 
\begin{equation*}%\label{k-id-down-2} 
 (u_{k+1}^*(\chi_{k+1, \eps_1}(\tilde{\tilde{g}}) )u_{k+1}) \chi_{k+2}(f) \approx_{7N\eps_1 + 3^{k+1}\delta_1}^{\norm{\cdot}_2} \chi_{k+2}(f).
\end{equation*}
This verifies Assumption \eqref{k-approx-for} for $k+1$ (which is void if $k+1 = N$).

With $\delta''$ sufficiently small, one has  
\begin{equation*}
\abs{\tau((u_{k+1}^* x u_{k+1})^j) - \tau(x^j)} < \delta_1,\quad j=1, ..., M,\  x\in (A)_1,\ \tau \in \tr(A). 
\end{equation*}
This verifies \eqref{k-match-tr-p} for $k+1$.

 
If $k+1 \leq N-1$, let us verify that (with $\delta''$ sufficiently small), the contraction $u_{k+1}$ satisfies the inductive assumptions \eqref{k-u-1-cond-2}, \eqref{k-down-to-D-3}, \eqref{k-down-to-D-4}, \eqref{k-down-to-D-5}, \eqref{k-lem-cond-1}, \eqref{k-left-nbhd-3}, and \eqref{k-left-nbhd-4} for $k+1$.
 
By \eqref{k+1-Xi-+} and noting that $ [u_k^*(\eta_{\delta_1}(\chi_{k, \eps_1}(\tilde{\tilde{g}}))u_k]$ and $ (\xi^+_{k+1, \eps_1}(f))$ commute (both are in $D$), one has
$$ \overline{\left< \xi_{k+1, \eps_1}^+(f) \right> A \left< \xi_{k+1, \eps_1}^+(f) \right>} \subseteq \overline{ \xi_{k+1, \eps_1}^+(f) A  \xi_{k+1, \eps_1}^+(f) }.$$
 Thus, Assumption \eqref{k-u-1-cond-2} for $k+1$ follows from \eqref{k-u-2-cond-2}. 
 
For the other assumptions, let us repeat the argument of $u_1$:  Set 
$$ \rho_{k+1}:=\min\{\tau(\rho_{k+1, \delta_1}(\tilde{\tilde{g}})): \tau \in \tr(A)\},$$
where
$$\rho_{k+1, \delta_1} = \left\{
\begin{array}{ll}
0, & t\leq \frac{k+1}{N}+\eps_1-\delta_1, \\
\mathrm{linear}, & \frac{k+1}{N} + \eps_1-\delta_1 \leq t \leq \frac{k+1}{N} + \eps_1-\delta_1/2, \\
1, & t = \frac{k+1}{N} + \eps_1-\delta_1/2, \\
\mathrm{linear}, & \frac{k+1}{N} + \eps_1-\delta_1/2 \leq t \leq \frac{k+1}{N} + \eps_1, \\
0, & t \geq \frac{k+1}{N}+\eps_1.
\end{array}
\right. 
$$
Since $(k+1)/N + \eps_1$ is not isolated from the left in $\mathrm{sp}(\tilde{\tilde{g}})$ and $A$ is simple, we have that $\rho_1>0$.

By \eqref{k+1-approx-close-D}, with a sufficiently small $\delta''$, there is a positive contraction 
$$ [u^*_{k+1}\chi_{k+1, \eps_1}(\tilde{\tilde{g}}) u_{k+1}] \in D $$
such that
\begin{equation*}%\label{down-to-D-1}
 \norm{ [u^*_{k+1}\chi_{k+1, \eps_1}(\tilde{\tilde{g}}) u_{k+1}] - u^*_{k+1}\chi_{k+1, \eps_1}(\tilde{\tilde{g}}) u_{k+1} }_{2, \tr(A)} < 3\delta''<\delta_1. 
\end{equation*}
With $\delta''$ sufficiently small, one has
$$\norm{ \eta_{\delta_1}([u^*_{k+1}\chi_{k+1, \eps_1}(\tilde{\tilde{g}}) u_{k+1}]) - u^*_{k+1}\eta_{\delta_1}(\chi_{k+1, \eps_1}(\tilde{\tilde{g}})) u_{k+1} }_{2, \tr(A)} < \min\{\rho_{k+1}/8, \delta_1 \}.$$
Define
\begin{equation}\label{k+1-defn-down-to-D-m}
[u^*_{k+1}\eta_{\delta_1}(\chi_{{k+1}, \eps_1}(\tilde{\tilde{g}})) u_{k+1}] := \eta_{\delta_1}([u^*_{k+1}\chi_{k+1, \eps_1}(\tilde{\tilde{g}}) u_{k+1}]) \in D.
\end{equation}
Then
\begin{equation}\label{k+1-down-to-D-2}
 \norm{ [u^*_{k+1}\eta_{\delta_1}(\chi_{k+1, \eps_1}(\tilde{\tilde{g}})) u_{k+1}] - u^*_{k+1}\eta_{\delta_1}(\chi_{k+1, \eps_1}(\tilde{\tilde{g}})) u_{k+1} }_{2, \tr(A)} < \min\{ \rho_{k+1}/8, \delta_1\}.
 \end{equation}
 Moreover (by \eqref{k+1-defn-down-to-D-m}), 
 \begin{equation}\label{k+1-down-to-D-3}
 \eta_{2\delta_1}([u^*_{k+1}\chi_{k+1, \eps_1}(\tilde{\tilde{g}}) u_{k+1}]) [u^*_{k+1}\eta_{\delta_1}(\chi_{k+1, \eps_1}(\tilde{\tilde{g}})) u_{k+1}]  = [u^*_{k+1}\eta_{\delta_1}(\chi_{k+1, \eps_1}(\tilde{\tilde{g}})) u_{k+1}]. 
 \end{equation}
 This verifies %\eqref{k-down-to-D-2} and 
 \eqref{k-down-to-D-3} for $k+1$.
 
 One should assume $\delta''$ is sufficiently small  that
 \begin{equation}\label{k+1-down-to-D-4} 
 \norm{ \eta_{\eps_1/4}([u_{k+1}^*\chi_{k+1, \eps_1}(\tilde{\tilde{g}})u_{k+1}]) - \eta_{\eps_1/4}(u_{k+1}^*\chi_{k+1, \eps_1}(\tilde{\tilde{g}})u_{k+1}) }_{2, \tr(A)} < \delta_1,
 \end{equation}
and
 \begin{equation}\label{k+1-down-to-D-5} 
 \norm{ \eta_{2\delta_1}([u_{k+1}^*\chi_{k+1, \eps_1}(\tilde{\tilde{g}})u_{k+1}]) - \eta_{2\delta_1}(u_{k+1}^*\chi_{k+1, \eps_1}(\tilde{\tilde{g}})u_{k+1}) }_{2, \tr(A)} < \delta_1. 
 \end{equation}
 This verifies \eqref{k-down-to-D-4} and \eqref{k-down-to-D-5} for $k+1$.
 
 Note that, by \eqref{k+1-down-to-D-2}, 
\begin{equation}\label{k+1-left-nbhd-1}
 \abs{\tau( [u_{k+1}^* \eta_{\delta_1}( \chi_{k+1, \eps_1}(\tilde{\tilde{g}})) u_{k+1}] ) - \tau( u_{k+1}^* \eta_{\delta_1}( \chi_{k+1, \eps_1}(\tilde{\tilde{g}})) u_{k+1} )} < \rho_{k+1}/8,\quad \tau \in\tr(A).
 \end{equation}
 
 With $\delta''$ sufficiently small, one has
\begin{equation}\label{k+1-left-nbhd-2}
 \tau(\eta_{\delta_1}(\chi_{k+1, \eps_1}(\tilde{\tilde{g}}))) \approx_{\rho_{k+1}/8} \tau(u^*_{k+1}\eta_{\delta_1}(\chi_{k+1, \eps_1}(\tilde{\tilde{g}}))u_{k+1}),\quad \tau \in \tr(A), 
 \end{equation}
and
\begin{equation}\label{k+1-swap-u-1}
\norm{ \eta_{\delta_1}(u_{k+1}^* \chi_{k+1, \eps_1}(\tilde{\tilde{g}})u_{k+1}) - u_{k+1}^* \eta_{\delta_1}(\chi_{k+1, \eps_1}(\tilde{\tilde{g}}))u_{k+1}}_{2, \tr(A)} < \delta_1/4.
\end{equation}
Hence, by \eqref{k+1-left-nbhd-1} and \eqref{k+1-left-nbhd-2},
\begin{eqnarray*}
\mathrm{d}_\tau(\chi_{k+2, \eps_1}(\tilde{\tilde{g}})) & \leq &  \tau(\eta_{\delta_1}(\chi_{k+1, \eps_1}(\tilde{\tilde{g}}))) -\rho_{k+1}/2 \\
& \approx_{\rho_{k+1}/8} & \tau(u^*_{k+1}\eta_{\delta_1}(\chi_{k+1, \eps_1}(\tilde{\tilde{g}}))u_{k+1}) - \rho_{k+1}/2 \\
& \approx_{\rho_{k+1}/8} & \tau([u^*_{k+1}\eta_{\delta_1}(\chi_{k+1, \eps_1}(\tilde{\tilde{g}}))u_{k+1}]) - \rho_{k+1}/2  \\
& < & \mathrm{d}_\tau([u^*_{k+1}\eta_{\delta_1}(\chi_{k+1, \eps_1}(\tilde{\tilde{g}}))u_{k+1}]) - \rho_{k+1}/2,
\end{eqnarray*}
for all $\tau \in \tr(A)$,
and therefore
\begin{equation}\label{k+1-lem-cond-1}
\mathrm{d}_\tau(\chi_{k+2, \eps_1}(\tilde{\tilde{g}})) < \mathrm{d}_\tau([u^*_{k+1}\eta_{\delta_1}(\chi_{k+1, \eps_1}(\tilde{\tilde{g}}))u_{k+1}]). 
\end{equation}
This verifies \eqref{k-lem-cond-1} for $k+1$.

By \eqref{k-u-2-cond-3} and \eqref{k-Xi-cut} (and \eqref{k+1-down-to-D-2}, \eqref{k+1-swap-u-1}), (note that $3\eps_1 < 1/N$ and $\delta''\ll \delta_1$)
\begin{eqnarray}\label{k+1-left-nbhd-3}
  & & [u^*_{k+1}\eta_{\delta_1}(\chi_{k+1, \eps_1}(\tilde{\tilde{g}}))u_{k+1}] (\xi^+_{k+2, \eps_1}(f)) \\
  & = & [u^*_{k+1}\eta_{\delta_1}(\chi_{k+1, \eps_1}(\tilde{\tilde{g}}))u_{k+1}]  (\xi^-_{k+1, \eps_1}(f))   (\xi^+_{k+2, \eps_1}(f))  \nonumber \\
  & \approx_{\delta_1}^{\norm{\cdot}_2}& (u^*_{k+1}\eta_{\delta_1}(\chi_{k+1, \eps_1}(\tilde{\tilde{g}}))u_{k+1})  (\xi^-_{k+1, \eps_1}(f))   (\xi^+_{k+2, \eps_1}(f))  \quad\quad ( \eqref{k+1-down-to-D-2}) \nonumber \\ 
  & \approx_{\delta_1}^{\norm{\cdot}_2}& \eta_{\delta_1}((u^*_{k+1}(\chi_{k+1, \eps_1}(\tilde{\tilde{g}}))u_{k+1}))  (\xi^-_{k+1, \eps_1}(f))   (\xi^+_{k+2, \eps_1}(f)) \quad \quad (\eqref{k+1-swap-u-1}) \nonumber \\ 
  & \approx_{3^k \delta_1}^{\norm{\cdot}_2}& \eta_{\delta_1}((u^*_{k+1}(\chi_{k+1, \eps_1}(\tilde{\tilde{g}}))u_{k+1}))  \left< \xi^-_{k+1, \eps_1}(f) \right>   (\xi^+_{k+2, \eps_1}(f)) \quad \quad (\eqref{k-Xi-cut}) \nonumber \\
  & \approx_{\delta''}^{\norm{\cdot}_2} & \left< \xi^-_{k+1, \eps_1}(f) \right>   (\xi^+_{k+2, \eps_1}(f)) \quad \quad (\eqref{k-u-2-cond-3}) \nonumber  \\
  & \approx_{3^k \delta_1}^{\norm{\cdot}_2} & ( \xi^-_{k+1, \eps_1}(f) )   (\xi^+_{k+2, \eps_1}(f)) \quad \quad (\eqref{k-Xi-cut}) \nonumber  \\
  & = &   \xi^+_{k+2, \eps_1}(f), \nonumber
  \end{eqnarray}
and therefore, one also has 
\begin{equation}\label{k+1-left-nbhd-4}
  [u^*_{k+1}\eta_{\delta_1}(\chi_{k+1, \eps_1}(\tilde{\tilde{g}}))u_{k+1}] (\xi^-_{k+2, \eps_1}(f)) \approx^{\norm{\cdot}_2}_{3^{k+1}\delta_1} \xi^-_{k+2, \eps_1}(f). 
\end{equation}  
This verifies \eqref{k-left-nbhd-3} and \eqref{k-left-nbhd-4} for $k+1$. Fix $\delta''$, and we obtain the desired $u_{k+1}$.

By induction, there are contractions $u_1, u_2, ..., u_{N} \in A$ such that (note that $3^N\delta_1 < \eps_1$)
\begin{equation}\label{N-approx-close-D}
\mathrm{dist}_{2, \tr(A)}(u_i^*\chi_{i, \eps_1}(\tilde{\tilde{g}}) u_i, (D)^+_1) < \min\{\eps_0, \delta_2, \eps/4\} \leq \eps_0,\quad i=1, ..., N,
\end{equation}
\begin{equation}\label{N-approx-bak}
 \norm{ \chi_{i-1}(f) (u_{i}^* \chi_{i, \eps_1}(\tilde{\tilde{g}}) u_{i}) - (u_{i}^* \chi_{i, \eps_1}(\tilde{\tilde{g}}) u_{i}) }_{2, \tr(A)} < 4 \delta_1 < \eps_0,\quad i=2, ..., N,
\end{equation}
\begin{equation}\label{N-approx-for}
\norm{ (u_{i}^* \chi_{i, \eps_1}(\tilde{\tilde{g}}) u_{i}) \chi_{i+1}(f) - \chi_{i+1}(f) }_{2, \tr(A)} < 7 N \eps_1 + 3^i\delta_1 < 8N\eps_1 < \eps_0,\quad i=1, ..., N-1, 
\end{equation}
%and
\begin{equation}\label{N-match-tr}
\norm{ (u_{i-1}^* \chi_{i, \eps_1}(\tilde{\tilde{g}}) u_{i-1})(u_{i}^* \chi_{i, \eps_1}(\tilde{\tilde{g}}) u_{i}) - (u_{i}^* \chi_{i, \eps_1}(\tilde{\tilde{g}}) u_{i}) }_{2, \tr(A)} < 6\delta_1 < \delta_0,\quad i=1, ..., N, 
\end{equation}
and
\begin{equation}\label{N-match-tr-p}
\abs{\tau((u_i^* x u_i)^j) - \tau(x^j)} < \delta_1 < \delta_0,\quad i =1, ..., N,\ j=1, ..., M,\  x\in (A)_1,\ \tau \in \tr(A).\end{equation}

Define
$$g:= \frac{1}{N}(u_1^*\chi_{1, \eps_1}(\tilde{\tilde{g}}) u_1 + \cdots + u_N^*\chi_{N, \eps_1}(\tilde{\tilde{g}}) u_N).$$
Then, by \eqref{N-approx-close-D}, \eqref{N-approx-bak} and \eqref{N-approx-for}, it follows from Lemma \ref{hereditary-sets-tr-app} that
\begin{equation}\label{almost-g-1}
 \norm{ f - g}_{2, \tr(A)} < \eps/2.
 \end{equation}

Note that
$$(\tilde{\tilde{g}} - \eps_1)_+ = \frac{1}{N}(\chi_{1, \eps_1}(\tilde{\tilde{g}}) + \cdots + \chi_{N, \eps_1}(\tilde{\tilde{g}}) ).$$
By \eqref{N-match-tr} and  \eqref{N-match-tr-p}, it follows from Lemma \ref{hereditary-sets-preserve-tr} that 
\begin{equation}\label{almost-g-2}
\abs{ \tau(\chi_{\frac{1}{2}+\eps_1, \delta}((\tilde{\tilde{g}} - \eps_1))) - \tau(\chi_{\frac{1}{2}+\eps_1, \delta}(g))} < \eps/4,\quad \tau\in \tr(A). 
\end{equation}

By \ref{N-approx-close-D} again, one has
$$\mathrm{dist}_{2, \tr(A)}(g, (D)_1^+) < \min\{\delta_2, \eps/4\};$$
together with \eqref{almost-g-1}, \eqref{almost-g-2}, and the choice of $\delta_2$ (\eqref{choice-delta-2}),
there is a positive contraction in $D$, still denoted by $g$, such that 
$$\norm{f - g }_{2, \tr(A)} < \eps/2 + \eps/4 = 3\eps/4$$ 
and,
$$ \abs{ \tau(\chi_{\frac{1}{2}+\eps_1, \delta}((\tilde{\tilde{g}} - \eps_1))) - \tau(\chi_{\frac{1}{2}+\eps_1, \delta}(g))} < \eps/4 + \eps/4 = \eps/2,\quad \tau\in \tr(A) $$
By \eqref{small-strech}, one has
$$ \tau(\chi_{\frac{1}{2}+\eps_1, \delta}(g)) <\eps/2 + \eps/2 = \eps.$$
Stretching $g$ to move $\frac{1}{2} + \eps_1$ to $\frac{1}{2}$ (and note that $\eps_1< \eps/2$), it satisfies the desired approximations \eqref{positive-s1} and \eqref{positive-s2}.
% 
%\vskip 1in
%
%%%%%%%%%%%%% no induction %%%%%%%%%%%%%%%%%%%
%Define
%$$\left< \xi^-_{2, \eps''}(f) \right> := [u^*_1\eta_{\delta'}(\chi_{1, \eps''}(\tilde{\tilde{g}}))u_1] (\xi^-_{2, \eps''}(f)) \in \mathrm{Her}([u_1^* \eta_{\delta'}(\chi_{1, \eps''}(\tilde{\tilde{g}}))u_1] ) \cap D.$$
%It follows from \eqref{left-nbhd-4} that
%\begin{equation}
% \norm{\left< \xi^-_{2, \eps''}(f) \right> - \xi^-_{2, \eps''}(f)}_2 < 3\delta'.  
% \end{equation}
%%choose $$[\xi^-_{2, \eps''}(f)] = [u_1\eta_{\delta'}(\chi_{1, \eps''}(\tilde{\tilde{g}}))u_1] (\xi^-_{2, \eps''}(f)) \in \mathrm{Her}([\eta_{\delta'}(u_1^*\chi_{1, \eps''}(\tilde{\tilde{g}})u_1)] ),$$ 
%Then
%$$ \tau(\theta_{2, \eps''}(f)) > d_\tau(\xi^-_{2, \eps''}(f)) + 2 \eps' \geq d_\tau([u^*_1\eta_{\delta'}(\chi_{1, \eps''}(\tilde{\tilde{g}}))u_1]\xi^-_{2, \eps''}(f)) + 2\eps'  = \mathrm{d}_\tau(\left< \xi^-_{2, \eps''}(f) \right>) + 2\eps',$$
%and hence 
%$$\tau(\theta_{2, \eps''}(\tilde{\tilde{g}})) \approx_{\eps'} \tau(\theta_{2, \eps''}(f)) > \mathrm{d}_\tau( \left< \xi_{2, \eps''}^-(f) \right>) + 2 \eps'. $$ 
%In particular,
%\begin{equation}\label{lem-cond-2}
%\tau(\theta_{2, \eps''}(\tilde{\tilde{g}})) >  \mathrm{d}_\tau( \left< \xi_{2, \eps''}^-(f) \right> ),\quad \tau \in \tr(A). 
%\end{equation}
%
%Also define $$\left< \xi^+_{2, \eps''}(f) \right> := [u_1^*(\eta_{\delta'}(\chi_{1, \eps''}(\tilde{\tilde{g}}))u_1] (\xi^+_{2, \eps''}(f)) \in \mathrm{Her}([u_1^*(\eta_{\delta'}((\chi_{1, \eps''}(\tilde{\tilde{g}})))u_1] ) \cap D.$$
%
%Then (note that $\delta' < \eps'/2$), using \eqref{left-nbhd-3}
%\begin{eqnarray*}
%&& \tau(\eta_{\eps_0}(\chi_{2, \eps''}( f )) ) \\
%& < & \tau(\xi_{2, \eps''}^+(f)) - 2\eps' \\
%& \approx_{\delta'} & \tau( [\eta_{\delta'}(u_1^*\chi_{1, \eps''}(\tilde{\tilde{g}})u_1)] (\xi^+_{2, \eps''}(f)) ) - 2\eps' \\
%&= & \tau(  \left< \xi^+_{2, \eps''}(f) \right> ) - \eps' \\
%& < & \mathrm{d}_\tau( \left< \xi^+_{2, \eps''}(f) \right> ) - 2\eps', 
%\end{eqnarray*}
%and therefore
%$$ \tau(\eta_{\eps_0}(\chi_{2, \eps''}(\tilde{\tilde{g}}))) \approx_{\eps'} \tau(\eta_{\eps_0}(\chi_{2, \eps''}( f )) )  < d_\tau( [\xi^+_{2, \eps''}(f)]) - 2\eps' + \delta'.$$
%In particular, 
%\begin{equation}\label{lem-cond-3}
%\tau(\eta_{\eps_0}(\chi_{2, \eps''}(\tilde{\tilde{g}}))) < d_\tau( [\xi^+_{2, \eps''}(f)]),\quad \tau \in\tr(A).
%\end{equation}
%
%With \eqref{lem-cond-1}, \eqref{lem-cond-2}, and \eqref{lem-cond-3}, by Lemma \ref{D-sandwich}, for any $\delta''>0$ (which will be fixed later), there is $u_2 \in A$ such that
%\begin{equation}\label{u-2-cond-1}
% u_2^*\chi_{2, \eps''}(\tilde{\tilde{g}})u_2 \in^{\norm{\cdot}_2}_{\delta''}   \overline{[u_1^*\eta_{\delta'}(\chi_{1, \eps''}(\tilde{\tilde{g}}))u_1] A [u_1^*\eta_{\delta'}(\chi_{1, \eps''}(\tilde{\tilde{g}}))u_1] }
% \end{equation}
%
%\begin{equation}\label{u-2-cond-2}
%\eta_{\eps_0/2}(u^*_2 \chi_{2, \eps''}(\tilde{\tilde{g}}) u_2) \in^{\norm{\cdot}_2}_{\delta''} \overline{\left< \xi_{2, \eps''}^+(f) \right> A \left< \xi_{2, \eps''}^+(f) \right>},
%\end{equation}
%
%\begin{equation}\label{u-2-cond-3}
% \eta_{\delta''}(u_2^*\chi_{2, \eps''}(\tilde{\tilde{g}})u_2) \left< \xi_{2, \eps''}^-(f) \right>  \approx_{\delta''}^{\norm{\cdot}_2} \left< \xi_{2, \eps''}^-(f) \right>,
% \end{equation}
%
%%and
%\begin{equation}\label{u-2-cond-4}
%\mathrm{dist}_{2}(u_2du_2^*, D_1) < \delta'',\quad \mathrm{dist}_{2}(u_2^*du_2, D_1) < \delta'',\quad d \in D_1,
%\end{equation}
%and
%\begin{equation}\label{u-2-cond-5}
%\norm{u_2u_2^* - 1}_{2}, \norm{u_2^*u_2 - 1}_{2} < \delta''.
%\end{equation}
%
%\begin{eqnarray*}
% & & u_2^*( \chi_{2, \eps''}(\tilde{\tilde{g}})) u_2 \\
% & \approx_{\delta''}^{\norm{\cdot}_2} & (u_2^*( \chi_{2, \eps''}(\tilde{\tilde{g}})) u_2) ( \eta_{2\delta'}([u_1^* \chi_{1, \eps''}(\tilde{\tilde{g}}) u_1] )) \quad \quad (\eqref{u-2-cond-1}, \eqref{down-to-D-3}) \\
% & \approx_{\delta'}^{\norm{\cdot}_2} & (u_2^*( \chi_{2, \eps''}(\tilde{\tilde{g}})) u_2) ( \eta_{2\delta'}(u_1^* \chi_{1, \eps''}(\tilde{\tilde{g}}) u_1 )) \quad \quad (\eqref{down-to-D-5}) \\
% & \approx_{2\delta'}^{\norm{\cdot}_2} & (u_2^*( \chi_{2, \eps''}(\tilde{\tilde{g}})) u_2) ( \eta_{2\delta'}(u_1^* \chi_{1, \eps''}(\tilde{\tilde{g}}) u_1 )) (u_1\chi_{1, \eps''}(\tilde{\tilde{g}}) u_1) \\
% &\approx_{\delta'+\delta''}& (u_2^*( \chi_{2, \eps''}(\tilde{\tilde{g}})) u_2) (u^*_1\chi_{1, \eps''}(\tilde{\tilde{g}}) u_1) \quad \quad (\eqref{u-2-cond-1}, \eqref{down-to-D-3}), \eqref{down-to-D-5})
%\end{eqnarray*}
%%By \eqref{u-2-cond-1} and \eqref{down-to-D-3},
%%\begin{eqnarray*}
%%& & (u_1^* \chi_{1, \eps''}(\tilde{\tilde{g}}) u_1)(u_2^*(\chi_{2, \eps''}(\tilde{\tilde{g}}) )u_2) \\
%% & \approx_{\delta'}^{\norm{\cdot}_2} & [u_1^* \chi_{1, \eps''}(\tilde{\tilde{g}}) u_1)](u_2^*(\chi_{2, \eps''}(\tilde{\tilde{g}}) )u_2) \\
%%& = & [u_1^* \chi_{1, \eps''}(\tilde{\tilde{g}}) u_1] \eta_{2\delta'}([u_1^* \chi_{1, \eps''}(\tilde{\tilde{g}}) u_1] )(u_2^*(\chi_{2, \eps''}(\tilde{\tilde{g}}) )u_2) \\
%%& \approx_{2\delta'} & \eta_{2\delta'}([u_1^* \chi_{1, \eps''}(\tilde{\tilde{g}}) u_1] )(u_2^*(\chi_{2, \eps''}(\tilde{\tilde{g}}) )u_2) \\
%%& = & (u_2^*(\chi_{2, \eps''}(\tilde{\tilde{g}}) )u_2).
%%\end{eqnarray*}
%In other words,
%\begin{equation}
% (u_1^* \chi_{1, \eps''}(\tilde{\tilde{g}}) u_1)(u_2^*(\chi_{2, \eps''}(\tilde{\tilde{g}}) )u_2) \approx_{6\delta'}^{\norm{\cdot}_2} u_2^*(\chi_{2, \eps''}(\tilde{\tilde{g}}) )u_2 .
% \end{equation}
% 
%By \eqref{u-1-cond-2} and \eqref{u-2-cond-1}, \eqref{down-to-D-3}, \eqref{down-to-D-4},
%\begin{eqnarray*}
%\chi_1(f) (u_2^*(\chi_{2, \eps''}(\tilde{\tilde{g}}) )u_2) & \approx_{\delta''} & \chi_1(f) \eta_{2\delta'}([u_1^*\chi_{1, \eps''}(\tilde{\tilde{g}})u_1]) (u_2^*(\chi_{2, \eps''}(\tilde{\tilde{g}}) )u_2) \quad \quad (\eqref{down-to-D-3}) \\
%& = & \chi_1(f) \eta_{\eps_0/2}([u_1^*\chi_{1, \eps''}(\tilde{\tilde{g}})u_1]) (u_2^*(\chi_{2, \eps''}(\tilde{\tilde{g}}) )u_2)\\
%& \approx^{\norm{\cdot}_2}_{\delta'} & (\chi_1(f) \eta_{\eps_0/2}(u_1^*\chi_{1, \eps''}(\tilde{\tilde{g}})u_1)) (u_2^*(\chi_{2, \eps''}(\tilde{\tilde{g}}) )u_2) \quad \quad (\eqref{down-to-D-4}) \\
%& \approx^{\norm{\cdot}_2}_{\delta'} & \eta_{\eps_0/2}(u_1^*\chi_{1, \eps''}(\tilde{\tilde{g}})u_1) (u_2^*(\chi_{2, \eps''}(\tilde{\tilde{g}}) )u_2) \quad \quad (\eqref{u-1-cond-2}) \\
%& \approx^{\norm{\cdot}_2}_{\delta'} & \eta_{\eps_0/2}([u_1^*\chi_{1, \eps''}(\tilde{\tilde{g}})u_1]) (u_2^*(\chi_{2, \eps''}(\tilde{\tilde{g}}) )u_2) \quad \quad (\eqref{down-to-D-4}) \\
%& = & u_2^*(\chi_{2, \eps''}(\tilde{\tilde{g}}) )u_2.
%\end{eqnarray*}
%So,
%$$ \chi_1(f) (u_2^*(\chi_{2, \eps''}(\tilde{\tilde{g}}) )u_2) \approx^{\norm{\cdot}_2}_{4\delta'} u_2^*(\chi_{2, \eps''}(\tilde{\tilde{g}}) )u_2.  $$
%
%One also has (with the same argument as \eqref{id-down-1}), by \eqref{u-2-cond-3}, 
%\begin{equation}\label{id-down-2} 
% (u_2^*(\chi_{2, \eps''}(\tilde{\tilde{g}}) )u_2) \chi_3(f) \approx_{7\eps''}^{\norm{\cdot}_2} \chi_3(f).
%\end{equation}
%
%%%%%%%. no induction %%%%%%%%%%%%%%%%
%
%Repeating this process, one obtains $u_1, u_2, ..., u_N \in A$ such that
%
%$$ \norm{ \chi_i(f) (u_{i+1}^* \chi_{i+1, \eps''}(\tilde{\tilde{g}}) u^{i+1}) - (u_{i+1}^* \chi_{i+1, \eps''}(\tilde{\tilde{g}}) u^{i+1}) }_2 < 4 \delta' < \eps'',\quad i=1, 2, ..., N-1,$$
%
%$$ \norm{ (u_{i}^* \chi_{i, \eps''}(\tilde{\tilde{g}}) u^{i}) \chi_{i+1}(f) - \chi_{i+1}(f) }_2 < 7 \eps'',\quad i=1, 2, ..., N-1,$$
%
%and
%
%$$ \norm{ (u_{i}^* \chi_{i, \eps''}(\tilde{\tilde{g}}) u^{i})(u_{i+1}^* \chi_{i+1, \eps''}(\tilde{\tilde{g}}) u^{i+1}) - (u_{i+1}^* \chi_{i+1, \eps''}(\tilde{\tilde{g}}) u^{i+1}) }_2 < 6\delta',\quad i=1, 2, ..., N-1.$$
%
%Define
%$$g:= \frac{1}{N}(u_1^*\chi_{1, \eps''} u_1 + \cdots + u_N^*\chi_{N, \eps''} u_N).$$
%It then follows from Lemma \ref{hereditary-sets} that
%$$ \norm{ f - g}_2 < (2N+1) 7 \eps'' + 6/N $$
%and
%$$ \abs{ \tau(\chi_{\frac{1}{2}+\eps'', \delta}((\tilde{\tilde{g}} - \eps''))) - \tau(\chi_{\frac{1}{2}+\eps'', \delta}(g))} < \eps,\quad \tau\in \tr(A). $$
%Hence,
%$$ \tau(\chi_{\frac{1}{2}+\eps'', \delta}(g)) <2\eps + \eps.$$
%Stretching $g$ to move $\frac{1}{2} + \eps''$ to $\frac{1}{2}$, one has the desired approximation.
%
%\vskip 1in 
%
%Moreover,
%$$\mathrm{d}_\tau(\phi_2) < \tau(\theta^{\eps''}_1(f)) \approx_{2\eps'} \tau(\theta^{\eps''}_1(\tilde{\tilde{g}})) < \mathrm{d}_\tau(\psi_1),\quad \tau \in \tr(A),$$
%where 
%$$ \theta^{\eps''}_1 = \left\{ 
%\begin{array}{ll}
%0, & t \leq \frac{1}{N} - \eps'',\\
%\mathrm{linear}, & t \in [\frac{1}{N} - \eps'', \frac{1}{N}],\\
%1, & t \geq \frac{1}{N}.
%\end{array}
%\right.$$
%
%By the choice of $\eps'$, one has
%\begin{equation}\label{int-eq-1-1}
%\mathrm{d}_\tau(\phi_2) < \tau(\theta^{\eps''}_1(\tilde{\tilde{g}})) < \mathrm{d}_\tau(\psi_1), \quad \tau \in \tr(A). 
%\end{equation}
%By Theorem \ref{sandwich-global}, there is a unitary $v_1 \in \mathrm{M}_n(\mathrm{C}(X))$  such that
%\begin{equation}
%\mathrm{dist}_{2, \tr(A)}(v_1\psi_1 v^*_1, \overline{\phi_1 A\phi_1}) < \delta',
%\end{equation}
%\begin{equation}\label{int-eq-1-2}
%\norm{(v_1\psi_1v^*_1)\phi_2 - \phi_2}_{2, \tr(A)} < 2\eps'',
%\end{equation}
%%$$ (v_1\psi_1v^*_1)\phi_2 \approx_{\eps''} \phi_2 $$
%and
%\begin{equation}\label{int-eq-1-3}
%\mathrm{dist}_{2, \tr(A)}(v_1dv^*_1, D_1) < \delta'',\quad d \in (D)_1,
%%v_1(x) \in P_n(\frac{2}{n}),\quad x \in X.
%\end{equation}
%where $\delta''>0$ is a constant such that if $a, b$ are positive contractions such that
%$$\norm{a - b}_{2, \tr(A)} < \delta''$$
%then
%$$ \norm{\eta_{\delta'}(a) - \eta_{\delta'}(b) }_{2, \tr(A)} < \eps''.$$
%
%Consider $\eta_{\delta'}(\psi_1)$, where
%$$
%\eta_{\delta'} = \left\{
%\begin{array}{ll}
%0, & t < 1-\delta' \\
%\mathrm{linear} & t \in [1-\delta', 1] \\
%1 & t \geq 1.
%\end{array}
%\right.
%$$
%Then
%$$(v_1\eta_{\delta'}(\psi_1)v_1^*) \phi_3 \approx_\eps \phi_3$$
%
%Indeed, by \eqref{int-eq-1-3}, there is $d_1 \in D_1$ such that
%$$ \norm{d_1 - v_1\psi_1v^*_1}_{2, \tr(A)} < \delta''.$$
%By \eqref{int-eq-1-2}, 
%$$\norm{d_1\phi_2 - \phi_2}_{2, \tr(A)} < 2\eps'' + \delta'',$$
%
%
%
%
%
%
%
%
%
%
%\vskip 1in
%Note that then
%$$(v_1\eta_{\delta'}(\psi_1) v_1^*) \phi_3 = (v_1\eta_{\delta'}(\psi_1) v_1^*) \eta_{\delta''}(\phi_2) \phi_3    \approx \phi_3   $$
\end{proof}


%\begin{thm}
%Let $C = \overline{c\mathrm{M}_n(\mathrm{C}(X))c}$, where $X$ is a compact Hausdorff space, and $c$ is a positive contraction. Let $f, g \in C \cap D$ be two diagonal positive contractions such that 
%$$\mathrm{rank}(f(x)) < \mathrm{rank}(g(x)).$$ Then, for any $\eps>0$, there is a unitary $u \in C+ 1_n\Comp$ such that 
%$$u^*f u \in_\eps \overline{gCg}$$
%and
%$$\mathrm{dist}_{2}((u^*hu)(x), D_n) < \frac{2}{n},\quad x \in X,\ h\in (D)_1.$$
%\end{thm}
%
%\begin{cor}
%Let $A = \mathrm{M}_n(\mathrm{C}(X))$, and let $D$ be the diagonal subalgebra. Let $N \in \mathbb N$ and $\eps>0$.
%
%Assume $f, g, h\in D$ are positive contractions such that $$gh = h,$$ 
%$$ \mathrm{rank}(g(x)) > \mathrm{rank}(f(x)) > \mathrm{Tr}(\chi_{1-\delta}(f)(x)) > \mathrm{rank}(h(x)),\quad x\in X,\ \delta>0.$$
%%and
%%$$\mathrm{tr}(\chi_{1-\delta}(f)(x)) > \mathrm{rank}(h(x)),\quad x \in X.$$ 
%Then there is a positive contraction $a \in \overline{gAg}$ such that
%$$\mathrm{dist}_2(a, D) < \frac{2}{n},$$
%$$ \norm{ah - h}_2 < \eps $$
%and
%$$ \abs{\mathrm{tr}(a^n(x)) - \mathrm{tr}(f^n(x))} < \eps,\quad i=1, 2, ..., N,\ x\in X. $$
%\end{cor}
%
%\begin{cor}
%
%Let $\eps>0$, and let $\Delta: \mathrm{C}^+([0, 1]) \to (0, +\infty)$ be a density function.  Then there is a finite set $\mathcal H \subseteq \mathrm{C}^+([0, 1])$ and $\delta>0$ such that if $f_0, f_1 \in D \subseteq \mathrm{M}_n(\mathrm{C}(X))$ are two positive contractions such that
%$$\abs{\tau(h(\phi_{\frac{\eps}{8}}(f_0))) - \tau(h(f_1))} < \delta,\quad h \in \mathcal H,$$
%and
%$$\tau(h(f_0)) > \Delta(h),\quad h \in \mathcal H,$$
%and let $P$ is a polynomial, 
%then there is a positive contraction $ g \in A$ such that 
%$$\norm{f_0 - g } < \eps$$
%$$\abs{\tau(P(g)) - \tau(P(f_1))} < \eps,$$
%and
%$$ \mathrm{dist}(g, D_n) < \frac{1}{n}.$$
%\end{cor}
%
%\begin{proof}
%
%
%Consider the testing functions $\chi_i$, $i=0, 1, ..., N$, where 
%$$\chi_i(t) = \left\{
%\begin{array}{ll}
%0, & t\in [0, \frac{i}{N}], \\
%Nt - i, & t \in [\frac{i}{N}, \frac{i+1}{N}], \\
%1, & t \in  [\frac{i+1}{N}, 1],
%\end{array} 
%\right.$$
%and $\tilde{\chi}_i$, $i=1, ..., N$, where 
%$$\tilde{\chi}_i(t) = \left\{
%\begin{array}{ll}
%0, & t\in [0, \frac{2i-1}{2N}] , \\
%2Nt - i, & t \in[\frac{2i-1}{2N}, \frac{i}{N}],  \\
%1, & t \in  [\frac{i}{N}, 1],
%\end{array} 
%\right.$$
%
%
%Choose $$\delta = \frac{1}{8}\min\{\Delta(\tilde{\chi}_i - \chi_{i}),\ i=1, ..., N\}.$$
%
%Let $f_0, f_1 \in D$ be positive contractions which satisfy the assumptions with $\mathcal H = \{\chi_i, \tilde{\chi}_i\}$. 
%Consider 
%$$ a_0:=\chi_0(f_0), \  a_1:=\chi_1(f_0), \  ..., \  a_N:=\chi_N(f_0) $$
%and
%$$ b_1:=\chi_1(f_1), \  b_2:= \chi_1(f_1), \  ..., \  b_N:=\chi_N(f_1). $$
%
%Then, by Corollary \ref{}, there are unitaries $u_1, ..., u_N$ such that
%$$u^*_i b_i u_i \in \overline{a_{i-1} A a_{i-1}}$$
%$$ (u^*_i b_i u_i) a_i = a_i $$
%and
%$$\mathrm{dist}_2((u^*_i b_i u_i), D) < \frac{2}{n}.$$
%
%Then
%$$\norm{f_0 - \frac{1}{N}( u^*_1 b_1 u_1 + \cdots + u^*_N b_N u_N)} < \frac{8}{N}$$
%and
%$$\mathrm{dist}_2(\frac{1}{N}( u^*_1 b_1 u_1 + \cdots + u^*_N b_N u_N), D) < \frac{2}{n}.$$
%
%Let us compare 
%$$ P(\frac{1}{N}( u^*_1 b_1 u_1 + \cdots + u^*_N b_N u_N)) $$
%and
%$$ P(\frac{1}{N}(b_1 + \cdots + b_N)).$$
%
%Note that there are $c_{i, j}$, $i=1, ..., N$, $j = 1, ..., m = \mathrm{deg}(P)$, such that 
%$$P(\frac{1}{N}(b_1+ \cdots + b_N)) = (\sum_{n=1}^m c_{1, n} b_1^n) + \cdots + (\sum_{n=1}^m c_{N, n} b_N^n)$$
%and
%$$P(\frac{1}{N}(u_1^*b_1u_1+ \cdots + u_Nb_Nu_N^*)) \approx (\sum_{n=1}^m c_{1, n} (u_1b_1u_1^*)^n) + \cdots + (\sum_{n=1}^m c_{N, n} (u_Nb_Nu_N^*)^n).$$
%In particular,
%$$\tau( P(\frac{1}{N}(b_1+ \cdots + b_N)) ) = \tau(P(\frac{1}{N}(u_1^*b_1u_1+ \cdots + u_Nb_Nu_N^*))),\quad \tau \in \mathrm{T}(A).$$
%But
%$$f_1 = \frac{1}{N}(b_1+\cdots + b_N).$$
%(We really need this one to be an exact equality, not approximation!)
%
%
%
%\end{proof}



As a consequence of Theorem \ref{prop-S}, one has the following corollary: %characterizations of $\mathcal Z$-absorption of AH algebras with diagonal maps or the crossed product C*-algebras $\mathrm{C}(X) \rtimes \Gamma$:

\begin{cor}\label{Z-SBP}
Let $A$ be a simple AH algebra with diagonal maps, or let $A = \mathrm{C}(X)\rtimes\Gamma$, where $(X, \Gamma)$ is free, minimal, and has the (URP) and (COS). Then $(D, A)$ has Property (C) if, and only if, $(D, \tr(A)|_D)$ has the (SBP).
%
%Assume that $(D, A)$ has Property (C), where $D$ is the canonical commutative sub-C*-algebra. Then $(D, \tr(A)|_D)$ has the (SBP) if $A \cong A\otimes \mathcal Z$.
\end{cor}

\begin{proof}
Assume that $(D, A)$ has Property (C). By Theorem \ref{Property-D}, the C*-algebra pair  $(D, A)$ has Property (E). Note that $A$ has Property (S) (Theorem \ref{S-holds}).  By Theorem \ref{prop-S}, $(D, \tr(A))$ has the (SBP).

The other implication follows from Theorem \ref{SBP=>C}.
\end{proof}

\begin{rem}
Let $A$ be an AH algebra with diagonal map or a crossed product C*-algebra as above. It would be an interesting question whether $\mathcal Z$-absorption of $A$, which is a sole C*-algebra property, would imply Property (C) of the pair $(D, A)$.  
\end{rem}

%\begin{thm}\label{main-thm}
%Let $A$ be a simple AH algebra with diagonal maps, or let $A = \mathrm{C}(X)\rtimes\Gamma$, where $(X, \Gamma)$ is free, minimal, and has the (URP) and (COS). If $(D, A)$ has Property (C), where $D$ is the canonical commutative subalgebra which has Property (C). Then the following conditions are equivalent:
%\begin{enumerate}
%\item[(1)]\label{main-1} $A$ has Property (S). 
%\item[(2)]\label{main-2} $(D, \tr(A))$ has the (SBP).
%\item[(3)]\label{main-3} $A \cong A\otimes \mathcal Z. $
%\item[(4)]\label{main-4} The strict order on $\mathrm{Cu}(A)$ is determined by traces.
%\item[(5)]\label{main-5} $\mathrm{qRR}(l^\infty(A)/J_{2, \omega, \tr(A)}) = 0$ (Definition \ref{qRR-defn}).
%\item[(6)]\label{main-6} $\mathrm{RR}(l^\infty(A)/J_{2, \omega, \tr(A)}) = 0$.
%\item[(7)]\label{main-7} $\mathrm{RR}(l^\infty(D)/J_{2, \omega, \tr(A)}) = 0$.
%\item[(8)]\label{main-8} $A$ has uniform property $\Gamma$ (Definition \ref{UGamma-defn}).
%\item[(9)]\label{main-9} $(D, A)$ has strong uniform property $\Gamma$ (Definition \ref{AD-defn}).
%\item[(10)]\label{main-10} $(D, \tr(A))$ is approximately divisible (Definition \ref{AD-defn}).
%\end{enumerate}
%In the case that $A = \mathrm{C}(X)\rtimes\Gamma$, each statement above is also equivalent to
%\begin{enumerate}
%\item[(11)]\label{main-11} $\mathrm{mdim}(X, \Gamma) = 0$.
%\end{enumerate}
%\end{thm}
%
%\begin{proof}
%$(1) \Rightarrow (2)$: By Theorem \ref{property-C-thm} and Theorem \ref{Property-D} respectively, the C*-algebra pair  $(D, A)$ has Properties (C) and (E). Since $A$ has Property (S), by Theorem \ref{prop-S}, $(D, \tr(A))$ has the (SBP).
%
%$(2) \Rightarrow (3)$: For the crossed product C*-algebras, this follows from Theorem 4.7 of \cite{EN-MD0} in the case $\Gamma = \Int$ and follows in general from Theorem 5.4 of \cite{Lindenstrauss-Weiss-MD} and Theorem 4.8 of \cite{Niu-MD-Z-absorbing}. For the AH algebras with diagonal maps, this implication follows from Proposition 4.10 of \cite{EN-SBP-I}.
%
%$(3) \Rightarrow (8)$: Theorem 5.6 of \cite{CETW-Gamma}.
%
%$(8) \Rightarrow (1)$: Proposition \ref{Property-S-Prop}.
%
%
%$(2) \Leftrightarrow (7)$: Theorem 2.12 of \cite{EN-SBP-I}.
%
%$(10) \Rightarrow (2)$: Theorem 3.5 of \cite{EN-SBP-I}.
%
%$(1) \Leftrightarrow (5) \Leftrightarrow (6) $:  Proposition \ref{character-S}.
%
%$(2) \Rightarrow (9)$: In the case of $\mathrm{C}(X)\rtimes\Gamma$, this follows from the proof of Theorem 9.4 of \cite{KS-comparison}. For AH algebras with diagonal maps, this follows from Theorem \ref{AH-SUPG} of the appendix.
%
%$(9) \Rightarrow (10)$: Trivial.
%
%$(3) \Leftrightarrow (4)$: In the case of $\mathrm{C}(X)\rtimes\Gamma$, this follows from Corollary 7.14 of \cite{LN-sr1}. In the case of AH algebras with diagonal maps, this follows from Theorem 4.1 of \cite{EHT-sr1} and Theorem 9.5 of \cite{Thiel-sr1}
%
%This shows the equivalence of Conditions (1)--(10).
%
%$(11) \Rightarrow (2)$: Theorem 5.1 of \cite{Niu-MD-Z-absorbing} (\cite{MR3614036} and \cite{GLT-Zk} for $\Int^d$-actions).
%
%$(2) \Rightarrow (11)$: Theorem 5.4 of \cite{Lindenstrauss-Weiss-MD}.
%\end{proof}
%

\bibliographystyle{plainurl}
\bibliography{operator_algebras}

%\begin{thebibliography}{10}
%
%\bibitem{CETW-Gamma}
%J.~Castillejos, S.~Evington, A.~Tikuisis, and S.~White.
%\newblock Uniform property {$\Gamma$}.
%\newblock {\em Int. Math. Res. Not. IMRN}, (13):9864--9908, 2022.
%
%\bibitem{CETWW-dim-n}
%J.~Castillejos, S.~Evington, A.~Tikuisis, S.~White, and W.~Winter.
%\newblock Nuclear dimension of simple {C*}-algebras.
%\newblock {\em Invent. Math.}, 224(1):245--290, 2021.
%
%\bibitem{CE-str1}
%A.~Ciuperca and G.~A. Elliott.
%\newblock A remark on invariants for {C*}-algebras of stable rank one.
%\newblock {\em Int. Math. Res. Not. IMRN}, (5):Art. ID rnm 158, 33, 2008.
%
%\bibitem{EGLN-ASH}
%G.~A. Elliott, G.~Gong, H.~Lin, and Z.~Niu.
%\newblock The classification of simple separable unital $\textrm{$\mathcal
%  Z$}$-stable locally {ASH} algebras.
%\newblock {\em J. Funct. Anal.}, (12):5307--5359, 2017.
%
%\bibitem{EGLN-DR}
%G.~A. Elliott, G.~Gong, H.~Lin, and Z.~Niu.
%\newblock On the classification of simple amenable {C*}-algebras with finite
%  decomposition rank, {II}.
%\newblock {\em J. Noncommut. Geom.}, 19(1):73--104, 2025.
%
%\bibitem{EHT-sr1}
%G.~A. Elliott, T.~M. Ho, and A.~S. Toms.
%\newblock A class of simple {C*}-algebras with stable rank one.
%\newblock {\em J. Funct. Anal.}, 256(2):307--322, 2009.
%
%\bibitem{ELN-Vill}
%G.~A. Elliott, C.~G. Li, and Z.~Niu.
%\newblock Remarks on {V}illadsen algebras.
%\newblock {\em J. Funct. Anal.}, 287(7):Paper No. 110547, 55, 2024.
%
%\bibitem{EN-K0-Z}
%G.~A. Elliott and Z.~Niu.
%\newblock On the classification of simple amenable {C*}-algebras with finite
%  decomposition rank.
%\newblock In R.~S. Doran and E.~Park, editors, {\em ``Operator Algebras and
%  their Applications: A Tribute to Richard V.~Kadison", Contemporary
%  Mathematics}, volume 671, pages 117--125. Amer. Math. Soc., 2016.
%
%\bibitem{EN-MD0}
%G.~A. Elliott and Z.~Niu.
%\newblock The {C{$^*$}}-algebra of a minimal homeomorphism of zero mean
%  dimension.
%\newblock {\em Duke Math. J.}, 166(18):3569--3594, 2017.
%
%\bibitem{EN-SBP-I}
%G.~A. Elliott and Z.~Niu.
%\newblock On the small boundary property, $\mathcal {Z}$-absorption, and
%  {B}auer simplexes.
%\newblock {\em arXiv:2406.09748}, 2024.
%
%\bibitem{GK-Dyn}
%J.~Giol and D.~Kerr.
%\newblock Subshifts and perforation.
%\newblock {\em J. Reine Angew. Math.}, 639:107--119, 2010.
%
%\bibitem{GLN-TAS-1}
%G.~Gong, H.~Lin, and Z.~Niu.
%\newblock Classification of finite simple amenable {$\mathcal Z$}-stable
%  {C*}-algebras, {I}. {C*}-algebras with generalized tracial rank one.
%\newblock {\em C. R. Math. Acad. Sci. Soc. R. Can.}, 42(3):63--450, 2020.
%
%\bibitem{GLN-TAS-2}
%G.~Gong, H.~Lin, and Z.~Niu.
%\newblock Classification of finite simple amenable {$\mathcal Z$}-stable
%  {C*}-algebras, {II}. {C*}-algebras with rational generalized tracial rank
%  one.
%\newblock {\em C. R. Math. Acad. Sci. Soc. R. Can.}, 42(4):451--539, 2020.
%
%\bibitem{Gromov-MD}
%M.~Gromov.
%\newblock Topological invariants of dynamical systems and spaces of holomorphic
%  maps. {I}.
%\newblock {\em Math. Phys. Anal. Geom.}, 2(4):323--415, 1999.
%
%\bibitem{MR3614036}
%Y.~Gutman.
%\newblock Embedding topological dynamical systems with periodic points in
%  cubical shifts.
%\newblock {\em Ergodic Theory Dynam. Systems}, 37(2):512--538, 2017.
%
%\bibitem{GLT-Zk}
%Y.~Gutman, E.~Lindenstrauss, and M.~Tsukamoto.
%\newblock Mean dimension of {$\Bbb{Z}^k$}-actions.
%\newblock {\em Geom. Funct. Anal.}, 26(3):778--817, 2016.
%
%\bibitem{JS-Z}
%X.~Jiang and H.~Su.
%\newblock On a simple unital projectionless {C*}-algebra.
%\newblock {\em Amer. J. Math.}, 121(2):359--413, 1999.
%
%\bibitem{KS-comparison}
%D.~Kerr and G.~Szab\'{o}.
%\newblock Almost finiteness and the small boundary property.
%\newblock {\em Comm. Math. Phys.}, 374(1):1--31, 2020.
%
%\bibitem{KR-CenSeq}
%E.~Kirchberg and M.~R{\o}rdam.
%\newblock Central sequence {C*}-algebras and tensorial absorption of the
%  {J}iang-{S}u algebra.
%\newblock {\em J. Reine Angew. Math.}, 695:175--214, 2014.
%
%\bibitem{LN-sr1}
%C.~G. Li and Z.~Niu.
%\newblock Stable rank of {$\mathrm{C}(X)\rtimes\Gamma$}.
%\newblock {\em arXiv: 2008.03361}, 2020.
%
%\bibitem{Li-interval}
%L.~Li.
%\newblock Simple inductive limit {C*}-algebras: spectra and approximations by
%  interval algebras.
%\newblock {\em J. Reine Angew. Math.}, 507:57--79, 1999.
%
%\bibitem{Lind-MD}
%E.~Lindenstrauss.
%\newblock Mean dimension, small entropy factors and an embedding theorem.
%\newblock {\em Inst. Hautes {\'E}tudes Sci. Publ. Math.}, (89):227--262 (2000),
%  1999.
%
%\bibitem{Lindenstrauss-Weiss-MD}
%E.~Lindenstrauss and B.~Weiss.
%\newblock Mean topological dimension.
%\newblock {\em Israel J. Math.}, 115:1--24, 2000.
%
%\bibitem{Matui-Sato-CP}
%H.~Matui and Y.~Sato.
%\newblock Strict comparison and {$\mathcal Z$}-absorption of nuclear
%  {C*}-algebras.
%\newblock {\em Acta Math.}, 209(1):179--196, 2012.
%
%\bibitem{Niu-MD}
%Z.~Niu.
%\newblock Mean dimension and {AH}-algebras with diagonal maps.
%\newblock {\em J. Funct. Anal.}, 266(8):4938--4994, 2014.
%
%\bibitem{Niu-MD-Z-absorbing}
%Z.~Niu.
%\newblock {$\mathcal{Z}$}-stability of transformation group {C*}-algebras.
%\newblock {\em Trans. Amer. Math. Soc.}, 374(10):7525--7551, 2021.
%
%\bibitem{Niu-MD-Z}
%Z.~Niu.
%\newblock Comparison radius and mean topological dimension: {R}okhlin property,
%  comparison of open sets, and subhomogeneous {C*}-algebras.
%\newblock {\em J. Analyse Math.}, 146:595--672, 2022.
%
%\bibitem{Niu-MD-Zd}
%Z.~Niu.
%\newblock Comparison radius and mean topological dimension: $\mathbb
%  {Z}^d$-actions.
%\newblock {\em Canad. J. Math.}, 76(4):1240--1266, 2024.
%
%\bibitem{RorUHF2}
%M.~R{\o}rdam.
%\newblock On the structure of simple {C*}-algebras tensored with a
%  {UHF}-algebra. {II}.
%\newblock {\em J. Funct. Anal.}, 107(2):255--269, 1992.
%
%\bibitem{Sato-CP}
%Y.~Sato.
%\newblock Trace spaces of simple nuclear {C*}-algebras with finite-dimensional
%  extreme boundary.
%\newblock {\em arXiv: 1209.3000}, 2012.
%
%\bibitem{Thiel-sr1}
%H.~Thiel.
%\newblock Ranks of operators in simple {C*}-algebras with stable rank one.
%\newblock {\em Comm. Math. Phys.}, 377:37--76, 2020.
%
%\bibitem{Thomsen-AI-Tr}
%K.~Thomsen.
%\newblock Inductive limits of interval algebras: the tracial state space.
%\newblock {\em Amer. J. Math.}, 116(3):605--620, 1994.
%
%\bibitem{TWW-QD}
%A.~Tikuisis, S.~White, and W.~Winter.
%\newblock Quasidiagonality of nuclear {C*}-algebras.
%\newblock {\em Ann. of Math. (2)}, 185(1):229--284, 2017.
%
%\bibitem{Toms-Ann}
%A.~S. Toms.
%\newblock On the classification problem for nuclear {C*}-algebras.
%\newblock {\em Ann. of Math. (2)}, 167(3):1029--1044, 2008.
%
%\bibitem{Toms-PLMS}
%A.~S. Toms.
%\newblock Stability in the {C}untz semigroup of a commutative {$
%  \mbox{C*}$}-algebra.
%\newblock {\em Proc. Lond. Math. Soc. (3)}, 96(1):1--25, 2008.
%
%\bibitem{Toms-Comp-DS}
%A.~S. Toms.
%\newblock Comparison theory and smooth minimal {C*}-dynamics.
%\newblock {\em Comm. Math. Phys.}, 289(2):401--433, 2009.
%
%\bibitem{TWW-Z}
%A.~S. Toms, S.~White, and W.~Winter.
%\newblock {$\mathcal{Z}$}-stability and finite-dimensional tracial boundaries.
%\newblock {\em Int. Math. Res. Not. IMRN}, (10):2702--2727, 2015.
%
%\bibitem{Vill-perf}
%J.~Villadsen.
%\newblock Simple {C*}-algebras with perforation.
%\newblock {\em J. Funct. Anal.}, 154(1):110--116, 1998.
%
%\bibitem{Vill-sr}
%J.~Villadsen.
%\newblock On the stable rank of simple {C*}-algebras.
%\newblock {\em J. Amer. Math. Soc.}, 12(4):1091--1102, 1999.
%
%\end{thebibliography}





%\appendix
%
%\section{Strong uniform property $\Gamma$ for AH algebras}
%
%\begin{thm}\label{AH-SUPG}
%Let $A$ be a simple unital AH algebra with diagonal maps, and let $D$ be the standard diagonal subalgebra. If $(D, \mathrm{T}(A))$ has the (SBP), then $(D, A)$ has strong uniform property $\Gamma$.
%\end{thm}
%
%\begin{proof}
%Let $\mathcal F \subseteq A$, $\eps>0$ and $n \in \mathbb N$, and let us construct mutually orthogonal  positive contractions $p_1, ..., p_n \in D$ such that 
%\begin{equation}\label{AH-SUPG-cond-1}
% \norm{1 - (p_1 + \cdots + p_n)}_{2, \tr(A)} < \eps,
% \end{equation}
%\begin{equation}\label{AH-SUPG-cond-2}
% \norm{ p_i f - f p_i} < \eps,\quad f \in \mathcal F,\ i=1, ..., n, 
% \end{equation}
% and
% \begin{equation}\label{AH-SUPG-cond-3}
%\abs{\tau (p_i f p_i) - \frac{1}{n} \tau(f)} < \eps,\quad i=1, ..., n,\ \tau \in\tr(A).
%\end{equation}
%Then, strong uniform property $\Gamma$ follows.
%
%To simplify the notation, let us only prove the theorem for the case that $A_i = \mathrm{M}_{n_1}(\mathrm{C}(X_i))$. The general case follows from the same argument.
%
%Let us first construct a new AH decomposition of $A$ with diagonal maps such that each building block contains $D$ as its diagonal subalgebra.
%
%Note that the AH algebra $A$ is generated by the diagonal subalgebra $D$ and the standard AF subalgebra $F$. Start with $A_1 = \mathrm{M}_{n_1}(\mathrm{C}(X_1))$. Consider the rank-one diagonal projections $p^{(1)}_1, ..., p^{(1)}_{n_1} \in \mathrm{M}_{n_1}(\Comp) \subseteq  A$ and the partial isometry $$ v^{(1)} = \left( 
%\begin{array}{cccc}
%0 & 1 & & \\
%& \ddots & \ddots &\\
%&&0 & 1 \\
%& & & 0
%\end{array}
%\right) \in \mathrm{M}_{n_1}(\Comp)\subseteq A,$$
%and consider the C*-algebra
%$$C_1 := \mathrm{C^*}\{p^{(1)}_iDp^{(1)}_i, v^{(1)}: i=1, ..., n_1\} \subseteq A.$$
%Note that
%$$C_1 \cong \mathrm{M}_{n_1}(p^{(1)}_1Dp^{(1)}_1).$$
%Moreover, since $p^{(1)}_1, ..., p^{(1)}_{n_1} \subseteq D$ and $p^{(1)}_1+ \cdots + p^{(1)}_{n_1} = 1$, one has that $D \subseteq C_1$ as the diagonal subalgebra, and there is a canonical homomorphism $\Lambda_1^*: A_1 \to C_1$ induced by the inductive limit $D_1 \to D_2 \to \cdots \to D$, where $D_i$ is the diagonal subalgebra of $A_i$, $i=1, 2, ...$.
%
%Similarly, consider the rank-one diagonal projections $p^{(2)}_1, ..., p^{(2)}_{n_2} \in \mathrm{M}_{n_2}(\Comp) \subseteq A$ and the partial isometry $$ v^{(2)} = \left( 
%\begin{array}{cccc}
%0 & 1 & & \\
%& \ddots & \ddots &\\
%&&0 & 1 \\
%& & & 0
%\end{array}
%\right) \in \mathrm{M}_{n_2}(\Comp)\subseteq A,$$
%and consider the C*-algebra
%$$C_2 := \mathrm{C^*}\{p^{(2)}_iDp^{(2)}_i, v^{(2)}: i=1, ..., n_2\} \subseteq A.$$
%Note that $$C_2 \cong \mathrm{M}_{n_2}(p^{(2)}_1Dp^{(2)}_1).$$
%Since the partition of unity $\{p_1^{(2)}, ..., p_{n_2}^{(2)}\}$ refines $\{p^{(1)}_1, ..., p^{(1)}_{n_1}\}$,  one has  that $C_1 \subseteq C_2$, and the following diagram commutes:
%\begin{equation*}
%\xymatrix{
%\mathrm{M}_{n_1}(p^{(1)}_1Dp^{(1)}_1) \ar@{^{(}->}[r] & \mathrm{M}_{n_2}(p^{(2)}_1Dp^{(2)}_1) \ar@{^{(}->}[r] & A \ar@{=}[d]  \\
%\mathrm{M}_{n_1}(\mathrm{C}(X_1)) \ar[r] \ar^{\Lambda_1^*}[u] & \mathrm{M}_{n_2}(\mathrm{C}(X_2)) \ar[r] \ar[u]^{\Lambda_2^*} & A .
%}
%\end{equation*}
%
%
%Repeating this process, one has the following commutative diagram:
%\begin{equation*}
%\xymatrix{
%\mathrm{M}_{n_1}(p^{(1)}_1Dp^{(1)}_1) \ar@{^{(}->}[r] & \mathrm{M}_{n_2}(p^{(2)}_1Dp^{(2)}_1) \ar@{^{(}->}[r] & \cdots \ar@{^{(}->}[r] & A \ar@{=}[d]  \\
%\mathrm{M}_{n_1}(\mathrm{C}(X_1)) \ar[r] \ar^{\Lambda_1^*}[u] & \mathrm{M}_{n_2}(\mathrm{C}(X_2)) \ar[r] \ar[u]^{\Lambda_2^*} & \cdots \ar[r]  &  A ,
%}
%\end{equation*}
%and therefore, we have the following AH decomposition of $A$: 
%\begin{equation}
%\xymatrix{
%\mathrm{M}_{n_1}(p^{(1)}_1Dp^{(1)}_1) \ar@{^{(}->}[r] & \mathrm{M}_{n_2}(p^{(2)}_1Dp^{(2)}_1) \ar@{^{(}->}[r] & \cdots \ar@{^{(}->}[r] & A.
%}
%\end{equation} 
%
%For each $i=1, 2, ...$, also write $$p_1^{(i)}Dp_1^{(i)} = \mathrm{C}(\tilde{X}_i), $$ where $\tilde{X}_i$ is metrizable and compact. 
%
%Without loss of generality, one may assume that $\mathcal F \subseteq \mathrm{M}_{n_1}(\mathrm{C}(\tilde{X}_1))$.
%%
%%\vskip 1in 
%%
%%
%%Denote by $F$ the standard AF subalgebra of $A$ generated by constant matrix-valued functions of each building blocks. Then 
%%\begin{enumerate}
%%\item $A = \textrm{C*}(D, F)$,
%%\item $v^*fv \in D$, where $v$ is a standard partial isometry of $F$.
%%\end{enumerate}
%%Note that, since $A_i = \mathrm{M}_{n_i}(\mathrm{C}(X_i))$, $i=1, 2, ...$, one has $F \cong \bigotimes_{i=1}^\infty \mathrm{M}_{n_i}(\Comp)$. 
%%
%%Without loss of generality, one may assume that $\mathcal F = \{\phi_1, ..., \phi_k, p_1, ..., p_{n_1}, v\}$, where $\phi_1, ..., \phi_k \in D_1$, $p_1, ..., p_{n_1}$ are rank-one diagonal projections of $\mathrm{M}_{n_1}(\Comp)$ (note that $p_1, ..., p_{n_1}$ form a partition of unity of $A$), and  
%%$$ v = \left( 
%%\begin{array}{cccc}
%%0 & 1 & & \\
%%& \ddots & \ddots &\\
%%&&0 & 1 \\
%%& & & 0
%%\end{array}
%%\right) \in \mathrm{M}_{n_1}(\Comp)\subseteq A. $$
%%
%%Consider the C*-algebra
%%$$C_1 = \mathrm{C^*}\{p_iDp_i, v: i=1, ..., n_1\} \subseteq A.$$ Since $p_1, ..., p_{n_1} \subseteq D$ and $p_1+ \cdots + p_{n_1} = 1$, there is a canonical map $\Lambda_1^*: A_1 \subseteq C_1$ induced by the the inductive limit $D_1 \to D_2 \to \cdots \to D$. Note that
%%$$C_1 \cong \mathrm{M}_{n_1}(p_1Dp_1).$$
%%
%%Choose $\delta>0$ such that
%%$$\norm{\phi_i(x) - \phi_i(y)} < \eps,\quad \mathrm{dist}(x, y)<\delta,\ i=1, ..., k,$$ where $\phi_1, ..., \phi_k$ are regarded as diagonal matrices on $X_{n_1}$. Then choose a finite set $\{y_1, ..., y_N\} \subseteq X_1$ which is $\delta$-dense.
%%
%%Choose $i$ sufficiently large, say $i=2$, such that, by the simplicity of $A$, for any $x \in X_2$, there is a partition $$\{1, 2, ..., m\} = \alpha_0(x) \sqcup \alpha_1(x) \sqcup \cdots \sqcup \alpha_N(x)$$ such that
%%$$ \mathrm{dist}(\lambda_s(x), y_i) < \delta,\quad i=1, ..., N,\ s \in \alpha_i(x),$$
%%$$ \frac{\abs{\alpha_0(x)}}{\abs{\alpha_1(x)}+ \cdots + \abs{\alpha_N(x)}} < \eps, \quad \mathrm{and} \quad \abs{\alpha_i(x)} \in n\mathbb N,\quad i=1, ..., N,$$
%%where $\lambda_1, ..., \lambda_m: X_2 \to X_1$ are the eigenvalue maps. 
%%For each $\alpha_i$, $i=1, ..., N$, decompose it into
%%$$\alpha_i(x) = \alpha_{i, 1}(x) \sqcup \cdots \sqcup \alpha_{i, n}(x)$$
%%such that $\abs{\alpha_{i, j_1}(x)} = \abs{\alpha_{i, j_2}(x)} $, $j_1, j_2 = 1, ..., n$.
%%
%%%$$ n_2 = m n_1, \quad  m = ns + r, \quad \mathrm{and} \quad \frac{r}{m} < \eps. $$ 
%%Consider the rank-one diagonal projections $q_1, ..., q_{n_2} \subseteq \mathrm{M}_{n_2}(\Comp),$
%%$$ w = \left( 
%%\begin{array}{cccc}
%%0 & 1 & & \\
%%& \ddots & \ddots &\\
%%&&0 & 1 \\
%%& & & 0
%%\end{array}
%%\right) \in \mathrm{M}_{n_2}(\Comp)\subseteq A, $$
%%and consider
%%$$C_2 = \mathrm{C^*}\{q_iDq_i, w: i=1, ..., n_2\} \subseteq A.$$
%%Since the partition of unity $\{q_1, ..., q_{n_2}\}$ refines $\{p_1, ..., p_{n_1}\}$ and $\mathrm{M}_{n_1}(\Comp) \subseteq \mathrm{M}_{n_2}(\Comp) $, one has $C_1 \subseteq C_2$. Note that $$C_2 \cong \mathrm{M}_{n_2}(q_1Dq_1).$$
%%We indeed have the following commutative diagram:
%%\begin{equation*}
%%\xymatrix{
%%\mathrm{M}_{n_1}(p_1Dp_1) \ar@{^{(}->}[r] & \mathrm{M}_{n_2}(q_1Dq_1) \ar@{^{(}->}[r] & A \ar@{=}[d]  \\
%%\mathrm{M}_{n_1}(\mathrm{C}(X_1)) \ar[r] \ar^{\Lambda_1^*}[u] & \mathrm{M}_{n_2}(\mathrm{C}(X_2)) \ar[r] \ar[u]^{\Lambda_2^*} & A 
%%}
%%\end{equation*}
%%
%%Write $q_1Dq_1 = \mathrm{C}(\tilde{X}_{n_2})$, and denote by $\Lambda: \tilde{X}_{n_2} \to X_2$ the map induced by $\Lambda_2^*$. Define $$\tilde{\lambda}_1 = \lambda_1\circ \Lambda,\  ...,\  \tilde{\lambda}_m = \lambda_m \circ\Lambda: \tilde{X}_{n_2} \to X_{n_1}.$$ 
%%Then, for each $x \in \tilde{X}_{n_2}$, there is a partition $$\{1, 2, ..., m\} = \alpha_0(x) \sqcup \alpha_1(x) \sqcup \cdots \sqcup \alpha_N(x)$$ such that
%%$$ \mathrm{dist}(\tilde{\lambda}_s(x), y_i) < \delta,\quad i=1, ..., N,\ s \in \alpha_i(x),$$
%%$$ \frac{\abs{\alpha_0(x)}}{\abs{\alpha_1(x)}+ \cdots + \abs{\alpha_N(x)}} < \eps, \quad \mathrm{and} \quad \abs{\alpha_i(x)} \in n\mathbb N,\quad i=1, ..., N.$$
%%By compactness of $\tilde{X}_{n_2}$, there is an open cover $\{U_1, ..., U_M\}$ of $\tilde{X}_{n_2}$ and points $x_1 \in U_1$, ..., $x_M \in U_M$ such that for each $j=1, ..., M$, 
%%$$ \mathrm{dist}(\tilde{\lambda}_s(x), y_i) < \delta,\quad i=1, ..., N,\ s \in \alpha_i(x_j),\ x\in U_j. $$
%%
%%Since $(D, A)$ has the (SBP), for each $j=1, ..., M$, there is a closed neighbourhood $V_j \ni x_j$ such that $V_j \subseteq U_j$, $V_1, ..., V_M$ are mutually disjoint, and
%%$$ \mu_\tau(\tilde{X}_{n_2} \setminus \bigcup_{j = 1}^M V_j) < \eps/m,\quad  \tau \in \tr(A).$$ 
%%
%%Then, for each $V_j$, $j=1, ..., M$, consider the tower
%%$$ \underbrace{ \underbrace{V_j, ..., \sigma^{n_1-1}(V_{j})}_{n_1}, ...,   \underbrace{\sigma^{(m-1)n_1}(V_j), ..., \sigma^{mn_1-1}(V_j)}_{n_1} }_m. $$
%%Also choose mutually orthogonal functions $h_j: \tilde{X}_2 \to [0, 1]$, $j=1, ..., M$, such that
%%$$ h_i|_{V_j} = 1,\quad j=1, ..., M.$$
%%
%%Then, for each $i=1, ..., n$, and $j=1, ..., M$, define the following elements in $C_2$:
%%$$p_{i, j} = \sum_{t=1}^N\sum_{s \in \alpha_{t, i}(x_j)} (e_{s} \otimes 1_{n_1}) h_j \quad \mathrm{and} \quad p_i = \sum_{j=1}^M p_{i, j}.$$
%%Then $p_1, ..., p_n$ have the desired property.
%%
%%\vskip 1in 
%%
%%
%%To construct $p_1, ... p_n$, without loss of generality, one may assume that $\mathcal F \subseteq A_1$. 
%%$$\mathcal F = \{\phi_1, ..., \phi_s, v_1, ..., v_t\},$$ where $\phi_1, ..., \phi_s \in D_1$ and $v_1, ..., v_t \in F_1$, $s, t \in \mathbb N$. 
%%
%Choose a finite open cover $\mathcal U$ of $\tilde{X}_{n_1}$  such that 
%\begin{equation}\label{dense-in-X1}
%\norm{f(x) - f(y)} < \eps,\quad x, y \in U,\ U \in \mathcal U,\ f \in \mathcal F.
%\end{equation} 
%Since $(D, \tr(A))$ has the (SBP), by Lemma 2.4 of \cite{EN-SBP-I}, there is an  open set $\hat{U} \subseteq U$ for each $U \in \mathcal U$ such that the subsets $$\hat{U},\quad U \in \mathcal U,$$ are mutually disjoint, and 
%\begin{equation}\label{SBP-cover-1}
%\mu_\tau(\tilde{X}_{1} \setminus \bigcup_{U \in \mathcal U} \hat{U}) < \frac{\eps}{n_1},\quad \tau \in \tr(A).
%\end{equation}
%Choose $x_U \in \hat{U}$ for each $U \in \mathcal U$.
%
%Choose $m_0 \in \mathbb N$ such that if $m_0 = nk + r$, where $ 0\leq r < n $, then $r/m_0 < \eps$.
%
%Move to the next stage sufficiently far out, say $C_2$, such that 
%by the simplicity of $A$, 
%$$ \abs{\{i=1, ...,m: \lambda_i(x) \in \hat{U} \}} \geq m_0,\quad U \in \mathcal U,\ x\in \tilde{X}_{2},$$
%and by \eqref{SBP-cover-1}, 
%$$ \frac{1}{m} \abs{\{i=1, ...,m: \lambda_i(x) \in \bigcup_{U \in \mathcal U} \hat{U} \}} > 1-\eps, \quad x\in \tilde{X}_2, $$
%where $\lambda_1, ..., \lambda_m: \tilde{X}_2 \to \tilde{X}_1$ are the eigenvalue maps of $C_1 \to C_2$. 
%
%
%By the compactness of $\tilde{X}_{n_2}$,  there is  an open cover $\mathcal V$ of $ \tilde{X}_{2}$ such that for each $V \in \mathcal V$ and each $U \in \mathcal U$, there is $\alpha_{U, V} \subseteq \{1, 2, ..., m\}$ such that 
%\begin{equation}\label{small-image-V}
%\lambda_i(V) \subseteq \hat{U}, \quad i \in \alpha_{U, V},
%\end{equation}
%\begin{equation}\label{SBP-large-multiplicity}
%\frac{1}{m}\sum_{U \in \mathcal U} \abs{\alpha_{U, V}} > 1-\eps \quad \mathrm{and}\quad \abs{\alpha_{U, V}} > m_0,\quad U \in \mathcal U, \  V \in \mathcal V. 
%\end{equation} 
%
%Since $(D, \tr(A))$ has the (SBP), by Lemma 2.4 of \cite{EN-SBP-I}, there is an open set $\hat{V} \subseteq V$ for each $V \in \mathcal V$ such that the open sets $\hat{V}$, $V \in \mathcal V$, are mutually disjoint, and 
%\begin{equation}\label{SBP-cover-2}
%\mu_\tau(\tilde{X}_{2} \setminus \bigcup_{V \in \mathcal V} \hat{V}) < \frac{\eps}{n_2},\quad \tau \in \tr(A).
%\end{equation}
%Then there are continuous functions $h_V: \tilde{X}_2 \to [0, 1]$, $V \in \mathcal V$, such that 
%$ h_V|_{\hat{V}^c} = 0$ and
%\begin{equation}\label{dense-h}
%\norm{1 - \sum_{V \in \mathcal V} h_V}_{2, \tr(A)} < \eps. 
%\end{equation}
%
%Inside each $\alpha_{U, V}$, by the choice of $m_0$, there are mutually disjoint sets $$\beta^{(1)}_{U, V}, ..., \beta^{(n)}_{U, V} \subseteq \alpha_{U, V}$$ 
%such that 
%\begin{equation}\label{equal-beta}
% \abs{ \beta^{(1)}_{U, V} } = \cdots =  \abs{ \beta^{(n)}_{U, V} } 
% \end{equation}
%and 
%\begin{equation}\label{large-beta}
%\frac{1}{\abs{ \alpha_{U, V} }} ( \abs{ \beta^{(1)}_{U, V} } + \cdots +  \abs{ \beta^{(n)}_{U, V} } ) > 1-\eps.
%\end{equation} 
%Then, define
%$$ p_1 = \sum_{V \in \mathcal V}\sum_{U \in \mathcal U}\sum_{i \in \beta_{U, V}^{(1)}} (h_V) e_i, \quad ..., \quad   p_n = \sum_{V \in \mathcal V}\sum_{U \in \mathcal U}\sum_{i \in \beta_{U, V}^{(n)}} (h_V) e_i,$$
%where, for each $i=1, ..., m$, $e_i$ denotes the diagonal projection of $C_2$ corresponding to the eigenvalue map $\lambda_i$. Then the contractions $p_1, ..., p_n$ have the desired property.
%
%Indeed, it is clear that $p_1, ..., p_n \in D$ and are mutually orthogonal. It follows from  \eqref{large-beta}, \eqref{SBP-large-multiplicity}, and \eqref{dense-h} that 
%\begin{eqnarray*}
%\norm{1 - (p_1 + \cdots + p_n)}_{2, \tr(A)} & = & \norm{1- \sum_{V \in \mathcal V}\sum_{U \in \mathcal U}(\sum_{i \in \beta_{U, V}^{(1)}} (h_V) e_i + \cdots + \sum_{i \in \beta_{U, V}^{(n)}} (h_V) e_i )}_{2, \tr(A)} \\
%& = & \norm{1- \sum_{V \in \mathcal V}\sum_{U \in \mathcal U} (\sum_{j = 1}^n\sum_{i \in \beta_{U, V}^{(j)}} e_i) h_V}_{2, \tr(A)} \\
%& \approx_\eps & \norm{1- \sum_{V \in \mathcal V}\sum_{U \in \mathcal U} (\sum_{i \in \alpha_{U, V}} e_i) h_V}_{2, \tr(A)}  \\
%& \approx_\eps & \norm{1- \sum_{V \in \mathcal V} h_V}_{2, \tr(A)} \\
%& \approx_\eps & 0.
%\end{eqnarray*}
%This verifies \eqref{AH-SUPG-cond-1}.
%
%Let us verify  \eqref{AH-SUPG-cond-2}. For each $f \in \mathcal F$, its image in $C_2$ is $$\sum_{j=1}^m (f \circ \lambda_j) e_j.$$
%Therefore, for each $i=1, ..., n$, noting that $h_V$, $V \in \mathcal V$, are in the center of $C_2$, we have 
%\begin{eqnarray*}
%(\sum_{j=1}^m (f \circ \lambda_j) e_j) p_i & = & (\sum_{j=1}^m (f \circ \lambda_j) e_j) (\sum_{V \in \mathcal V}\sum_{U \in \mathcal U}\sum_{j \in \beta_{U, V}^{(i)}} (h_V) e_j) \\
%& = & \sum_{V \in \mathcal V}\sum_{U \in \mathcal U}\sum_{j \in \beta_{U, V}^{(i)}} (h_V) (f\circ\lambda_j) e_j \\
%& = & (\sum_{V \in \mathcal V}\sum_{U \in \mathcal U}\sum_{j \in \beta_{U, V}^{(i)}} (h_V) e_j) (\sum_{j} (f \circ \lambda_j) e_j) \\
%&= & p_i (\sum_{j=1}^m (f \circ \lambda_i) e_j).
%\end{eqnarray*}
%This verifies  \eqref{AH-SUPG-cond-2}.
%
%Let us verify \eqref{AH-SUPG-cond-3}. 
%For each $f \in \mathcal F$, its image in $C_2$ is
%$$\sum_{j = 1}^m (f \circ \lambda_j) e_j $$
%and then, for each $i=1, ..., n$ and $\tau \in \tr(A)$, by \eqref{small-image-V} and \eqref{dense-in-X1}, 
%\begin{eqnarray*}
%\tau(p_i(\sum_{j = 1}^m (f \circ \lambda_j) e_j)) & = & \tau((\sum_{V \in \mathcal V}\sum_{U \in \mathcal U}\sum_{j \in \beta_{U, V}^{(i)}} (h_V) e_j) (\sum_{j = 1}^m (f \circ \lambda_j) e_j)) \\
%& = & \tau(\sum_{V \in \mathcal V}\sum_{U \in \mathcal U} \sum_{j \in \beta_{U, V}^{(i)}} (h_V) (f \circ \lambda_j) e_j ) \\
%& \approx_\eps & \tau(\sum_{V \in \mathcal V}\sum_{U \in \mathcal U} \sum_{j \in \beta_{U, V}^{(i)}} (h_V) (f(x_U)) e_j ). 
%\end{eqnarray*}
%By \eqref{equal-beta} and \eqref{large-beta}, it follows that
%$$ \sum_{V \in \mathcal V}\sum_{U \in \mathcal U} \tau (f(x_U)  \sum_{j \in \beta_{U, V}^{(i)}} h_V e_j ) \approx_\eps  \sum_{V \in \mathcal V}\sum_{U \in \mathcal U}  \frac{1}{n}\tau(f(x_U) \sum_{j \in \alpha_{U, V}} h_V e_j ),\quad U\in \mathcal U,\ V \in \mathcal V,\ \tau \in\tr(A),$$
%and therefore, by \eqref{small-image-V}, \eqref{dense-in-X1}, \eqref{SBP-cover-1}, and \eqref{dense-h},
%\begin{eqnarray*}
%\tau( \sum_{V \in \mathcal V}\sum_{U \in \mathcal U} \phi(x_U)  \sum_{j \in \beta_{U, V}^{(i)}} h_V e_j ) 
%& \approx_\eps & \frac{1}{n} \tau(\sum_{V \in \mathcal V}\sum_{U \in \mathcal U} f(x_U) \sum_{j=1}^m h_V e_j ) \\
%& = & \frac{1}{n}\tau( \sum_{V \in \mathcal V}\sum_{U \in \mathcal U}  \sum_{j=1}^m (h_V)(f(x_U)) e_j ) \\
%&\approx_\eps & \frac{1}{n}\tau( \sum_{V \in \mathcal V}\sum_{U \in \mathcal U}  \sum_{j=1}^m (h_V)(\phi\circ \lambda_j) e_j ) \\
%& = & \frac{1}{n}\tau( (\sum_{V \in \mathcal V}\sum_{U \in \mathcal U} h_V e_j)  (\sum_{j=1}^m (f\circ\lambda_j) e_j ) ) \\
%& \approx_{2\eps} & \frac{1}{n} \tau( \sum_{j=1}^m (f\circ\lambda_j) e_j ).
%\end{eqnarray*}
%Thus,
%$$\tau(p_ifp_i) \approx_{2\eps} \tau(p_i f) \approx_{5\eps} \frac{1}{n} \tau(f),\quad i=1, ..., n,\ \tau \in \tr(A),$$
%as desired.
%\end{proof}




\end{document}
